\documentclass[11pt,letterpaper]{article}
\usepackage[T1]{fontenc}
\usepackage[utf8]{inputenc}
\usepackage{lmodern}
\usepackage[margin=1in,headheight=15pt]{geometry}
\usepackage{amsmath,amssymb,amsthm,mathtools,mathrsfs}
\usepackage{microtype,booktabs,longtable,array,enumitem}
\usepackage[dvipsnames]{xcolor}
\usepackage{fancyhdr,listings,xurl,needspace}
\usepackage[colorlinks=true,linkcolor=MidnightBlue,citecolor=MidnightBlue,urlcolor=MidnightBlue,breaklinks=true]{hyperref}
\usepackage[nameinlink,noabbrev]{cleveref}
\hypersetup{%
  pdftitle={Foundations for Surreal and Surcomplex Mathematics},%
  pdfauthor={Merged research exposition},%
  pdfsubject={Sets, classes, universes, type theories, and a pinned Lean audit}}
\pagestyle{fancy}\fancyhf{}
\fancyhead[L]{\small Foundations for surreal and surcomplex mathematics}
\fancyhead[R]{\small\thepage}
\renewcommand{\headrulewidth}{0.3pt}
\setlength{\parskip}{0.3em}
\setlength{\parindent}{1.25em}
\setlength{\emergencystretch}{3em}
\setlist{leftmargin=2em,itemsep=0.2em,topsep=0.35em}
\clubpenalty=10000
\widowpenalty=10000
\displaywidowpenalty=10000
\setcounter{tocdepth}{2}
\numberwithin{equation}{section}
\allowdisplaybreaks[2]
\newtheorem{theorem}{Theorem}[section]
\newtheorem{proposition}[theorem]{Proposition}
\newtheorem{lemma}[theorem]{Lemma}
\newtheorem{corollary}[theorem]{Corollary}
\theoremstyle{definition}
\newtheorem{definition}[theorem]{Definition}
\newtheorem{example}[theorem]{Example}
\newtheorem{principle}[theorem]{Design principle}
\theoremstyle{remark}
\newtheorem{remark}[theorem]{Remark}
\newcommand{\No}{\mathbf{No}}
\newcommand{\SC}{\mathbf{SC}}
\newcommand{\Ord}{\mathbf{Ord}}
\newcommand{\Oz}{\mathbf{Oz}}
\newcommand{\R}{\mathbb R}
\newcommand{\C}{\mathbb C}
\newcommand{\N}{\mathbb N}
\newcommand{\Q}{\mathbb Q}
\newcommand{\Z}{\mathbb Z}
\newcommand{\U}{\mathcal U}
\newcommand{\OO}{\mathcal O}
\newcommand{\HH}{\mathcal H}
\newcommand{\Ag}{\mathscr A}
\newcommand{\Fg}{\mathscr F}
\newcommand{\mm}{\mathfrak m}
\newcommand{\supp}{\operatorname{supp}}
\newcommand{\st}{\operatorname{st}}
\newcommand{\rk}{\operatorname{rank}}
\newcommand{\Spec}{\operatorname{Spec}}
\newcommand{\cis}{\operatorname{cis}}
\newcommand{\bd}{\operatorname{b}}
\newcommand{\Con}{\operatorname{Con}}
\newcommand{\Type}{\mathsf{Type}}
\newcommand{\Prop}{\mathsf{Prop}}
\newcommand{\ETR}{\mathsf{ETR}}
\newcommand{\abs}[1]{\lvert#1\rvert}
\newcolumntype{P}[1]{>{\raggedright\arraybackslash}p{#1}}
\lstdefinestyle{leanstyle}{basicstyle=\small\ttfamily,columns=fullflexible,%
  keepspaces=true,breaklines=true,showstringspaces=false,frame=single,%
  framerule=0.3pt,rulecolor=\color{black!25},xleftmargin=0.5em,xrightmargin=0.5em,%
  aboveskip=0.7em,belowskip=0.7em,commentstyle=\color{black!65},%
  morecomment=[l]{--},morecomment=[s]{/-}{-/},keywordstyle=\bfseries,%
  morekeywords={universe,inductive,structure,where,def,theorem,variable,namespace,%
  end,import,noncomputable,section,by,exact,intro,have,let,fun,Type,Prop,forall,axiom,example}}
\lstset{style=leanstyle}

\begin{document}
\hypersetup{pageanchor=false}
\begin{titlepage}
\centering
\vspace*{0.2cm}
{\large\scshape A merged foundational and formalization audit\par}
\vspace{0.55cm}
{\Huge\bfseries Foundations for Surreal\par and Surcomplex Mathematics\par}
\vspace{0.45cm}
{\LARGE Sets, Classes, Universes, Type Theories,\par and a Pinned Lean Audit\par}
\vspace{0.5cm}
{\large September 21, 2026\par}
\vspace{0.5cm}
\begin{minipage}{0.94\textwidth}\small
\begin{abstract}
\noindent
The phrase ``formalize the surreal numbers'' hides at least four distinct
projects: construct canonical numbers and prove the ordered-field laws;
characterize the resulting simplicity hierarchy; develop algebra over
sufficiently large set-sized surreal subfields; and formalize functions,
supports and analysis over the full class. They share mathematics but not
size requirements and not proof-assistant dependencies.

Surreal numbers are individually set-sized and collectively a proper class.
Their complexification $\SC=\No[i]$ inherits that size without introducing a
new algebraic paradox. The genuine foundational issues are the permitted sizes
of cuts, the coding of quotients, the strength of recursion, the scope of
support and function quantifiers, and the interpretation of topology. This
report compares definable classes in ZF and ZFC, explicit class theories
(GB, GBC, specified NBG), Kelley--Morse, Grothendieck universes and
Tarski--Grothendieck, Feferman reflection, birthday fragments, set-sized Hahn
workspaces, constructive set theory, classical dependent type theory, higher
inductive-inductive and univalent approaches, and an internal set-theoretic
model inside Lean. It restates, size-correctly and with proofs, the workspace,
summability, topology and phase principles of five supplied surcomplex
research manuscripts; it gives an elementary finite-algebra descent theorem, a
support-local nonlinear recursion proposition, two near-mirror examples over
$\C((t^\Q))$ with opposite exponent signs and opposite conclusions, and a
family of topological counterexamples that pin down exactly where the
degeneracy theorems live. It then reports one dated, pinned, read-only
inspection of a contemporary Lean surreal repository, separating what the
inspected source establishes from what remains an unverified bridge, and
proposes a layered implementation architecture.

The central lesson is the discipline of three meanings of ``all'': all members
of a specified set; all values at a specified universe level subject to
smaller-index rules; and the entire definable class. The unifying rule behind
that discipline is that \emph{an individual object can have small data without
the totality of all such objects being small}. Smallness, algebraic closure,
summability and ambient topology must be separate stated hypotheses rather than
connotations of the word ``surreal.''
\end{abstract}
\end{minipage}
\vfill
\end{titlepage}
\hypersetup{pageanchor=true}
\pagenumbering{roman}
\thispagestyle{empty}
\noindent
\begin{minipage}{\textwidth}\small
\textbf{Status.} This report merges three independently written audits. Their
three status declarations are reproduced here as a union rather than rewritten
into one.

\smallskip
\emph{From the first audit.} ``This is a researched mathematical exposition and
implementation blueprint. The supplied manuscripts are treated as motivating
research texts, not as machine-certified libraries. The Lean examples
accompanying this article have not been compiled in this environment. No
unconditional consistency proof, completed formalization of the manuscripts, or
exhaustive survey of every repository is claimed.''

\smallskip
\emph{From the second audit.} ``All consistency assurances are relative to the
chosen foundational theory and its metatheory. This article supplies
constructions, explicit size restrictions, and proof obligations; it does not
claim an absolute proof of the consistency of set theory or Lean.'' Its
illustrative Lean files and source-audit record ``are not represented as an
executed formal verification.''

\smallskip
\emph{From the third audit.} ``This is a researched mathematical exposition and
implementation design, not a machine-checked formalization of the attached
research manuscripts. `Paradox-free' means avoiding the identified invalid
constructions within explicitly stated foundations; no absolute consistency
proof of those foundations is claimed. Lean source observations are pinned
where possible and dated. Illustrative Lean snippets have not been compiled in
this environment.''

\smallskip
\emph{For the merged document.} Three of the five shipped Lean files were
compiled and axiom-audited for this report, under a toolchain that is not the
one the sources pin; nothing else here was machine-checked. See
\cref{found:sec:boundary}, which states the verification boundary in one place
and in plain language.
\end{minipage}
\clearpage
\begingroup\small\setlength{\parskip}{0pt}\tableofcontents\endgroup
\clearpage
\pagenumbering{arabic}

\section{Scope, the five source manuscripts, and the recommended separation}
\label{found:sec:scope}

\subsection{Four distinct projects, three layers, four meanings of ``rigorous''}
\label{found:sub:fourprojects}
There are at least four different projects hidden in the phrase ``formalize the
surreal numbers.'' One can construct canonical numbers and prove their
ordered-field laws; characterize the resulting simplicity hierarchy; develop
algebra over sufficiently large set-sized surreal subfields; or formalize
functions, supports and analysis over the full class. These projects share
mathematics, but they do not have the same size requirements or the same
proof-assistant dependencies.

Cutting the same material a different way, the problem has three layers that
should not be merged. The first is a \emph{foundational construction}: what
objects are surreal numbers, what does it mean to quantify over them, and why
are recursive definitions legitimate? The second is a \emph{mathematical
interface}: which properties of ordered fields, cuts, simplicity, normal forms
and exponentiation are available? The third is \emph{formal verification}: how
can a proof assistant represent these objects and check arguments without
silently changing their size or their topology?

Correspondingly, ``rigorous'' names four different achievements. A precise
mathematical construction specifies its objects and permissible operations. An
axiomatic construction additionally names the assumptions from which existence
and theorems follow. A proof-assistant development supplies proof terms
accepted by a specified kernel with an audited axiom set. A relative-consistency
analysis compares the assumptions with another theory. A Lean declaration
asserting an ordered field together with arbitrary-cut filling can be
syntactically accepted even when the asserted axioms imply a contradiction;
conversely a class construction in ZFC can be rigorous with no proof-assistant
implementation at all.

The surcomplex numbers of the supplied work are
\[
  \SC=\No[i]=\{x+iy:x,y\in\No,\ i^2=-1\},
\]
where the braces denote a \emph{class}, not a set asserted to contain all such
pairs. The algebraic input is that $\No$ is real closed, and consequently
$\SC$ is algebraically closed in the finite-polynomial sense. These are
established results of surreal theory, not consequences of a notational
definition \cite{found:Conway,found:Gonshor,found:Alling}. Neither statement
supplies a complex-analytic topology, a general summation operation, or a set
of all class functions.

\subsection{What ``paradox-free'' is allowed to mean here}
\label{found:sub:paradoxfree}
A rigorous treatment answers four separate questions. \emph{Formation}: are the
purported objects legal sets, classes or types in the chosen foundation?
\emph{Construction}: do the definitions actually produce the stated ordered
field and its complexification? \emph{Metatheory}: what assumptions justify
confidence in the underlying formal system? \emph{Verification}: has a proof
assistant checked the particular definitions and proofs, with a recorded axiom
dependency? A positive answer to the first does not imply a positive answer to
the second. An axiomatic structure with impossible fields can be syntactically
well formed and uninhabited. Conversely, an inconvenient universe error is
usually a size mismatch, not evidence that surreal numbers are inconsistent.

In this report, \emph{paradox-free formalization} means a construction with
explicit formation rules, legitimate recursion and comprehension, and no hidden
assumptions that contradict its own order or universe laws. Consistency remains
relative to the foundation used. More sharply still, ``paradox-free'' denotes
only the explicit exclusion of the specific invalid constructions analysed
below --- the self-cut, the proper-class-member quotient, the unrestricted
support sum, and incompatible equality principles --- and nothing more. For a
consistent, effectively axiomatized theory strong enough for arithmetic, the
second incompleteness theorem prevents the theory from proving its own
standard arithmetized consistency statement under that theorem's hypotheses.
This obstructs neither the construction of surreal numbers nor any ordinary
mathematical existence theorem. It tells us not to advertise a construction in
ZFC or Lean as an absolute proof that its foundation cannot derive a
contradiction \cite[\S5]{found:ShulmanSize}.

The preferred architecture for the written mathematics is
\begin{equation}\label{found:eq:architecture}
 \begin{gathered}
 \boxed{\text{set-sized algebra and analysis}}\\[0.35em]
 +\ \boxed{\text{explicit size-aware comparison maps}}\\[0.35em]
 +\ \boxed{\text{a class or universe interpretation}}.
 \end{gathered}
\end{equation}
The first component carries most proofs. The second prevents accidental changes
of meaning. The third supplies the genuinely unbounded interpretation of $\No$
and $\SC$.

\subsection{The five supplied manuscripts, and who read which}
\label{found:sub:ledger}
This report merges three independently prepared audits, referred to throughout
as \textbf{M1}, \textbf{M2} and \textbf{M3}. They were given \emph{different}
collections of source manuscripts. Five distinct surcomplex research
manuscripts are involved, and \emph{only one of them was read by all three
audits}. The full ledger --- titles, filenames, recorded hashes, and which
audit read which text --- is \cref{found:app:ledger}. The essential facts are
these.

\begin{description}
\item[A, the Analysis manuscript.] \emph{Surcomplex Analysis: Infinitesimal
Calculus, Hahn-Coherent Holomorphy, and Contour Theory}
\cite{found:SrcA}. Read by all three audits. M1 and M3 cite it under the
filename \texttt{surcomplex\_analysis(3).tex}; M2 cites it under the filename
\texttt{surcomplex\_analysis(2).tex}. The SHA-256 values recorded independently
by M2 and by M3 are identical
(\texttt{e6e2ef4a\ldots17748f4}), so this is one file carrying two labels. We
do not silently pick a filename and we assert nothing about which label is
correct.
\item[AG, the Analytic Geometry manuscript.] \emph{Infinitesimal Analytic
Geometry over the Surcomplex Numbers} \cite{found:SrcAG}. Read by \textbf{M1
only}. No hash was recorded for it by any audit.
\item[T, the Trigonometry manuscript.] \emph{Trigonometry on the Surcomplex
Plane} \cite{found:SrcTrig}. Read by \textbf{M2 only}; hash recorded by M2.
\item[P, the Polynomial Algebra manuscript.] \emph{Surcomplex Polynomial
Algebra: Order Geometry, Newton Profiles, and Scale-Resolved Stability}
\cite{found:SrcP}. Read by \textbf{M3 only}; hash recorded by M3.
\item[R, the Hartogs manuscript.] \emph{Hartogs Coherence, Finite Maps, and
Residue Duality over the Surcomplex Numbers} \cite{found:SrcR}. Read by
\textbf{M3 only}; hash recorded by M3.
\end{description}

\noindent\textbf{Tagging convention.} Because the audits read different texts,
every claim in this report that is derived from AG, T, P or R carries a visible
tag naming the single audit that read that source. No sentence in this report
says ``the supplied manuscripts'' when it means a text that only one member
saw. Claims derived from A may be attributed to the three audits jointly, and
are so marked. This convention is not decorative: it is the only thing
preventing a merged document from reading as an audit of one coherent corpus
when it is in fact an audit of five texts with a single common element.

Two further non-claims about the sources travel with this ledger. First, P
refers to a separate \emph{global-divisors} manuscript that was never supplied
for any of the three tasks; no assessment of that fourth text is offered here
(M3). Second, other archives named in the bibliographies of the supplied
manuscripts --- in particular in the trigonometry bibliography --- were not
examined (M2). All five manuscripts are user-supplied research expositions, not
independently certified libraries or refereed publications, and their original
files are not duplicated in any of the three archives.

\subsection{What each manuscript assumes}
\label{found:sub:assumptions}
\emph{A} (all three audits) explicitly chooses NBG with global choice and
stipulates that individual supports, index families, coefficient sequences and
the data of its constructions are sets. It distinguishes fine-local,
Hahn-coherent and all-scale theories. Its statements about arbitrary small
rescaling and about all-scale rigidity use the full surreal class, not merely
one fixed field. Its contour functional is coefficientwise; it is not
integration along arbitrary fine-continuous real-parameter paths. Its fine
topology is described by balls with arbitrary positive surreal radii, not by a
purported set of all open subclasses.

\emph{AG} (M1 only) works with a nonzero divisible set-sized ordered subgroup
$\Gamma\subset\No$ and writes
\begin{equation}\label{found:eq:workspaces}
 t^\gamma=\omega^{-\gamma},\qquad
 K_\Gamma=\C((t^\Gamma)),\qquad
 R_n=\C\{z_1,\ldots,z_n\},\qquad
 \Ag_{n,\Gamma}=R_n((t^\Gamma)).
\end{equation}
Its coefficients are convergent complex germs \emph{without} a shared ordinary
radius. It claims several-variable division, a Nullstellensatz, finite
complete-intersection deformation, residue duality and finite-zero workspace
invariance, with explicit mathematical proofs and without machine verification.

\emph{T} (M2 only) retains the same support convention. Its circular functions
are initially defined on finite surreal arguments, its geometric lengths range
over all surreals, and its extensions to infinite circular arguments require
extra phase data. These are substantive domain restrictions, not presentational
conveniences.

\emph{P} (M3 only) confines polynomials, supports, matrices and the relevant
recursions to sets; its degrees are ordinary natural numbers; it introduces
divisible set-sized exponent groups $\Gamma$ and algebraically closed fields
$K_\Gamma$.

\emph{R} (M3 only) uses the same workspaces but also considers an arbitrary
class function on a full infinitesimal thickening, and invokes class
replacement after restricting that function to an ordinary set of complex
points.

The three manuscripts M3 read state that they have not been independently
proof-assistant verified; the same holds of the two M1 read and the two M2
read. Their theorems are therefore \emph{targets and supplied arguments}, not
imported formal certificates. Size-correctness of $\Ag_{n,\Gamma}$ does not
certify its Noetherianity, and formalization of an abstract field does not
certify a proposed analytic residue formula (M1).

\subsection{Three levels of assertion, maintained throughout}
\label{found:sub:threelevels}
The following discipline, taken from M1, is used in every section.
\begin{enumerate}
\item Classical construction facts are attributed to established literature.
\item Statements about A, AG, T, P and R are explicitly marked source-derived,
with the tag naming the audit that read the source.
\item Propositions proved in this report are elementary foundational or
algebraic deductions, included to make the proposed architecture precise; no
novelty claim, no priority claim and no resolution of a published open
conjecture is attached to any of them.
\item Statements about software describe \emph{inspected declarations or
documentation}, not a successful independent rebuild. Proposed module names and
future APIs are design suggestions.
\end{enumerate}

\subsection{A practical answer before the details}
\label{found:sub:practical}
For a conventional mathematical exposition, a carefully specified version of
GBC is a convenient language: individual numbers and supports are sets; $\No$,
its operations and particular global analytic functions are classes. The basic
theory can instead be expressed in ZFC through definable classes, with no new
large-cardinal axiom merely for surreal arithmetic.

For Lean, the most effective architecture is usually different. Work with a
universe-indexed surreal type, prove theorems about \emph{small} families at
that level, and implement the analytic core over ordinary set-sized Hahn
fields. Reuse the existing game-based surreal field rather than starting its
algebra from zero. Prove explicit comparison theorems between the game, sign
and Hahn presentations. The full normal-form theorem and the analytic
enrichments should be tracked as mathematical dependencies, not inferred from
the existence of suggestively named modules.

This \emph{split} --- one recommendation for the mathematics, a different and
compatible one for the implementation --- is the spine of the recommendation in
\cref{found:sec:comparison}. It is the only framing among the three audits that
names the two questions separately, and we adopt it for that reason.

\section{The verification boundary of this report}
\label{found:sec:boundary}

\textbf{None of the three source audits compiled any Lean file, rebuilt any
repository, or ran \texttt{\#print axioms}. For the merged report that changed
in exactly one respect, recorded in \cref{found:sub:compiled}: three of the five
shipped Lean files were compiled, and \texttt{\#print axioms} was run on every
declaration in them. Nothing else here was machine-checked. No repository was
rebuilt, and no supplied manuscript was verified.}

Those sentences are the whole of the verification status of the present
document, and they are stated here, early, in their own section, rather than
distributed across footnotes where a skimming reader would miss them. The
remainder of this section records exactly what \emph{was} done, by whom, and
what follows from it.

\subsection{One inspection, three readers}
\label{found:sub:onethree}
All three audits inspected the \emph{same} public Lean repository at the
\emph{same} immutable revision, on the same day, by reading source through the
public web interface. All three record the same commit hash, the same
toolchain, the same field instance, the same small-support Hahn type, and the
same absent bridges. Printed three times, or printed once in confident merged
prose with three citations, that agreement would be indistinguishable from
independent confirmation. It is not.

Three agreeing readings of one source corroborate each other only about what
that source says. They do not corroborate that the source compiles, that its
theorems are what their names suggest, or that its axiom closure is clean.
Accordingly the findings of that inspection are printed \emph{once}, in
\cref{found:sec:ecosystem}, framed as a single act of source reading with three
concurring readers, and nowhere in this report are they presented as three
pieces of evidence.

\subsection{What was executed, and by whom}
\label{found:sub:executed}
The only step any of the three preparation environments actually executed was a
\LaTeX{} build of its own article.

\begin{longtable}{@{}P{0.13\textwidth}P{0.40\textwidth}P{0.40\textwidth}@{}}
\toprule
Audit & Executed & Not executed\\
\midrule\endhead
M1 & pdf\LaTeX{} build; PDF compiled and rendered for layout inspection.
   & Every Lean file. No Lean or Lake executable was available. No build log,
     no exit status, no \texttt{\#print axioms} output anywhere in the archive.\\[0.5em]
M2 & pdf\LaTeX{} 3.141592653-2.6-1.40.26 (TeX Live 2025/dev/Debian); 32 pages,
     17 main sections, 2 appendices, 28 references; zero unresolved reference or
     citation warnings; zero overfull or underfull box warnings; all 32 pages
     rendered and inspected.
   & Lean compilation recorded in machine-readable form as
     ``NOT EXECUTED; Lean/Lake were not installed in the preparation
     environment.''\\[0.5em]
M3 & \texttt{latexmk}/pdf\LaTeX{} completed successfully; 40 pages, 94 PDF
     bookmarks, 36 references; cross-references stabilized with no unresolved
     references; 0 overfull boxes; 0 missing glyphs; all pages rendered and
     visually inspected. SHA-256 recorded for both its own \texttt{.tex} and its
     own \texttt{.pdf}.
   & Lean compilation not run; external Lean project rebuild not run; transitive
     \texttt{\#print axioms} audit not run; Mizar execution not run; supplied-manuscript
     formal verification not performed.\\
\bottomrule
\end{longtable}

\noindent In M3's own words, its build report ``certifies document-production
checks only.'' M2's build report states its scope as ``PDF production checks
and source inspection, not formal mathematical verification.''

\subsection{Three archives, three different code inventories, none compiled}
\label{found:sub:archives}
The three archives do not ship the same things, and the difference matters for
what a future reader could actually check. The ``State'' column below records
each file \emph{as its archive delivered it}; ``Not compiled'' is a statement
about that archive, not about the present report, which compiled three of the
five (\cref{found:sub:compiled}).

{\small
\begin{longtable}{@{}P{0.08\textwidth}P{0.25\textwidth}P{0.20\textwidth}P{0.34\textwidth}@{}}
\toprule
Audit & Lean files shipped & Import footprint & State as delivered\\
\midrule\endhead
M1 & \path{FoundationsSketch.lean} (38 lines)
   & \texttt{Init} \emph{only}
   & Not compiled. Contains \texttt{noUniversalStrictBound}, \texttt{RawGame},
     \texttt{SmallCutData}, and one \texttt{\#print axioms} command whose output
     is recorded nowhere.\\[0.4em]
M1 & \path{ComplexifySketch.lean} (35 lines)
   & \path{Mathlib.Algebra.Ring.Basic}
   & Not compiled. Proves exactly one theorem, $i^2=-1$. No zero, no addition,
     no subtraction, no inverse, no \texttt{Ring} and no \texttt{Field} instance.\\[0.4em]
M1 & \path{InspectCurrentLibrary.lean} (24 lines)
   & Pinned repository plus Mathlib
   & Not run. Seven \texttt{\#check}/\texttt{\#synth} commands whose output is
     recorded nowhere.\\[0.4em]
M2 & \path{LogicalGuards.lean}
   & \texttt{import Mathlib} (all of it)
   & Not compiled. Contains \texttt{noUnrestrictedUpperBounds}, \texttt{RawGame},
     \texttt{QuadraticPair} with \texttt{normSq\_mul} and \texttt{gramIdentity}
     proved \texttt{by ring}, plus three unexecuted \texttt{\#print axioms} lines.\\[0.4em]
M2 & \path{SourceAudit.lean}
   & Pinned repository
   & Not run.\\[0.4em]
M3 & \emph{none}
   & ---
   & M3 ships no Lean files at all. Its five code blocks exist only as
     \texttt{lstlisting} snippets typeset inside its PDF, each marked in the
     document itself as an illustrative design.\\
\bottomrule
\end{longtable}
}

Note the asymmetry recorded in the third column. M1's \texttt{FoundationsSketch}
imports \texttt{Init} only and is therefore checkable without any Mathlib build;
M2's \texttt{LogicalGuards} imports the whole of Mathlib. So even the one
theorem both archives contain --- the machine-level form of the unrestricted-cut
obstruction --- has two entirely different dependency footprints. Neither was
run.

M2's README lists a \texttt{SHA256SUMS} file among the package contents, and
the delivered archive did contain one. It is not present in this repository:
the collection carries no checksum manifests, so the file was removed when the
package was unpacked. Nothing was dropped in transit, and the README did not
overstate its archive's contents. M2's README also gives clone/checkout/build/run commands for
both POSIX \texttt{sh} and PowerShell and explicitly labels them ``a
reproduction recipe, NOT commands executed while preparing this package,'' with
the instruction: build the upstream project with its original pinned dependency
manifest, and \emph{do not run \texttt{lake update} and then describe that as a
check of this exact revision}. That instruction is reproduced here because it
is the single most likely way for a future reader to produce a result that
looks like a check of the pinned revision and is not one.

\subsection{What the merged document is, and what its length does not buy}
\label{found:sub:notacertificate}
This report is longer than any of its three sources. It cites five source
manuscripts, proves more than thirty propositions, and carries more than fifty
references. None of that is execution. The union of three read-only audits is a
larger blueprint, not a certificate. Specifically:

\begin{itemize}
\item zero of the three audits compiled any Lean file;
\item zero rebuilt the inspected repository;
\item zero recorded any \texttt{\#print axioms} output;
\item zero performed a transitive axiom audit;
\item zero ran Mizar;
\item zero machine-verified any of the five supplied manuscripts.
\end{itemize}

\noindent Every one of those six remains true of the three source audits. The
merged report changes the first and the third, and only for the three files
named in \cref{found:sub:compiled}. The other four are unchanged: no
repository was rebuilt, no transitive axiom audit was performed, no Mizar was
run, and no supplied manuscript was machine-verified.

\subsection{What was compiled for this report}
\label{found:sub:compiled}

The three audits each stated that no Lean or Lake executable was available where
they were prepared. Both were available where this merge was prepared, so the
shipped files were put through them. This subsection is the record.

\emph{The toolchain was not the pinned one.} The three audits pin
\texttt{leanprover/lean4:v4.35.0-rc2} with Mathlib
\texttt{dec5b2b780537b6eaf7f5e5f000c12f7387fb24d}. That toolchain is not
installed on the machine used here. The compilations below ran under Lean
\texttt{4.32.0} against a built Mathlib at the matching release. This is a
\emph{nearby but different} toolchain, and nothing below licenses any claim
about the pinned revision or about code that depends on it.

\begin{center}\small
\begin{tabular}{@{}p{0.29\textwidth}p{0.17\textwidth}p{0.20\textwidth}p{0.26\textwidth}@{}}
\toprule
\textbf{File} & \textbf{Imports} & \textbf{As shipped} & \textbf{Axioms reported}\\
\midrule
M1 \texttt{FoundationsSketch} & \texttt{Init} only & compiles, exit 0 &
  \texttt{noUniversalStrictBound}: none at all\\[0.35em]
M1 \texttt{ComplexifySketch} & Mathlib & compiles, exit 0 &
  three of five none; \texttt{neg} and \texttt{imaginaryUnit\_sq}:
  \texttt{propext}\\[0.35em]
M2 \texttt{LogicalGuards} & Mathlib & \textbf{does not compile} &
  after repair: \texttt{noUnrestrictedUpperBounds} none; two others
  \texttt{propext}, \texttt{Quot.sound}\\[0.35em]
M1 \texttt{InspectCurrentLibrary} & \texttt{CombinatorialGames} & not attempted &
  repository absent\\[0.35em]
M2 \texttt{SourceAudit} & \texttt{CombinatorialGames} & not attempted &
  repository absent\\
\bottomrule
\end{tabular}
\end{center}

\textbf{No \texttt{sorryAx} and no \texttt{Classical.choice} appears anywhere.}
Every axiom reported is \texttt{propext} or \texttt{Quot.sound}, both ordinary
kernel axioms of Lean~4 rather than holes. In particular the two size-obstruction
theorems that carry the argument of \cref{found:sec:size} --- M1's
\texttt{noUniversalStrictBound} and M2's \texttt{noUnrestrictedUpperBounds} ---
depend on no axioms whatever, so both are fully constructive.

\emph{M2's \texttt{LogicalGuards.lean} does not compile as shipped.} The kernel
reports \texttt{invalid 'import' command, it must be used in the beginning of the
file}. The file opens with a \texttt{/-!\,\ldots\,-/} module docstring and only
then writes \texttt{import Mathlib}. A module docstring is a command, and no
command may precede an import. M1's files open with a plain
\texttt{/-\,\ldots\,-/} comment, which is legal because a comment is not a
command; the entire difference is the exclamation mark. Moving the import above
the docstring makes the file compile, and the three axiom probes it already
contains then produce the results tabulated above. The file ships here exactly as
delivered, defect included.

\emph{What this does and does not establish.} It establishes that three small
illustrative files are syntactically and logically well formed under a current
toolchain, and that the two impossibility theorems are constructive. It
establishes nothing about the pinned repository, about the surreal field
instance, about the absent normal-form bridge or real-closedness instance, or
about any theorem in the five supplied manuscripts. Compiling a
thirty-eight-line elementary impossibility theorem verifies that theorem and
nothing beyond it.

Source inspection and compilation are different verification levels. A
successful execution of the small illustrative examples in the M1 and M2
archives --- which is now on record for three of them --- does not constitute
verification of the full surreal theory or of any manuscript. The deliverable here is a mathematical exposition, a foundational
audit and an implementation blueprint. \Cref{found:app:limitations} collects the
complete union of the three limitation lists.
\section{The size obstructions every foundation must respect}
\label{found:sec:size}

\subsection{Set-sized numbers, a proper-class totality}
\label{found:sub:properness}
One canonical presentation identifies a surreal with a sign sequence
\[
 s:\alpha\longrightarrow\{-,+\},\qquad \alpha\in\Ord,
\]
with birthday $\bd(s)=\alpha$. Each such function is a set. At every fixed
ordinal bound $\beta$, the sequences of length less than $\beta$ form a set:
the relevant function sets exist, and their union is indexed by the set
$\beta$. The total collection over all ordinals is a definable class. The empty
sequence is zero; the constantly positive sequence of length $\alpha$
represents the ordinal $\alpha$ \cite{found:Gonshor,found:Conway}.

All three audits prove the following by the same Replacement argument on the
birthday map. We print it once, in M3's slightly stronger form, which carries
the bounding clause as part of the statement.

\begin{proposition}[There is no set of all surreal numbers]\label{found:prop:proper}
In the sign-sequence construction the class $\No$ is not a set. More strongly,
every \emph{set} of surreal numbers has its birthdays bounded by an ordinal.
\end{proposition}
\begin{proof}
For a set $B$ of sign sequences, Replacement gives the set
$\{\operatorname{dom}(s):s\in B\}$ of their domains. The ordinal
\[
 \beta=\sup\{\bd(x)+1:x\in B\}
\]
bounds their birthdays. If $B$ contained every surreal it would contain the
constantly positive sequence of length $\beta$, whose birthday is $\beta$ and
so is not less than $\beta$: a contradiction. The empty case is harmless.
\end{proof}

The same argument is the size mechanism behind the Burali-Forti obstruction: a
collection containing a distinct canonical object at every ordinal stage cannot
itself be one of the bounded sets in that construction. It is not a defect in
surreal arithmetic. The ordered-pair construction $\SC=\No\times\No$ with the
intended complex arithmetic is likewise a class: an individual pair is a set,
but the embedding $x\mapsto(x,0)$ precludes a set containing all pairs.

\subsection{A size ledger}
\label{found:sub:sizeledger}
The following seven-row ledger is M1's, and is a pedagogic artifact rather than
a theorem. It pairs each notation with the reading that is safe and the reading
that is not.

\Needspace{18\baselineskip}
\begin{longtable}{@{}P{0.25\textwidth}P{0.31\textwidth}P{0.36\textwidth}@{}}
\toprule
Expression & Safe reading & Dangerous reading\\
\midrule\endhead
$x\in\No$ & A set code satisfying a definable predicate & An element of a universal set containing all surreal codes\\[0.4em]
$\{L\mid R\}$ & $L,R$ are sets, with every left option below every right option & A cut of two arbitrary proper classes\\[0.4em]
$\sum r_\alpha\omega^{a_\alpha}$ & A set-length normal form with suitable support & An unspecified sum over all ordinals\\[0.4em]
$\SC[[X]]$ & A class of set-coded countable coefficient sequences & A set-sized ring without a size restriction\\[0.4em]
$\Ag_{n,\Gamma}$ & A set ring for a fixed set $\Gamma$ & All possible exponent groups silently combined into one set\\[0.4em]
A class function & A class of ordered pairs of sets, satisfying functionality & A set-valued element of an unrestricted space of all class functions\\[0.4em]
A class-valued sheaf & Coded section and restriction relations, or a larger-universe sheaf & A small-category-valued functor with unacknowledged proper-class values\\
\bottomrule
\end{longtable}

\subsection{The unrestricted-cut obstruction: the interface regression test}
\label{found:sub:cutobstruction}
A Conway presentation $\{L\mid R\}$ uses \emph{sets} of previously available
numbers with $l<r$ for every $l\in L$ and $r\in R$, and denotes the simplest
number strictly between the two sides. It does not assert that every pair of
subclasses with that separation property admits a separator. All three audits
prove the following, by substituting the whole carrier on the left against the
empty right side, and all three identify it as the single most useful
regression test in the subject.

\begin{proposition}[The unrestricted cut axiom is inconsistent]\label{found:prop:allcuts}
Let $X$ carry an irreflexive relation $<$. There is no operation assigning to
every subcollection $A\subseteq X$ an element $b(A)\in X$ with $a<b(A)$ for
every $a\in A$. Equivalently: for a nonempty strict order $S$, the assertion
that every pair of subclasses $L,R\subseteq S$ with $L<R$ has an element
strictly between them is false once the permitted pairs include
$(S,\varnothing)$.
\end{proposition}
\begin{proof}
Take $A=X$ and then $a=b(X)$. The required inequality is $b(X)<b(X)$,
contradicting irreflexivity. In the two-sided form, separation of $S$ from the
empty right side is vacuous, so a separator $a\in S$ would satisfy $x<a$ for
every $x\in S$; take $x=a$.
\end{proof}

Thus the surreal cut $\{\No\mid\varnothing\}$ is not a new, extremely large
surreal. It is an \emph{illegal input}. The valid theorem bounds every
\emph{set} of surreal numbers by another surreal; in a universe-indexed
implementation it bounds every sufficiently \emph{small} family, not every
subset of the entire larger carrier. Moving from sets to classes does not
validate the unrestricted axiom; it makes the excluded input easier to
describe.

This is stronger guidance than ``beware of Russell's paradox.'' It supplies a
one-line regression test for an axiom interface: \emph{any interface granting
strict upper bounds for all subsets of its own nonempty carrier is already
inconsistent, regardless of its name}. The same reasoning works internally in
type theory whenever a rule accepts the entire carrier as an indexing family
(M2), and it recurs in \cref{found:sec:types} as the reason the universe shift
in the raw-game constructor is the construction rather than an inconvenience.

\begin{lemma}[Bounds for sets of surreals; the small positive lower bound]\label{found:lem:bounds}
Every set $B\subseteq\No$ is bounded above and below in $\No$. For every set
$A\subseteq\No_{>0}$ there is $\delta>0$ with $\delta<a$ for all $a\in A$.
\end{lemma}
\begin{proof}
The legal cut $\{B\mid\varnothing\}$ gives an upper bound and its negation a
lower bound. The legal cut $\{0\mid A\}$ gives the final assertion; for
$A=\varnothing$ take $\delta=1$.
\end{proof}

Here each input is a set, so \cref{found:prop:allcuts} does not apply. There is
no greatest surreal, yet every set of surreals has an upper bound; the
collection of all surreals is precisely not one of those sets. All three audits
give this lemma, in the second form under the names ``small positive
lower-bound principle'' (M2) and ``a lower bound for a set of positive
tolerances'' (M1). Two further remarks travel with it. The hypothesis that $A$
is a set is essential. And $\delta$ need not lie in any previously specified
subfield containing $A$ (M2) --- a point that becomes the whole content of
\cref{found:prop:twotopologies}.

\subsection{Four different assertions involving ``all''}
\label{found:sub:fourall}
The following four statements have different meanings (M2):
\begin{enumerate}
\item Every ordinal of the ambient set universe embeds in the surreal class.
\item Every ordinal small relative to a specified universe embeds in its
associated surreal type.
\item Every set of surreal coefficients fits in a sufficiently large set-sized
workspace.
\item There is one set-sized workspace containing every surreal number.
\end{enumerate}
The first three have rigorous formulations. The fourth contradicts the first. A
theorem producing a workspace \emph{after} a set of inputs has been given does
not produce a universal workspace. This is the same distinction as the three
meanings of ``all'' identified by M1 --- all members of a specified set; all
values at a specified universe level subject to smaller-index rules; the entire
definable class --- which that audit identifies as the actual source of
apparent paradox.

\subsection{Bounds are not suprema}
\label{found:sub:notdedekind}
All three audits prove that the surreal ordered field is not Dedekind complete,
with three different witnesses for one fact: M1 halves an upper bound of $\N$,
while M2 and M3 subtract one.

\begin{proposition}\label{found:prop:incomplete}
The surreal ordered field is not Dedekind complete, even for set-sized bounded
subsets.
\end{proposition}
\begin{proof}
The set $\N$ is bounded above, for instance by $\omega$. If $h$ is an upper
bound then $h\ge 2n+1$ for every $n\in\N$, so $h/2>n$ for every $n$ and $h/2$
is a strictly smaller upper bound. Equivalently, $s\ge n+1$ for all $n$ makes
$s-1$ another upper bound. Either way there is no least upper bound.
\end{proof}

More generally (M3), if a nonempty set $B$ has no greatest member and $b$ is an
upper bound, then every member of $B$ is strictly below $b$, and the legal cut
$\{B\mid\{b\}\}$ produces a smaller upper bound. The cut construction selects
the \emph{simplest} separator; it does not select a least upper bound in the
numerical order. Consequently ``complete'' must always be qualified: real
closedness, saturation for small cuts, valuation completeness and Dedekind
completeness are not interchangeable. A Lean declaration asserting
ordered-completeness for all surreals must therefore be rejected, and the phrase
``complete ordered field'', which occurs in introductory software prose, must
not be read as a formal Dedekind-completeness theorem. All three audits make
this last point; \cref{found:sec:ecosystem} records where the phrase actually
occurs.

\subsection{Russell's paradox, universal function spaces, and universal sets}
\label{found:sub:russell}
In ZF, Separation forms $\{x\in a:\varphi(x)\}$ from a pre-existing set $a$; it
does not form a universal set. In an explicit class theory the Russell class is
allowed as a class but cannot be a set member. Similarly, the graph of a class
function $\No\to\No$ may be a class, but the collection of \emph{all} such
class functions is not automatically another class whose elements are those
graphs. Classes in NBG are not eligible to be members merely because they have
been named (M3). The same warning applies to ``all open subclasses,'' ``the
category of all classes,'' and ``all equivalence classes of a class-sized
relation.'' Each is useful informal language; a formalization needs a level of
coding, a universe bound, a scheme of assertions, or a stronger framework.

Theories with a universal set, such as NFU and other stratified theories, are
not a shortcut (M1, M2). They restrict comprehension, type assignment, or the
class of admissible constructions. A theory with a universal carrier and a
strict order cannot also prove that every subset of that carrier has a strict
upper bound inside it: \cref{found:prop:allcuts} is untouched. One cannot
combine the existence of a universal set from one foundation with unrestricted
ZFC-style subset and cut rules from another. A full surreal development in such
a theory would require its own translation of ordinal, support and smallness
conditions. \emph{No claim of such a completed translation is made here} (M1).

\section{Constructions of the carrier}
\label{found:sec:constructions}

\subsection{Raw games and the numeric quotient}
\label{found:sub:rawgames}
A raw game is a well-founded, set-branching tree with designated left and right
options at each node. Working with raw games first avoids requiring a numeric
comparison predicate in the constructor itself. Define comparison recursively on
pairs of trees, then define which games are numeric. For numeric raw games the
comparison law is
\begin{equation}\label{found:eq:comparison}
x\le y\quad\Longleftrightarrow\quad
\bigl(\forall x^L\in L_x,\ \neg(y\le x^L)\bigr)
\ \land\
\bigl(\forall y^R\in R_y,\ \neg(y^R\le x)\bigr).
\end{equation}
A raw game is numeric when its options are numeric and every left option is
strictly less than every right option. Classical comparison and numeric
induction establish that numerical equivalence
$x\sim y\iff x\le y\land y\le x$ is an equivalence relation and that the
quotient order is linear. Equal values can have different presentations: both
$\{\varnothing\mid\varnothing\}$ and $\{-1\mid1\}$ represent zero. The
distinction between the preorder on presentations and equality of numbers must
remain visible until quotienting \cite{found:Conway}.

The recursive operations begin with
\begin{align}
 -x&=\{-R_x\mid -L_x\},\label{found:eq:negcut}\\
 x+y&=\{L_x+y,\ x+L_y\mid R_x+y,\ x+R_y\},\label{found:eq:addcut}\\
 xy&=\{l_xy+xl_y-l_xl_y,\ r_xy+xr_y-r_xr_y\mid\nonumber\\[-0.3em]
 &\hspace{3.9em}l_xy+xr_y-l_xr_y,\ r_xy+xl_y-r_xl_y\},\label{found:eq:mulcut}
\end{align}
where each displayed family ranges over the appropriate option sets.
Numericity, compatibility with $\sim$, associativity, distributivity and the
inverse construction are separate proof obligations. Merely printing
\eqref{found:eq:mulcut} does not provide a field structure, and inversion is not
justified by displaying the multiplication cut.

A rigorous development records three independent claims at each recursive
definition (M2): the recursive calls are well founded; the option families are
sets; and the constructed left options lie below the right options.
Well-foundedness alone does not imply that a game is a number.

\subsection{Recursion measures: what is and is not a valid substitute}
\label{found:sub:recursionmeasure}
There are two distinct orders. Numerical $<$ is not well founded --- it has the
descending chain $0,-1,-2,\ldots$ --- while simplicity, that is, taking a strict
prefix of a sign sequence or passing to an option, is well founded. Recursive
comparison and arithmetic must decrease the latter kind of complexity.

For a pair $(x,y)$, the Hessenberg natural-sum measure
\begin{equation}\label{found:eq:natsum}
 \bd(x)\mathbin{\#}\bd(y)
\end{equation}
is strictly increasing in each argument, and replacing either argument by an
option strictly decreases it. Two \emph{different} wrong candidates are
excluded by the two audits that considered the question, and both exclusions
survive here because they eliminate different errors:
\begin{itemize}
\item the ordinary \textbf{maximum} of the two birthdays is not a valid general
substitute, since decreasing the smaller argument can leave the maximum
unchanged (M1);
\item ordinary \textbf{ordinal addition} is not a valid substitute, since the
natural sum is strictly increasing in each coordinate while ordinary ordinal
addition is not strictly increasing in its left coordinate (M3).
\end{itemize}
M2 records only the weaker statement that some well-founded relation on pairs of
birthdays or option trees justifies the recursion. Formal termination arguments
must name such a relation; ``one of the inputs became simpler'' is not a checked
termination certificate.

There is also a route that avoids a global pair-recursion theorem entirely
(M1). Fix the two inputs and collect all their subpositions; their descendant
closures are sets. Recursion then takes place on a set-sized well-founded
relation, and uniqueness proves agreement as the input-dependent domains
enlarge. For implementations with indexed option presentations this
dependency-local argument is preferable to assuming that all syntactic forms of
a given numerical birthday form a small collection. The general principle
behind it is \cref{found:prop:recursion}.

\subsection{Why a literal proper-class quotient is unsafe}
\label{found:sub:classquotient}
The equivalence class of a numerical value among all raw presentations can be a
proper class. All three audits use the same motivating failure: the cuts
$\{-\alpha\mid\alpha\}$ for positive ordinals $\alpha$ all have numerical value
zero, while their presentations have unbounded rank. So ``the class of
equivalence classes'' cannot mean a class whose elements are those proper
classes; in GB or NBG, proper classes are not elements.

M1 names three safe alternatives: canonical sign sequences; restricting the
prequotient to a set or a fixed type universe before quotienting; and Scott's
trick. M2 and M3 both promote the third to a numbered proposition with proof.
We print that proof once.

\begin{proposition}[Set codes for a definable class quotient]\label{found:prop:scott}
Let $E$ be a definable equivalence relation on a definable class $D$ of sets.
For $x\in D$ let $\rho_x$ be the least rank of a set $E$-equivalent to $x$, and
put
\[
 C_x=\{y\in D\cap V_{\rho_x+1}:E(x,y)\ \text{and}\ \rk(y)=\rho_x\}.
\]
Then $C_x$ is a nonempty set, and $C_x=C_z$ if and only if $E(x,z)$.
\end{proposition}
\begin{proof}
The ranks under consideration form a nonempty definable class of ordinals
containing $\rk(x)$; a least one is obtained by searching the set $\rk(x)+1$.
All candidates of that rank lie in $V_{\rho_x+1}$, and Separation forms $C_x$.
If $E(x,z)$ then the two classes of equivalent representatives agree, hence so
do their least ranks and their codes. Conversely a member of their common
nonempty code is equivalent to both $x$ and $z$, so transitivity gives
$E(x,z)$.
\end{proof}

The honest caveat, stated by both M2 and M3, is that the resulting collection of
codes can still be a proper class: the repair is that every \emph{code} is a
set, not that the quotient has become small. Scott coding preserves an
extensional quotient with no global choice of representatives, but transporting
recursive functions and simplification through these codes adds overhead. That
is an engineering reason to prefer canonical signs or a native type-theoretic
quotient, not a foundational defect.

\subsection{Sign expansions as the canonical carrier}
\label{found:sub:signs}
Sign sequences avoid the quotient at the level of the carrier. For two distinct
sequences, compare at their first disagreement, treating termination as an
intermediate symbol:
\[
 -\ <\ \text{termination}\ <\ +.
\]
A set of disagreement positions has a least member, so this is a classical
linear order. Each sign sequence has canonical left and right options obtained
by truncating immediately before a plus or a minus respectively; these families
are indexed by subsets of its ordinal length and are therefore sets, and they
reconstruct the sequence as the simplest separator. This is a particularly
clean foundation for birthdays, initial embeddings and no-size-collapse
theorems \cite{found:Gonshor,found:Ehrlich2001}.

\begin{remark}[Simplicity is not birthday precedence]\label{found:rem:simplicity}
Write $x\prec_s y$ for ``$x$ is a proper initial segment of $y$'' and
$x\prec_b y$ for ``$\bd(x)<\bd(y)$.'' The second does not imply the first: the
sequence $+$ is shorter than $--$ but is not its initial segment (M3). Initial
segment extension is the simplicity relation, \emph{not} the numerical order
(M2). This distinction returns in \cref{found:sub:announcement}, where the
birthday-precedence relation of a recent research announcement must be kept
apart from sign-prefix simplicity.
\end{remark}

\begin{remark}[The ordinal embedding is not an arithmetic homomorphism]\label{found:rem:ordinals}
The ordinal embedding $\alpha\mapsto{}$(all-plus sequence of length $\alpha$)
preserves the canonical order but does not identify ordinal arithmetic with
field arithmetic: in $\No$, $1+\omega=\omega+1$ by commutativity, while in
ordinary ordinal arithmetic $1+\omega=\omega$ (M2, M3). The positive fact is
that on embedded ordinals, surreal addition and multiplication agree with the
\emph{natural} (Hessenberg) operations, not generally with the usual
noncommutative ordinal operations (M3) \cite{found:Conway}. A formal library
must use separate operations with explicit comparison results rather than
treating an ordinal coercion as a homomorphism.
\end{remark}

The tradeoff is that addition, multiplication and the normal-form theorem are
less immediate on strings. A productive formalization uses both descriptions and
proves an order-preserving equivalence between numeric-game quotients and sign
sequences. The Mizar development discussed in \cref{found:sec:ecosystem} is an
established example of using more than one presentation rather than requiring
one representation to serve every proof \cite{found:PakKaliszyk}; its existence
is also evidence that a proof assistant need not have a primitive surreal
inductive-inductive constructor to formalize the classical theory (M2).

\subsection{Hahn fields as workspaces, not a circular definition}
\label{found:sub:hahnworkspace}
A Hahn field $\R((t^\Gamma))$ can be constructed from an independently given
ordered abelian group $\Gamma$. Its coefficients, support and convolution have
direct set-sized definitions, which makes Hahn fields excellent algebraic and
analytic workspaces.

The slogan $\No\cong\R((\omega^{\No}))$ has a different status. Its exponent
class is the object being described, the supports are reverse well ordered and
set-sized, and the evaluation map uses the Conway monomial construction. It
summarizes a \emph{normal-form theorem}; it is not a nonrecursive set
definition that solves its own size problem, and it does not define $\No$ as an
unrestricted set of functions from $\No$ to $\R$. A normal form is a set-coded
expression
\begin{equation}\label{found:eq:normalform}
 x=\sum_{\xi<\lambda}r_\xi\omega^{a_\xi},\qquad
 \lambda\in\Ord,\quad r_\xi\in\R\setminus\{0\},\quad
 a_\xi>a_\eta\ \text{for }\xi<\eta,
\end{equation}
whose coding data --- the ordinal $\lambda$, the coefficient function and the
exponent function --- are sets, even though exponent \emph{values} range through
the class $\No$. It is not an arbitrary proper-class function $\No\to\R$.
Existence and uniqueness of \eqref{found:eq:normalform} constitute the Conway
normal-form theorem, not a consequence of size bookkeeping
\cite{found:Conway,found:Gonshor}. A formal development may use this theorem as
a proved bridge; it must not postulate the bridge merely because a Hahn-series
type already exists.

In a class theory an individual class-Hahn series should be stored by a set
support and a set coefficient function on that support; its zero extension to
all surreal exponents is a class function, not a set function with a
proper-class domain. In a larger type universe that zero extension may be an
ordinary function type, but the small-support condition still has to be imposed
(M2). These are different encodings of one size discipline.

\section{ZF, ZFC, and definable classes}
\label{found:sec:zfc}

\subsection{Definable classes require no new objects}
\label{found:sub:definable}
In first-order ZFC, ``$x$ is a surreal sign sequence'' is a formula and class
notation is an abbreviation. A definable class addition is represented by a
formula $\operatorname{Add}(x,y,z)$ together with the theorem that for all
surreal $x,y$ there is exactly one surreal $z$ satisfying it. A class function
$f:\No\to\No$ is represented by a formula $\varphi_f(x,y)$ for which one proves
\[
 \forall x\bigl(\operatorname{Surreal}(x)\Rightarrow
 \exists!y\,[\operatorname{Surreal}(y)\land\varphi_f(x,y)]\bigr).
\]
There is no set of all inputs or of all graph pairs, yet each restriction to a
set of inputs has a set of outputs by Replacement. Statements about finitely
many such operations are ordinary first-order assertions about sets. This
suffices for the classical field construction
\cite{found:Conway,found:Gonshor}.

For an ordinal $\alpha$ the birthday-bounded collection
\[
 \No_{<\alpha}=\bigcup_{\beta<\alpha}\{-,+\}^{\beta}
\]
is a set, by Power Set, Separation, Replacement and Union, and every set of
sign sequences sits in some such stage. The full $\No$ is the definable class
union of these stages, not a set obtained by applying Union to a nonexistent
set of all stages. No inaccessible cardinal is needed for this class
construction.

\subsection{Which axioms do what}
\label{found:sub:whichaxioms}
The relevant axioms are concrete (M3). Power Set and Separation produce bounded
collections of sign sequences and supports. Replacement and Union collect
birthdays, coefficient families and recursively generated sets. Foundation
supports rank and well-founded induction. Infinity provides the ordinary
natural numbers. These are sufficient resources for the usual set-theoretic
construction.

Replacement is not decorative (M2): it bounds the birthdays and collects the
supports of a set family. Power Set is visible in unrestricted function sets and
in the set of all Hahn coefficient functions on a fixed exponent set (M1).
Weakening the set theory requires rechecking these steps; even removing Power
Set has subtleties, since Replacement and Collection cease to have some of
their familiar equivalences \cite{found:ZFCminus}. Zermelo set theory without
Replacement, Kripke--Platek set theory, intuitionistic ZF and constructive ZF
are not interchangeable alternatives: bounded finite-day constructions may
require very little, while collection of arbitrary ordinal-indexed supports and
the full normal-form machinery needs more.

\subsection{Schema versus class quantification}
\label{found:sub:schema}
ZFC has no variables ranging over classes. A theorem for a \emph{specified
definable} $F$ can use its defining formula in Replacement. A theorem
quantified over all class functions is instead interpreted as a schema over
definitions, or stated in an explicit class theory. These formulations are
related, but their quantifiers are not literally identical, and confusing a
schema about definable functions with quantification over all classes can
strengthen or weaken a statement unintentionally (M2).

For example, R (M3) takes a class function $F$ on an infinitesimal thickening,
restricts it to an ordinary set $U$, and collects $F(c)$ for $c\in U$. In ZFC
this is legitimate when the graph of $F$ is given by a formula for which the
required functionality has been proved. An unformalized phrase ``arbitrary
external assignment'' does not supply that formula. \Cref{found:sub:hartogs}
returns to this step, which is the decisive size step in R's automatic Hartogs
support theorem.

\subsection{Choice, representatives and normal forms: a choice audit}
\label{found:sub:choice}
ZF is classical logic without choice; it is not constructive set theory.
Canonical sign sequences, ordinal rank bounds and unique normalized
constructions should not be assigned a global-choice dependency merely because
they are discussed with classes. The deterministic sign-sequence representation
avoids choosing a raw representative from each class. Adding choice permits the
usual classical field theory, ordinary complex analysis, and the selection
arguments used in the supplied support proofs without a separate weak-choice
analysis, which is why ZFC is a practical baseline for the set-sized layers.

The three audits agree on the following itemized audit for this subject.
\begin{itemize}
\item The simplest element of a separated surreal cut is unique, so its
selection is canonical (M2).
\item Unions of supports need no representative choice once normal forms exist
(M2); the union of a specified set of supports, its rational span, and
coefficientwise extension of a fixed linear map are all canonical (M1).
\item Selecting roots of a \emph{set} of polynomials uses ordinary choice (M2).
AG's choice of a complex-linear splitting is a choice concerning a set-sized
vector-space surjection (M1).
\item Choosing representatives or bases across a proper class of unrelated
structures can require global choice (M2). A global family of splittings for
\emph{all} workspaces is a different request and should not be added unless
actually needed (M1).
\item A chosen section of a quotient in Lean often uses classical choice, as
the inspected \texttt{Surreal.out} does. That is a convenience of the
implementation, not a proof that every mathematical account of $\No$ needs
global choice (M1).
\end{itemize}

None of this entails that every proof using global choice essentially needs it.
Conversely, a \emph{particular} proof using arbitrary representatives, a
well-ordering of a coefficient field, or extraction of a decreasing sequence
from a non-well-ordered set may require a choice principle not analysed here
(M3). This report does not claim a sharp reverse-mathematical calibration of
every step in normal forms, real closedness, Hahn-field algebraic closedness or
analytic duality (M1), and identifying the weakest sufficient subsystem would be
a separate investigation (M2).
\section{Class theories and recursion strength}
\label{found:sec:classes}

\subsection{Fix the convention before invoking a theory name}
\label{found:sub:gbconvention}
Names for class theories vary, and all three audits insist that the convention
be fixed before the letters are used. Here \emph{GB} means a two-sorted
G\"odel--Bernays-style theory with set and class variables, elementary class
comprehension (class parameters allowed, quantified variables restricted to
sets), and the usual class-replacement principle; choice assumptions are stated
separately. \emph{GBC} denotes the version with Global Choice and the usual ZFC
set part. NBG is a closely related presentation whose naming conventions differ
over whether choice or global choice is built in; when translating A's ``NBG
with global choice,'' GBC is the intended working convention. None of the
comparisons below depends on treating the letters NBG as a guarantee of an
unstated choice convention.

A class consists of sets, and only sets can be members. The elements of a
proper class are still sets, not other proper classes. Consequently $\No$,
$\SC$, the graph of surreal addition, the graph of a definable global function,
the finite elements $\OO\subset\No$, the infinitesimals $\mm\subset\No$, the
surcomplex unit circle, and the purely infinite part $\Pi$ are legitimate
classes; their individual values remain sets. Elementary class comprehension
allows set quantifiers and existing class parameters but not arbitrary bound
class quantifiers in the defining formula; class replacement ensures that the
image of a set under a class function is a set. These principles are exactly
suited to recording a surreal carrier, its operations, a class function $F$,
and the set of its values on an ordinary domain.

\begin{remark}[Second-sort quantification is not a class of classes]\label{found:rem:secondsort}
The assertion ``for every class homomorphism $\chi:\Pi\to\mathbb T(\No)$'' is a
legitimate second-sort quantification. It is \emph{not} the assertion that
there exists a class whose elements are all such homomorphisms; the latter
would require classes to be elements, hence a higher-order framework or a
coding restriction. This is the proper reading of the global phase
classification in T (M2), and it is the same discipline that
\cref{found:sub:sheaves} applies to a class-valued sheaf.
\end{remark}

\begin{remark}[Conservativity, stated exactly]\label{found:rem:conservativity}
The standard predicative class theories are conservative over the corresponding
set theory for sentences about sets; in particular GBC is conservative over ZFC
for first-order set sentences. This is a relative proof-theoretic fact. It does
\emph{not} make the theories identical: GBC can quantify over classes, and
models can have different collections of classes over the same underlying set
universe. It must not be strengthened to the assertion that every particular
model of ZFC has an expansion to GBC with exactly the same sets, nor to the
assertion that every model of ZFC already has a definable global well-order;
expansion questions are subtler, especially for global choice
\cite{found:Williams,found:WilliamsMin,found:ShulmanSize}.
\end{remark}

\subsection{The recursion design principle}
\label{found:sub:recursionprinciple}
The words ``transfinite'' and ``class-length'' alone do not determine
proof-theoretic strength.

\begin{principle}[Audit the output size as well as the index size]\label{found:prin:recursion}
A recursion with a proper-class domain can still be a local, set-valued
recursion. Conversely, a recursion of ordinal length can ask for a new
\emph{class} at each stage and require a substantially stronger principle. The
length of the index alone does not determine the strength needed.
\end{principle}

For the first kind, ordinary well-founded recursion, rank bounds and class
comprehension often suffice: each stage $\No_{<\alpha}$ is a set, each sign
sequence is a set, and the predecessors involved in constructing the value of a
fixed arithmetic expression can be gathered into a set. A set-like well-founded
relation --- one with only set-many predecessors per element --- is the standard
favourable case \cite{found:HamkinsRecursion}. Operations obtained by agreement
of such local recursions do not require an axiom of unrestricted class
recursion. M3 makes this precise as the only numbered theorem in the group on
recursion strength.

\begin{proposition}[Local recursion and coherent assembly]\label{found:prop:recursion}
Suppose every input $x$ has a set $D_x$ of dependencies, closed under the
predecessors required by a given recursive rule; suppose the predecessor
relation on $D_x$ is well founded, the rule has a unique set-valued output once
predecessor outputs are known, and the recursions on different $D_x$ agree
wherever both are defined. Then the value at each input is obtained by set
recursion. If the data and rule are definable, the assembled operation is a
definable class function; in GB the same argument permits fixed class
parameters.
\end{proposition}
\begin{proof}
Perform well-founded recursion on the set $D_x$; its values and graph are sets
by Replacement. For two dependency sets, uniqueness of recursively computed
values proves agreement by induction on the common required predecessors.
Define the global graph by saying that $y$ is the value at $x$ in such a local
recursion; existence and compatibility make this graph functional. No stage of
this argument constructs a class-valued approximation as an element of a
sequence of classes.
\end{proof}

For a raw game, the transitive closure of its coded option tree is set-sized.
For the positive-support recursions in P and R (M3), the dependency domain is
the explicit set $E^*$ of \cref{found:sub:positivesupport}. These are the
relevant localization certificates. \Cref{found:prop:recursion} is explicitly
\emph{not} a blanket certification of every informal class recursion: more
elaborate surreal operations still require their actual recursive dependency
proofs.

M1 supplies the corresponding audit checklist: a formalization audit should
record (i) the domain of recursion, (ii) the codomain of each stage, (iii) the
decreasing relation, and (iv) the size of the information inspected at a step.

\subsection{ETR, class forcing, and two invalid moves}
\label{found:sub:etr}
\emph{Elementary transfinite recursion} (ETR) concerns elementary definitions
that may recursively produce class-sized sections along class well-orders. Its
general availability is genuine extra strength over GBC in its usual forms.
Iterated satisfaction and truth constructions are the standard examples: at a
stage one may need a relation on all set assignments rather than one set. Even
an $\Ord$-indexed construction can exhibit this difficulty when the stages are
classes, and results on class forcing locate important recursions at
$\ETR_{\Ord}$ \cite{found:ClassForcing,found:GitmanHamkins,found:HamkinsETR}.

Two moves are therefore invalid, and all three audits reject both:
\begin{itemize}
\item ``NBG permits transfinite recursion, therefore every recursion along all
ordinals is allowed'' --- the conclusion is false;
\item ``we must use Kelley--Morse because $\No$ is a proper class'' --- the
strength is not justified by the size of the carrier.
\end{itemize}
One must identify the actual recursion and prove that its dependency domains and
outputs have the advertised sizes. In particular, invoking ETR for a recursion
on one well-ordered Hahn support is unnecessary; conversely, an iteration of
class truth predicates is not made a set recursion by using the same word
``ordinal.''

\subsection{What Kelley--Morse adds, and what it still does not}
\label{found:sub:km}
Kelley--Morse (KM, also MK) permits impredicative class comprehension: class
quantifiers may occur in the defining formulas. This supports stronger class
recursion and satisfaction constructions, and its proof-theoretic and
model-theoretic behaviour differs from GBC's; it is not a conservative
notational extension of ZFC \cite{found:Williams,found:ClassForcing}. It is a
legitimate ambient theory, especially for a project about classes of
structures, iterated truth, or sophisticated class recursions.

But KM is still a theory whose classes have sets as members. It does not make
proper classes into ordinary members, does not grant unrestricted power classes
of all classes, and does not solve the universal-cut contradiction. It does not
automatically provide a \emph{hyperclass} whose elements are all proper-class
functions $\SC\to\SC$; if arbitrary collections of such functions are genuinely
required, one must add another level, move to larger universes, or work with
codes and definable families. Further strengthening by class collection, class
choice, or higher-order objects must also be stated separately; variants of
second-order set theory differ in these principles \cite{found:Williams}.
Nothing in the set-sized ring $\Ag_{n,\Gamma}$ requires KM solely because its
interpretation lies in the surreal class (M1). The present polynomial
workspaces, set-supported Hahn recursions, and restriction of a class function
to an ordinary set do not by themselves require KM (M3).

\subsection{Global choice is optional infrastructure, not a cure for size}
\label{found:sub:globalchoice}
Global Choice supplies a class-wide choice function, for example when a
construction really asks for simultaneous selections across a proper class of
nonempty sets. It is useful for globally chosen representatives and uniformly
selected algebraic constructions. It cannot make a proper class into a set, it
cannot form $\{\No\mid\varnothing\}$, and it cannot create the set of all class
functions (M1). It is also \emph{not} needed merely to pass from
$\forall r\,\exists H_r$ to a proof that fixes one arbitrary $r$ and one witness
for that $r$, nor to define a graph whose output is already uniquely
characterized (M3). These two observations can substantially reduce an
unnecessarily strong axiom ledger.

\section{Universes, reflection, and bounded surreal fields}
\label{found:sec:universes}

\subsection{Grothendieck universes make smallness relative}
\label{found:sub:universes}
A Grothendieck universe $\U$ is a transitive set closed under the usual set
operations and appropriately indexed unions; we assume it contains the ordinary
infinite sets needed here. In a familiar model, $V_\kappa$ for a strongly
inaccessible $\kappa$ is such a universe, and a Tarski--Grothendieck universe
principle provides arbitrarily large ones. A set in $\U$ is called
$\U$-small; finite products, power sets of members, and unions indexed by
$\U$-small families of members remain in $\U$ \cite[\S8]{found:ShulmanSize}.
This is additional foundational strength, not a disguised theorem of ZFC.

Define externally
\begin{equation}\label{found:eq:NoU}
 \No_\U=\{s:\alpha\to\{-,+\}:\alpha<\kappa\},
\end{equation}
the sign sequences whose lengths and data belong to $\U$. From outside, this is
a set; from the small world represented by $\U$ it is the analogue of a proper
class. If $I\in\U$ indexes a family of its elements, the union of their ordinal
lengths is still bounded by an ordinal of $\U$, so their small Conway cut can be
filled inside $\No_\U$. An externally arbitrary subset of $\No_\U$ need not be
an admissible small option set.

Two of the three audits give a numbered proposition here, in \emph{opposite
logical directions}. They are not in conflict: one says what $\No_\U$ does, the
other says what it cannot be, and the merge keeps both. M3's is positive and
inventorial; it justifies the construction.

\begin{proposition}[Why universe relativization avoids the cut contradiction]\label{found:prop:universe}
For $\U=V_\kappa$ with $\kappa$ strongly inaccessible, $\No_\U$ contains a copy
of every ordinal $\alpha<\kappa$, has no upper bound in itself, and fills
separated cuts indexed by members of $\U$. The external set of all its elements
is not an element of $\U$.
\end{proposition}
\begin{proof}
The constantly positive sign sequence gives each $\alpha<\kappa$. Every
$\U$-small family of birthdays has supremum below $\kappa$, by the closure
properties of $\U$, so the ordinary small-cut construction remains below
$\kappa$. If $\No_\U$ were itself in $\U$, its family of all birthdays would be
$\U$-small and hence bounded below $\kappa$, contradicting the presence of
every $\alpha<\kappa$. Finally, appending a plus sign to the sign sequence of a
proposed upper bound $x$ gives a larger element still in $\No_\U$.
\end{proof}

M1's is contrapositive and negative; it is the form that functions as an
interface regression test.

\begin{proposition}[The missing self-cut]\label{found:prop:universecut}
If $\No_\U$ admits surreal cuts of all $\U$-small separated option families,
then $\No_\U$ is not itself $\U$-small.
\end{proposition}
\begin{proof}
Otherwise $L=\No_\U$, $R=\varnothing$ would be an admissible cut, and its value
would be strictly larger than every element of its own carrier, contradicting
\cref{found:prop:allcuts}.
\end{proof}

M2 makes the same point in prose only. One should not say that $\No_\U$ contains
``all surreals'' without the qualification relative to $\U$: enlarging the
universe produces more sign sequences and new scales. When $\U\subset V$, a
formal development should provide an embedding $\No_\U\to\No_V$, preserve
canonical sign descriptions, and prove compatibility of order, arithmetic and
normal forms, rather than identifying the two carriers definitionally; a series
with support merely small relative to $V$ need not correspond to a member of
$\No_\U$ (M1). Arithmetic closure at the smaller level is justified by
performing the classical construction internally in the transitive model
$V_\kappa$, whose individual inputs, output codes and small option families
belong to $V_\kappa$; the interpreted field operations agree with the external
ones on common inputs by uniqueness of their recursive definitions, and the
same method gives the real-closed-field theorem at that level (M1).

\begin{remark}[The cost]\label{found:rem:universecost}
The existence of a strongly inaccessible cardinal is not a theorem of ZFC, and
arbitrarily many universes is stronger still. None of it is needed merely to
discuss the definable proper class $\No$ in ZFC. The inaccessible bound in
\eqref{found:eq:NoU} is a convenient sufficient condition, not a claim of
minimality. A semantic model of a type theory is also not a theorem that its
exact proof-theoretic strength equals that of one familiar large-cardinal
axiom (M2, M3, M1).
\end{remark}

\subsection{Tarski--Grothendieck and Mizar}
\label{found:sub:tg}
Tarski--Grothendieck set theory supplies sufficiently large universe-like sets
for arbitrary starting data, giving a flexible external setting in which one
level's classes become the next level's sets. It is the foundation of the Mizar
tradition, including the surreal formalization discussed in
\cref{found:sub:mizar} \cite{found:PakKaliszyk}. A completed Mizar proof
establishes a result \emph{in that formal environment}; extracting a proof in a
weaker theory requires dependency analysis and does not follow simply because
the informal theorem is usually stated in ZFC (M2). The extra universe strength
is not forced by the elementary assertion that $\No$ is a definable proper class
in ZFC (M1).

\subsection{Feferman reflection without a blanket universe axiom}
\label{found:sub:reflection}
M1 is the only audit that treats reflection frameworks, and its treatment
survives in full. Feferman's reflection-based smallness frameworks, often
written ZFC/S, provide another way to formalize some large-object reasoning: an
added smallness constant satisfies appropriate reflection instances. Their
metatheoretic conservativity relies on the fact that any one proof uses only
finitely many instances. It is \emph{not} the assertion that ZFC proves the
existence of one transitive set model satisfying all of ZFC in a single
internal satisfaction statement \cite[\S11]{found:ShulmanSize}.

For the present application this is a possible way to organize finitely
specified mathematical work without postulating arbitrarily many inaccessible
cardinals. It is not interchangeable with a Grothendieck universe: closure under
arbitrary externally specified small-domain families may fail, or may require a
definability condition. A normal-form or workspace theorem quantifying over all
such families has to preserve that qualification. Replacing a universe by a
reflection constant while retaining all of its closure rules would lose the
point of the weaker framework.

\subsection{Birthday cutoffs: the negative fact and the positive theorem}
\label{found:sub:cutoffs}
For an ordinal $\lambda$ put $\No_{<\lambda}=\{x:\bd(x)<\lambda\}$, a set in
ZFC. All three audits record the negative fact: an arbitrary cutoff is not a
field. The finite-birthday surreals are exactly the dyadic rationals; they are
closed under addition and multiplication but not inversion, since $3$ has finite
birthday and $1/3$ does not. Choosing an ordinal bound and declaring the
resulting subtype a field is invalid without a closure theorem, and a
formalization must not insert a field instance on a birthday subtype merely
because its ambient type is a field (M2).

M3 alone supplies the positive result, and it is the group's only positive
theorem on birthday fragments. The corrected work of van den Dries and Ehrlich
identifies the cutoffs that \emph{are} fields as the epsilon numbers, the
ordinals satisfying $\omega^\lambda=\lambda$, and establishes strong closure and
elementary-extension properties for the resulting fields
\cite{found:vdDE,found:vdDEerratum}. An epsilon number need not be an
inaccessible cardinal: $\varepsilon_0$ is already one. Hence \emph{no
large-cardinal universe axiom is needed to produce useful set-sized real-closed
surreal fields}. The accompanying erratum must be taken into account when
reusing detailed ordinal support estimates from the original paper --- a
non-claim recorded independently by M3 and M2, the latter adding that exact
ordinal closure bounds for exponential subfields are mathematical theorems from
the literature, not consequences of notation, and must be ported together with
their hypotheses and the published correction
\cite{found:EhrlichKaplanExp}.

What a birthday cutoff does \emph{not} automatically provide is closure under
every family one calls ``small'' in an ambient set theory. Its own full set of
elements cannot have an upper bound inside itself, and the cofinality of the
cutoff and the lengths of the input families become explicit constraints.

\subsection{Closure under a fixed finitary signature}
\label{found:sub:closure}
M2 supplies the following positive closure proposition, together with its
essential caveat.

\begin{proposition}[Closure under a fixed set of finitary operations]\label{found:prop:closure}
Let $A$ be a set of surreal codes and let $\mathcal F$ be a uniformly specified
set-indexed family of finitary class operations with set values on their
domains --- more precisely, one class graph encodes the operations together
with their set indices and finite arities, the operations themselves not being
members of a set. Then there is a set $B\supseteq A$ closed under every
operation of $\mathcal F$ whenever the inputs belong to $B$.
\end{proposition}
\begin{proof}
Put $A_0=A$. Given $A_n$, adjoin every value of every operation of $\mathcal F$
on tuples from $A_n$ lying in its domain, obtaining $A_{n+1}$; finite Cartesian
powers, Replacement and Union make this a set. Let $B=\bigcup_{n<\omega}A_n$.
Every finite tuple from $B$ already lies in one $A_n$, so its output lies in
$A_{n+1}\subseteq B$.
\end{proof}

For field operations include $0,1,+,-,\cdot$ and inversion away from zero; for
canonical exponential and logarithmic closure include those particular
operations. \emph{This proposition does not produce closure under all
set-indexed Hahn sums}: those operations have unbounded arity and their
admissible input families are not one fixed finitary signature. A full Hahn
field is the more suitable workspace when such sums are central.

\subsection{Three distinct closure requirements}
\label{found:sub:threeclosures}
For a field or carrier $F$, M3 distinguishes three requirements that are easily
conflated:
\begin{enumerate}
\item closure under finitely many arithmetic operations and finite algebraic
extensions;
\item closure under specified set-indexed strongly summable families;
\item the availability of a number above or below every member of each
permitted small family.
\end{enumerate}
Real closedness supplies only the first. A fixed Hahn field has a natural
summation operation for suitable supports, but its value group may have no
element above a given unbounded sequence. A universe-relative surreal
construction has the third property only for the lower-universe-small families
built into its interface. This three-way distinction is indispensable for
translating A's full-scale arguments, and it is exactly what
\cref{found:sec:twoseries} exhibits in two sharp examples.

The three devices solve different problems. A birthday cutoff gives a concrete
algebraic subfield with ordinal closure theorems; a Grothendieck universe gives
systematic closure under a broad family of constructions; a Hahn workspace gives
a specified value group and transparent supports. Their carriers may overlap
substantially, but their universal properties are not the same.

\subsection{External countability versus internal properness}
\label{found:sub:countable}
M2 alone records this warning, and it survives in full. A countable model of set
theory can internally believe that its surreal collection is a proper class
while externally its domain and its internal surreal codes are countable. There
is no contradiction: the external enumeration need not be an internal set or
class function available in the model. Similarly, ``every set family has a
bound'' means every family represented by a set \emph{in the relevant model}; it
is not an absoluteness claim about every family visible from a stronger
metatheory. The same applies to normal forms and supports: set membership,
well-foundedness, completeness and the collection of all reals must be
interpreted in the stated ambient model.

\subsection{Universes in Lean are not set-theoretic axioms renamed}
\label{found:sub:leanuniverses}
Lean's universe levels are part of its type-formation rules. They are not
variables ranging over a Lean datatype of inaccessible cardinals. A
set-theoretic model using inaccessible cardinals is a sufficient semantic
account of a hierarchy of universes, but it is \emph{not} a theorem that the
kernel literally assumes the corresponding ZFC large-cardinal sentences (M1).
Exact proof-theoretic comparisons require a specified type theory, universe
rules, inductive principles and metatheory. The practical rule is simpler: never
collapse a larger universe just to remove an inconvenient type mismatch
\cite{found:LeanUniverses}. \Cref{found:sec:types} develops the type-theoretic
side in detail.

\section{Surcomplexification: finite algebra, no new size principle}
\label{found:sec:complex}

\subsection{Pairs with the correct multiplication}
\label{found:sub:pairs}
For an ordered field $F$ define $F[i]$ on pairs by
\begin{align}
 (a,b)+(c,d)&=(a+c,b+d),\label{found:eq:pairadd}\\
 (a,b)(c,d)&=(ac-bd,ad+bc),\label{found:eq:pairmul}\\
 \overline{(a,b)}&=(a,-b),\qquad i=(0,1),\ 1=(1,0).\label{found:eq:conj}
\end{align}
Then $i^2=(-1,0)$, and if $(a,b)\ne(0,0)$ ordered-field positivity gives
$a^2+b^2>0$ and
\begin{equation}\label{found:eq:pairinv}
 (a,b)^{-1}=\left(\frac{a}{a^2+b^2},\frac{-b}{a^2+b^2}\right).
\end{equation}
The field construction needs only an \emph{ordered} field, not yet real
closedness: the denominator is positive whenever the pair is nonzero. All three
audits make this point and prove the following, which we print once.

\begin{proposition}[Finite data complexify safely]\label{found:prop:complex}
The pair construction produces a field from any ordered field $F$, and an
algebraically closed field when $F$ is real closed. Its carrier is a set
whenever $F$ is a set and a class of set-coded pairs whenever $F$ is a class
field of set-coded elements. The same coordinate formulas define the surcomplex
class field from the surreal class field.
\end{proposition}
\begin{proof}
The ring identities are finite polynomial identities, and
\eqref{found:eq:pairinv} is verified by direct multiplication using the nonzero
denominator. The equivalence between real closedness of $F$ and algebraic
closedness of $F[i]$ is the standard real-closed-field theorem. For a class
carrier, each pair and each polynomial coefficient list is still a set-coded
finite object, and the graphs of the operations are class relations defined by
the same formulas; pairing and finite products introduce no new proper-class
membership operation.
\end{proof}

Equivalently, in a set-sized workspace one may form $F[X]/(X^2+1)$, whose
irreducibility follows because an ordered field has no square equal to $-1$.
This is an ordinary ideal quotient; there is no need to form a quotient whose
members are proper-class equivalence classes. The quotient method integrates
well with generic field-extension libraries; the pair method makes conjugation,
real and imaginary parts and the modulus more transparent.

``$\SC$ is algebraically closed'' means: every nonconstant polynomial given by a
finite coefficient list in $\SC$ has a zero in $\SC$. It does not assert the
existence of a set of all such polynomials (M1). Real closedness of $\No$ is a
substantive construction theorem; once available, adjoining $i$ introduces no
new paradox and no class-level choice.

\begin{remark}[Not the product ring]\label{found:rem:productring}
The carrier $F\times F$ with its \emph{default componentwise multiplication} is
a product ring with zero divisors: $(1,0)(0,1)=(0,0)$. Its identity is $(1,1)$,
not $(1,0)$. It is not $F[i]$. A proof-assistant implementation must use a
dedicated structure, a type synonym with carefully transported operations, or
the quadratic-extension construction, and must not reuse the default
product-ring instance. All three audits flag this; M3 states it explicitly as a
proof-assistant trap, and it appears as a row in the regression table of
\cref{found:app:shortcuts}.
\end{remark}

\subsection{Real closed is not algebraically closed}
\label{found:sub:realclosed}
No nontrivial ordered field is algebraically closed, because $X^2+1$ has no root
in it. Thus ``$\No$ is real closed'' and ``$\SC$ is algebraically closed'' must
not be merged. There is likewise no order on $F[i]$ making it an ordered field
with the same multiplication: in an ordered field a nonzero square is positive,
whereas $i^2=-1$. Statements about a ``real'' component of a surcomplex number
mean an element of $\No$ and must not be confused with ordinary $\R$ (M2).

This is also an audit rule for documentation: an abstract or a comment
containing ``algebraically closed surreal field'' does not override the
elementary obstruction, nor does it identify a formally proved theorem. M2 and
M3 both flag exactly this wording in the abstract of the published Mizar
development; see \cref{found:sub:mizar}.

\subsection{The modulus is base-field-valued, not a real norm}
\label{found:sub:modulus}
For $z=(a,b)$ over a real closed $F$ put
\[
 |z|=\sqrt{a^2+b^2}\in F_{\ge0}.
\]
Multiplicativity $|zw|=|z||w|$ follows from the sum-of-two-squares identity. The
triangle inequality reduces to
\begin{equation}\label{found:eq:cauchyschwarz}
 (ac+bd)^2\le(a^2+b^2)(c^2+d^2),
\end{equation}
whose difference is the square $(ad-bc)^2$ (M1); equivalently it is the
two-dimensional Gram identity
\begin{equation}\label{found:eq:gram}
 (ac+bd)^2+(ad-bc)^2=(a^2+b^2)(c^2+d^2),
\end{equation}
which M2 ships as a Lean one-liner proved \texttt{by ring} (uncompiled in the
archive; compiled here, depending on \texttt{propext} and \texttt{Quot.sound}
--- see \cref{found:sub:compiled}). These are order-valued estimates, valid over every
real-closed field at infinite as well as infinitesimal scales, and they need no
reconstruction from Conway cuts.

For surcomplex numbers the modulus takes \emph{surreal} values. It is not a
real-valued norm satisfying the hypotheses of Lean's ordinary
\texttt{NormedField} infrastructure, and replacing an infinite modulus by a real
number would change the notion of error tolerance. Do not install incompatible
order or norm instances merely to reuse classical-analysis lemmas. The geometry
and the topology must be formalized using their actual scalar field.

\subsection{Conjugation, modulus and valuation are three separate structures}
\label{found:sub:valuation}
M3 alone supplies the separator that makes this precise, and it survives here.
In a Hahn workspace the natural additive valuation is
\[
 v\Bigl(\sum_\gamma a_\gamma t^\gamma\Bigr)=\min\supp(a),\qquad v(0)=+\infty,
\]
with an ordered-group-valued target, and it is ultrametric. The modulus is
$F$-valued and satisfies an ordinary triangle inequality. Even the targets
differ. Concretely:
\begin{equation}\label{found:eq:separator}
 2t^\rho\ \text{lies in the valuation ball}\ v(z)\ge\rho,
 \ \text{but not in the order disk}\ |z|\le t^\rho.
\end{equation}
A surcomplex implementation should therefore expose a conjugation map, a
positive quadratic form or modulus, and a valuation as \emph{independent} pieces
of structure. An ordinary complex-number type specialized to $\R$ should not be
assumed to support an arbitrary real-closed coefficient field without checking
its hypotheses.

\subsection{The Hahn complexification isomorphism}
\label{found:sub:hahncomplex}
For a set-sized ordered abelian group $\Gamma$, real and imaginary coefficients
give a natural isomorphism
\begin{equation}\label{found:eq:hahncomplex}
 \R((t^\Gamma))[i]\ \cong\ \C((t^\Gamma)).
\end{equation}
The union of two well-ordered supports in a linear order is well ordered, so a
pair of real Hahn series determines a complex Hahn series; conversely taking
real and imaginary parts only shrinks support. The coefficientwise maps preserve
addition, and the finite convolution formulas preserve
\eqref{found:eq:pairmul}. If $\Gamma$ is divisible, the characteristic-zero
Hahn-field theorem gives algebraic closedness of $\C((t^\Gamma))$
\cite{found:Poonen}, and hence real closedness of $\R((t^\Gamma))$.

This is a local theorem about a \emph{set} field. Embedding that field into
$\SC$ requires the surreal normal-form theorem and compatibility with
$t^\gamma=\omega^{-\gamma}$; \emph{the mere existence of an abstract Hahn field
does not establish the embedding} (M1).
\section{Set-sized workspaces and theorem transfer}
\label{found:sec:workspaces}

\subsection{The localization theorem}
\label{found:sub:localization}
Write $t^\gamma=\omega^{-\gamma}$, as in the supplied manuscripts, so that a
Hahn support is well ordered in the \emph{increasing} $t$-exponent order --- the
reverse of the customary decreasing-powers-of-$\omega$ convention. For a
set-sized divisible ordered additive subgroup $\Gamma\subseteq\No$ put
\[
 F_\Gamma=\R((t^\Gamma)),\qquad K_\Gamma=\C((t^\Gamma)).
\]
The normal-form theorem embeds these fields in $\No$ and $\SC$, and the
classical Hahn-field theorem makes $K_\Gamma$ algebraically closed and hence
$F_\Gamma$ real closed
\cite{found:Conway,found:Gonshor,found:Poonen}. These embeddings and closure
statements are substantive prerequisites for localization, not consequences of
the notation $((t^\Gamma))$.

All three audits prove the localization theorem by unioning normal-form supports
and taking the $\Q$-span. We print the proof once, in M3's three-clause form,
which is the strongest statement any of them reaches.

\begin{theorem}[Workspace localization]\label{found:thm:workspace}
Assume the surreal normal-form theorem, the usual arithmetic compatibility of
its monomials, and the algebraic-closure theorem for Hahn fields with divisible
exponent group and algebraically closed coefficient field. Then:
\begin{enumerate}
\item every \emph{set} $A\subseteq\SC$ is contained in one algebraically closed
set-sized field $K_\Gamma$, for some nonzero divisible set-sized ordered
subgroup $\Gamma\subseteq\No$;
\item any set of such workspaces is contained in a larger one;
\item a nonzero polynomial that splits in one workspace acquires no new roots
in a larger surcomplex extension.
\end{enumerate}
\end{theorem}
\begin{proof}
For (1), take for each $a=x+iy\in A$ the union of the exponent supports of the
normal forms of $x$ and $y$, with signs reversed to match
$t^\gamma=\omega^{-\gamma}$. Replacement over $A$ followed by Union gives a set
$S$ of all required exponents. Inside the additive ordered group of $\No$ put
\[
 \Gamma=\operatorname{span}_\Q(S\cup\{1\})
 =\Bigl\{\textstyle\sum_{j=1}^{n}q_js_j:n\in\N,\ q_j\in\Q,\
 s_j\in S\cup\{1\}\Bigr\}.
\]
Finite lists from a set form a set, and their images under finite sums form a
set, so $\Gamma$ is a set by Replacement. It is nonzero, divisible, and inherits
the order from $\No$, and every normal form used for an element of $A$ has
support in $\Gamma$. The cited normal-form and Hahn-field theorems give the
inclusion $A\subseteq K_\Gamma\subseteq\SC$ with $K_\Gamma$ algebraically
closed.

For (2), take the union of the exponent groups of the given workspaces and
repeat the rational-span construction.

For (3), if $P(X)=a\prod_{j=1}^{n}(X-\alpha_j)$ with $a\ne0$ in one field, the
same equality holds in every extension field. At a root $\beta$ the finite
product is zero, so $\beta=\alpha_j$ for some $j$. There is no additional root
location.
\end{proof}

\begin{remark}[The companion non-claim]\label{found:rem:nocover}
M1 attaches the explicit non-claim that \cref{found:thm:workspace} \emph{does
not} say that a single set field contains all surcomplex numbers, and does not
provide a set of workspaces covering the whole class: such a union would be a
set by Replacement and Union, contrary to \cref{found:prop:proper}. The two
statements are compatible --- clause (2) enlarges a \emph{set} of workspaces,
which is a different quantifier from covering a proper class. M2 adds that the
proof presupposes the normal-form theorem and is not a new elementary proof of
it, and that an abstract ordered group $\Gamma$ has no canonical embedding into
$\No$ merely because the notation $t^\Gamma$ is written; the canonical embedding
above uses the actual inclusion $\Gamma\subseteq\No$.
\end{remark}

\Cref{found:thm:workspace} is the foundational content of P's workspace
proposition (M3), made explicit as a dependency theorem, and it already supports
finite polynomial factorization, finite matrices, Hermite interpolation and the
residue-pairing constructions without passing a proper-class ring to an
ordinary algebraic-geometry functor.

\subsection{Enlargement as a system of compatible maps}
\label{found:sub:enlargement}
For $\Gamma\subseteq\Delta$, extend a coefficient function by zero outside
$\Gamma$ to obtain an injective field homomorphism
\[
 j_{\Gamma\Delta}:K_\Gamma\hookrightarrow K_\Delta .
\]
These maps satisfy identity and composition laws and preserve addition,
multiplication, order on the real part, valuation with the corresponding target
embedding, and coefficient extraction. Strong sums are preserved when the same
summability certificate is transported. The same coefficientwise construction
gives maps on $\Ag_{n,\Gamma}$, formal-series coefficient rings, and finite
polynomial or matrix data (M1).

\begin{remark}[Not a small diagram]\label{found:rem:colimit}
The collection of all set-sized exponent subgroups of the full surreal class is
not a set-indexed diagram. Its ``direct limit'' is therefore not automatically an
instance of a library theorem about colimits of small diagrams (M3). Either the
union is definable in class theory, or one fixes a larger universe in which the
index category is an actual type, or one states compatible local theorems
without forming the colimit as a single object (M1). This prevents a
well-intended categorical reformulation from reintroducing the original size
problem.
\end{remark}

The safe pattern for globalizing a local theorem is: gather the parameters into
a workspace; apply the theorem there; if a conclusion quantifies over a proposed
additional surcomplex point, enlarge again to contain that point and the
original parameters; and use a proved compatibility or descent theorem to
compare the two conclusions. Mere inclusion of carriers is not a proof of such
compatibility.

\subsection{A sound localization argument has two halves}
\label{found:sub:twohalves}
All three audits converge on the same warning, stated most crisply by M2: to
reduce a theorem to a workspace one must first prove that all input data fit in
it, and \emph{then} prove that every construction used by the proof stays there
or specify a controlled enlargement. \Cref{found:thm:workspace} supplies the
first half only.

Finite identities descend particularly well: a polynomial has finitely many
coefficients, a matrix finitely many entries, a quotient-algebra computation
finitely many structure constants, and one common workspace contains all of
them. A support recursion can also be set-sized even if it produces infinitely
many coefficients, provided its entire input and recursion support is a set.
Polynomial factorization stays in a divisible algebraically closed Hahn
workspace.

What does not descend automatically: a statement about \emph{every} point of
$\SC$ is not equivalent to its restriction to one $K_\Gamma$; a
positive-dimensional algebraic variety can acquire new points over a larger
field even when its defining equations do not change; a function defined
separately on different infinitesimal monads can behave outside the chosen
workspace in ways its restriction does not determine; and a radius required to
dominate every member of an unbounded subset of $\Gamma$ may lie outside
$\Gamma$ altogether. A rescaling that dominates every exponent in a coefficient
sequence need not stay in the workspace --- which is exactly
\cref{found:ex:rescaling}.

A safe descent proof therefore names the mechanism: persistence of a polynomial
identity; invariance of a finite root multiset; base change of a finite free
algebra; coefficientwise uniqueness; or an explicitly applicable
model-theoretic transfer theorem. ``The problem has finitely many symbols'' is
not a replacement for that argument when those symbols include quantification
over a full function class. Statements about all sequences, all new analytic
functions or all neighbourhoods should never be globalized by deleting a
subscript $\Gamma$ from a theorem.

Embedding compatibility is algebraic but not topological (M2): in
$\R((t^\Q))$ the sequence $t^n$ tends to $0$ in the intrinsic valuation
topology, while after embedding into the full surreal field the same sequence
does not converge in the full fine topology, whose radii include scales below
all its terms. A library must never synthesize the latter property from the
former.

\subsection{Real-closed-field quantifier elimination as a restricted tool}
\label{found:sub:qe}
M3 alone develops this, and it survives in full. For a fixed finite degree
bound, complex polynomial identities and inequalities can be encoded by pairs of
real coordinates, and root counts with multiplicity can be expressed using a
finite factorization list and finite case distinctions. Quantifier elimination
for real-closed fields then permits transfer of such first-order statements
between real-closed workspaces \cite{found:Tarski}. P's order-circle polynomial
Rouch\'e theorem is a scheme indexed by a fixed ordinary degree bound; its
hypotheses quantify over pairs of field coordinates on an algebraically
specified circle, and the count is encoded by a finite root list rather than by
a quantifier over arbitrary functions or subsets. That makes the transfer
legitimate in $F_\Gamma$ once coefficients, centre and radius have been
localized (M3).

The inventory of what the field language does \emph{not} contain is the point.
It does not distinguish the ordinary integer subset from arbitrary field
elements, and it contains no symbol for
\begin{itemize}
\item birthday, \item simplicity, \item Hahn support,
\item a particular exponential, \item a predicate for ordinary holomorphic
functions.
\end{itemize}
It therefore cannot automatically transfer an assertion about all Taylor
coefficients, all ordinal stages, or all coherent charts. Adding such symbols
changes the theory and requires new arguments. In particular, pure
real-closed-field axioms cannot characterize the canonical surreal simplicity
hierarchy --- the same point made from the axiomatic side in
\cref{found:sec:axioms}.

\subsection{Finite algebra over an algebraically closed base}
\label{found:sub:finitealgebra}
The following is M1's, and it is the only substantial commutative-algebra
theorem in the group. It explains why finite algebraic zero sets do not acquire
new points at larger surreal scales.

\begin{theorem}[Finite algebra over an algebraically closed base]\label{found:thm:finitepoints}
Let $K\subseteq L$ be fields with $K$ algebraically closed, and let $B$ be a
finite-dimensional commutative unital $K$-algebra. Then every $K$-algebra
homomorphism $\phi:B\to L$ has image in $K$, and the finite local lengths of $B$
are preserved after extension to $L$.
\end{theorem}
\begin{proof}
For $b\in B$, multiplication by $b$ is a linear operator on the
finite-dimensional space $B$. Its characteristic polynomial $p_b\in K[T]$
annihilates $b$, by Cayley--Hamilton applied to multiplication and then
evaluated at $1$. Hence $p_b(\phi(b))=0$. The polynomial splits over $K$, and
since $L$ is a field, $\phi(b)$ must equal one of its roots in $K$.

For the length statement, decompose the Artinian algebra as $B=\prod_j B_j$
with each $B_j$ local. Its residue field is a finite extension of $K$, hence
$K$. If $\mm_j$ is its nilpotent maximal ideal, then $\mm_j\otimes_K L$ is
nilpotent and the quotient of $B_j\otimes_K L$ by it is $L$; thus the
base-changed algebra is still local. Moreover
$\dim_L(B_j\otimes_K L)=\dim_K B_j$. Since both residue fields equal the
corresponding base field, these dimensions are the local algebraic lengths. The
zero algebra is a harmless empty-product case with no points.
\end{proof}

\begin{remark}[Finite base change is not global tensor equality]\label{found:rem:finitebasechange}
To apply \cref{found:thm:finitepoints} to AG (M1) one must \emph{first}
establish the specific base-change isomorphism
\begin{equation}\label{found:eq:finitebasechange}
 \Ag_{n,\Delta}/(F)\ \cong\
 K_\Delta\otimes_{K_\Gamma}\bigl(\Ag_{n,\Gamma}/(F)\bigr)
\end{equation}
for the finite complete-intersection deformation under consideration. AG obtains
it by using the same finite ordinary basis and support-controlled multiplication
constants in both workspaces. \emph{Neither \cref{found:thm:finitepoints} nor
AG's theorem asserts $\Ag_{n,\Delta}=K_\Delta\otimes_{K_\Gamma}\Ag_{n,\Gamma}$
for the entire infinite-dimensional analytic ring.} Finite quotient
compatibility is a much more specific statement, and the distinction matters
when choosing a Lean abstraction for workspace extension.
\end{remark}

Once \eqref{found:eq:finitebasechange} is available,
\cref{found:thm:finitepoints} explains why any proposed full-surcomplex zero
already has coordinates in $K_\Gamma$. The theorem needs no global surreal
metric space and no theory of fine-continuous paths; it is generic algebra, and
\cref{found:sub:finitefirst} recommends implementing it before any analytic
topology.

\subsection{Characters and actual analytic fibers}
\label{found:sub:characters}
R (M3) explicitly avoids assuming that an arbitrary algebra homomorphism
commutes with an infinite Hahn sum. Instead it proves finite identities such as
\begin{equation}\label{found:eq:divideddifference}
 H(w)-H(\xi)=\sum_j(w_j-\xi_j)G_j(w)
\end{equation}
inside the coherent ring; a character sending $w_j$ to $\xi_j$ must then send
$H$ to $H(\xi)$ by an ordinary finite homomorphism calculation. This is a model
formalization pattern: \emph{prove a finite identity that carries the intended
semantic information} instead of silently assigning continuity or strong-sum
preservation to every homomorphism. It appears as a row in the regression table
of \cref{found:app:shortcuts}.

Similarly, a finite algebra $K[X_1,\ldots,X_d]/I$ can be treated with ordinary
ideal quotients, finite-dimensional modules and localization; a statement about
$\Spec$ should use that set-sized ring and should not treat the full class field
as an object of a small category of commutative rings without a universe
specification.

\subsection{The residue pairing}
\label{found:sub:residue}
M3 alone records the following, and it is a strong candidate for a generic
reusable theorem before any full-class analytic layer exists. For monic $P$ of
degree $n$ over a field $K$, the perfect residue pairing on $K[X]/(P)$ uses only
finite remainders and a finite Gram matrix. With
\begin{equation}\label{found:eq:residue}
 \lambda_P(H)=[X^{n-1}](H\bmod P),
\end{equation}
the matrix in the power basis has ones on its antidiagonal and zeros on one side
of that antidiagonal; reversing columns makes it triangular, and its determinant
is $(-1)^{n(n-1)/2}$. This argument depends neither on surreal representations,
nor on a valuation, nor even on $K$ being algebraically closed.

By contrast, an initial-polynomial root count needs a splitting field, the
valuation and residue maps, and compatibility of those maps with finite
polynomial operations. Those assumptions must not be silently discharged by a
bare \texttt{Field Surreal} instance; P's explicitly algebraically closed Hahn
workspace is the intended source of the stronger hypotheses (M3).

\section{Hahn support as a logical interface}
\label{found:sec:support}

\subsection{Three independent certificates}
\label{found:sub:certificates}
A Hahn construction requires separate answers to three questions (M3). Is the
data indexed by a permitted set or small type? Is the resulting support well
ordered? Does each coefficient receive only finitely many contributions? None of
these follows from the other two in the generality needed here.

For a set-indexed family $(A_j)_{j\in J}$ the manuscripts define \emph{strong
summability} by
\begin{equation}\label{found:eq:summability}
 S=\bigcup_{j\in J}\supp(A_j)\text{ is well ordered},\qquad
 \{j:\gamma\in\supp(A_j)\}\text{ is finite for every }\gamma,
\end{equation}
and the sum has coefficient $\sum_j[t^\gamma]A_j$, a finite sum after the second
condition is applied. Zero coefficients can be discarded, so the actual support
of the result is a subset of $S$. All three audits state both clauses and stress
that neither implies the other. The second clause cannot be replaced by a lower
valuation bound, and positivity alone cannot replace well-ordering: the positive
set $\{1/n:n\ge1\}\subset\Q$ is not well ordered (M1).

\begin{example}[Two failed simplifications]\label{found:ex:failed}
Each series $t^{1/n}$, $n\ge1$, has singleton support, but the family is not
strongly summable because $\{1/n:n\ge1\}$ is not well ordered in the increasing
exponent order. Conversely the constant family $A_n=1$ has well-ordered union of
supports $\{0\}$, but infinitely many terms contribute at exponent zero. It is
not a formal Hahn sum either.
\end{example}

An ordinary convergent series of complex constants may have a perfectly good
complex sum while failing the second condition in \eqref{found:eq:summability};
that operation belongs to the coefficient field's analysis, not to formal Hahn
summation, and mixing the two requires an explicit additional definition.

\subsection{The positive-support mechanism and \texorpdfstring{$E^*$}{E-star}}
\label{found:sub:positivesupport}
If $E$ is a well-ordered set of strictly positive elements of an ordered abelian
group, let
\[
 E^*=\{e_1+\cdots+e_k:k\in\N,\ e_j\in E\},
\]
including the empty sum. The Hahn--Neumann lemma says that $E^*$ is well ordered
and that each exponent has only finitely many representations by nondecreasing
finite words in $E$ \cite{found:Neumann}. Combined with a well-ordered input
support $S$, this gives a controlled support $S+E^*$.

M3 isolates \emph{two different uses} of this lemma, and the distinction is
load-bearing. It proves that the output expression denotes a Hahn series; and it
proves that recursive coefficient equations really involve finite sums. A proof
that checks only well-ordering has not yet justified the coefficient
calculation; a pointwise finite formula does not guarantee that the union of its
output exponents is well ordered.

For instance, when $B$ has positive support the formal geometric identity
$(1+B)^{-1}=\sum_{k\ge0}(-B)^k$ is justified by the support lemma and
coefficientwise cancellation, with no metric convergence theorem applied. For
AG's positive-support operator $T$ (M1), the useful certificate is not that
valuations tend to increase but that the expansion of $T^rG$ has contributions
indexed by
\[
 s+e_1+\cdots+e_r=\gamma,\qquad s\in S,\quad e_j\in E\subset\Gamma_{>0},
\]
with $S$ well ordered and $E$ positive and well ordered; the Hahn--Neumann lemma
gives the well-ordered union $S+E^*$ and finitely many contributing words at
each exponent. This licenses coefficientwise telescoping for
$\sum_{r\ge0}(-T)^rG$, rearrangement of residues, and comparison with a finite
ordinary parameter family. A proof of a trace identity at a fixed coefficient
should retain the finite contributing data explicitly; it must \emph{not} claim
that there are finitely many exponents below that coefficient, which in a
higher-rank Hahn group can be false.

\subsection{A support-local nonlinear recursion}
\label{found:sub:nonlinear}
M3 alone proves the following, which isolates the common mechanism in P's Hensel
factorization and R's preparation and matrix division. The isolation is useful
because it removes both surreal birthdays and ordinary complex topology from the
recursion proof.

\begin{proposition}[Support-local nonlinear recursion]\label{found:prop:positive}
Let $A$ be a commutative coefficient ring, $V=A^m$, let $L:V\to V$ be an
invertible $A$-linear map, and let $N:V\to V$ be a polynomial map all of whose
monomials have total degree at least two. Let $X$ be a Hahn family in $V$ with
well-ordered positive support contained in $E$. Then
\[
 LY+N(Y)=X
\]
has a unique positive-support Hahn solution $Y$, and $\supp Y\subseteq
E^*\setminus\{0\}$.
\end{proposition}
\begin{proof}
Use well-founded recursion on $E^*\setminus\{0\}$. At an exponent $\nu$ each
nonlinear contribution is a product of at least two already indexed positive
coefficients whose exponents sum to $\nu$, so every exponent in such a product
is strictly less than $\nu$; the Hahn--Neumann lemma makes the coefficient sum
finite. Define
\[
 Y_\nu=L^{-1}\bigl(X_\nu-[t^\nu]N(Y_{<\nu})\bigr).
\]
The constructed family is a Hahn family because its support is a subset of the
well-ordered set $E^*$, and exponents outside that set receive no input and no
contribution from the constructed nonlinear products; thus the equation holds
coefficientwise.

For uniqueness among positive-support solutions, take the least exponent where
two candidate vectors differ; their finitely many component supports have
well-ordered union. A difference in a nonlinear coefficient at that exponent
would require a difference at a strictly smaller exponent, so the nonlinear
differences vanish there, and injectivity of $L$ forces the alleged first
difference to vanish as well.
\end{proof}

\begin{corollary}[Fixed domain]\label{found:cor:fixeddomain}
If $A$ is a ring of holomorphic functions on one ordinary complex domain and
$L^{-1}$ preserves that ring, the same proof keeps every coefficient on the
original domain.
\end{corollary}

That is the crucial fixed-domain property in R (M3). The proof uses set
recursion on $E^*$, finite polynomial convolution and a supplied inverse
$L^{-1}$; it requires no Noetherian completion theorem, no fine-topological
limit, and no ETR axiom. It is the concrete instance of
\cref{found:prop:recursion}.

\subsection{Formal summability is not \texttt{tsum}}
\label{found:sub:tsum}
Two examples are needed here, and both survive, because they fail in different
ways.

\begin{example}[The geometric series that never enters the neighbourhood]\label{found:ex:geometric}
Let $\Gamma=\Q+\Q\omega$ and consider $1+t+t^2+\cdots$. It is a valid Hahn sum,
equal to $(1-t)^{-1}$. But the tail after the $N$-th partial sum has valuation
$N+1<\omega$, so the partial sums never enter the valuation neighbourhood
demanding an error of valuation greater than $\omega$; they do not converge to
that sum in this higher-rank valuation topology. AG uses this example explicitly
to rule out unjustified numerical-contraction arguments (M1).
\end{example}

\begin{example}[Natural-number iteration need not reach the required scale]\label{found:ex:archimedean}
Let the exponent group contain $1$ and an element $\Omega$ greater than every
integer. The strongly summable series $H=\sum_{n\ge0}t^n$ has a remainder
beginning at $t^{N+1}$ after $N$ terms, and that remainder is \emph{not} smaller
than $t^\Omega$ for any ordinary $N$. So finite partial sums need not converge
in a topology whose neighbourhoods include that scale (M3).
\end{example}

\begin{remark}[The Archimedean-argument warning]\label{found:rem:archimedean}
An argument asserting that $n\delta$ eventually exceeds every group element is
an \emph{Archimedean} argument and fails in \cref{found:ex:archimedean} (M3).
This is why P's coprime factor lifting is naturally formulated as well-founded
recursion through a support monoid rather than as an unspecified limit of Newton
iterates, and it is why formal Hahn summability and topological summability need
different interfaces.
\end{remark}

A formal implementation should therefore have a dedicated admissible-family
structure and a coefficientwise sum operation. Mathlib's
\path{HahnSeries.SummableFamily} supplies exactly this kind of infrastructure,
separately from topological summability; see \cref{found:sub:mathlibhahn}.
Choosing a default value for a sum outside its admissible domain, as some APIs
do for totality, does not justify the intended sum equation without its support
or positive-order hypothesis (M1).

\subsection{The full surreal exponent type needs a further restriction}
\label{found:sub:toolarge}
At a fixed type-theoretic level, a coefficient function on the entire surreal
type can have a well-ordered support that is too large to be a normal form for a
surreal at that same level. The support must be \emph{small relative to the
universe used to construct those surreals}, not merely a subtype of their larger
carrier. The direct set-theoretic analogue is the proper-class expression
\begin{equation}\label{found:eq:toolarge}
 \sum_{\alpha\in\Ord}t^\alpha,
\end{equation}
which has a well-ordered class of exponents, each receiving one coefficient, and
is nevertheless not a surreal normal form because its support is not a set. Well
ordering and finite coefficient dependence do not repair the missing size
condition.

M2 states the type-theoretic form of the same obstruction: a generic
\texttt{HahnSeries Surreal R} with no small-support restriction admits series
too large to represent a surreal at the chosen level, since the family of
singleton monomials indexed by every ordinal of the smaller universe has
individually small disjoint supports whose union is not small. \emph{A proof
that the union is well ordered does not prove its sethood.} That is exactly why
a small-support subtype is needed, and the inspected Lean library does introduce
one; see \cref{found:sub:hahnlayer}.

\section{Two series over \texorpdfstring{$\C((t^\Q))$}{the field C((t\^{}Q))} with opposite signs}
\label{found:sec:twoseries}

\subsection{The two series, side by side}
\label{found:sub:sidebyside}
This section contains the two sharpest examples in the report. They live over
the \emph{same} field $K=\C((t^\Q))$, with ordered real part
$F=\R((t^\Q))$, and they differ only in the sign of an exponent:
\begin{align}
 A(X)&=\sum_{n\ge0}t^{\mathbf{-}n^2}X^n,
 \label{found:eq:minusseries}\\[0.3em]
 f(z)&=\sum_{n\ge0}t^{\mathbf{+}n^2}z^n.
 \label{found:eq:plusseries}
\end{align}
Their conclusions are opposite in every respect.
\Cref{found:eq:minusseries} is strongly summable at \emph{no} nonzero argument
of $K$, because its term valuations $-n^2+nq$ are unbounded \emph{below} in
$\Q$; it becomes summable only after the value group is enlarged.
\Cref{found:eq:plusseries} \emph{is} strongly summable everywhere in $K$, is
Hahn-coherent at every positive radius $r\in F$, and is nonpolynomial.

\begin{center}\small
\begin{tabular}{@{}P{0.20\textwidth}P{0.34\textwidth}P{0.37\textwidth}@{}}
\toprule
 & $A(X)=\sum t^{-n^2}X^n$ & $f(z)=\sum t^{+n^2}z^n$\\
\midrule
Field & $K=\C((t^\Q))$ & $K=\C((t^\Q))$\\
Exponent sign & minus & plus\\
Term valuations at $v(h)=q$ & $-n^2+nq$, unbounded \emph{below} & $n^2+n\sigma$, unbounded \emph{above}\\
Strongly summable in $K$? & at no nonzero argument & at every argument\\
Repair & enlarge to $\Delta=\Q+\Q\omega$ & none needed\\
Conclusion & fixed workspace too poor to \emph{evaluate} & fixed workspace too poor to \emph{rule out}\\
Audits & M1 and M2 & M3\\
\bottomrule
\end{tabular}
\end{center}

\noindent\textbf{Read this before reading the proofs.} These are \emph{two
different series}, not two variants of one example. The sign of the exponent is
the whole difference, and it reverses every conclusion. A reader who conflates
them will come away believing that this report contradicts itself, or will
quote the wrong one. \Cref{found:sub:complementary} states the relation between
them explicitly.

\subsection{The rescaling obstruction, with the minus sign}
\label{found:sub:rescaling}
M1 and M2 both prove the following, by the same valuation computation. We print
it once.

\begin{proposition}[Universal rescaling fails in a fixed Hahn field]\label{found:ex:rescaling}
Over $K=\C((t^\Q))$ the formal series $A(X)=\sum_{n\ge0}t^{-n^2}X^n\in K[[X]]$
cannot be evaluated by strong summation at \emph{any} nonzero $h\in K$.
\end{proposition}
\begin{proof}
Let $q=v(h)\in\Q$. The leading exponent of $t^{-n^2}h^n$ is $-n^2+nq$, and the
successive difference is $q-2n-1$, which is eventually negative. So for all
sufficiently large $n$ these exponents strictly decrease, without a lower bound
in $\Q$. Each lies in the support of its term, so the union of term supports
contains an infinite strictly decreasing sequence and is not well ordered. This
violates the first clause of \eqref{found:eq:summability}.
\end{proof}

The conclusion is \emph{negative}, and it is about A: the small-scale evaluation
theorem of the analysis manuscript does not survive restriction to a fixed Hahn
workspace, and no nonzero radius in this field has the universal small-scale
property asserted for the full class. In the full surreal field one may instead
choose an exponent $\Delta$ dominating the set of coefficient supports and use
$\rho=t^\Delta$; the separated-band proof in A uses precisely that additional
freedom (M2). A formal port to $K_\Gamma$ must either assume an appropriate
dominating exponent in $\Gamma$ or enlarge the workspace. It must not assert the
full-class theorem uniformly inside every fixed Hahn field, and no formalization
should weaken A's rescaling theorem to the assertion that every series has a
usable nonzero argument in whatever field happened to contain its coefficients
(M1).

M1 alone supplies the \emph{positive half}, and it survives here in full.

\begin{proposition}[The enlarged group repairs it]\label{found:prop:rescalingpositive}
After enlarging the exponent group to $\Delta=\Q+\Q\omega\subset\No$,
evaluation of $A$ at $h=t^\omega$ \emph{is} strongly summable.
\end{proposition}
\begin{proof}
The exponents are $n\omega-n^2$, with consecutive difference
$\omega-(2n+1)>0$ since $\omega$ exceeds every integer. They therefore strictly
increase, and exactly one term sits at each exponent, so both clauses of
\eqref{found:eq:summability} hold.
\end{proof}

M2 proves only the negative half. \Cref{found:prop:rescalingpositive} is what
makes \cref{found:ex:rescaling} a genuine \emph{obstruction justifying A's
permission to enlarge the workspace}, rather than a defect in A.

\subsection{All-scale polynomial rigidity, in support-theoretic form}
\label{found:sub:rigidityargument}
M3 reconstructs the core of A's polynomial-rigidity argument in a form small
enough to audit. Suppose a function has Taylor coefficients $a_n$ at zero and a
coherent chart at every positive surreal radius $r$; for each $r$ the supports
of $a_nr^n$ lie in one well-ordered chart support. If infinitely many $a_n$ are
nonzero their valuations form a \emph{set}, so one may choose $\Delta>0$ with
\begin{equation}\label{found:eq:deltachoice}
 |v(a_n)|<\Delta/4\qquad(a_n\ne0).
\end{equation}
At $r=t^{-\Delta}$, successive nonzero coefficients indexed by $m>n$ satisfy
\begin{equation}\label{found:eq:rigiditycomputation}
 v(a_mr^m)-v(a_nr^n)=v(a_m)-v(a_n)-(m-n)\Delta<0 .
\end{equation}
Their valuations form an infinite strictly decreasing sequence inside one
well-ordered support: a contradiction. Coefficientwise identity then identifies
the function with its finite Taylor polynomial on every chart.

The foundational requirement is the ability to dominate this \emph{whole set} of
valuations. The proof does not need a global choice of charts for every radius
at once --- it fixes the one radius constructed from the hypothetical infinite
Taylor support --- but it does need that radius to be available in the same
ambient full-scale structure. A universe-relative formulation retains the
argument for small coefficient families. \emph{A fixed arbitrary Hahn field
cannot supply it}, and the next subsection shows why not.

\begin{remark}[Two unrelated theorems are called ``rigidity''] \label{found:rem:rigidity}
The word ``rigidity'' names two different things in the merged literature, and
this report never uses the bare word. The result reconstructed here is
\emph{all-scale polynomial rigidity}, an analytic statement taken from A. A
different result, \emph{sign-tree categoricity}
(\cref{found:prop:signtree}), is a foundational statement due to M2 about the
uniqueness of realizations of the sign-tree specification. They share no
content. Each occurrence of one carries a pointer to the other.
\end{remark}

\subsection{The nonpolynomial internally coherent example, with the plus sign}
\label{found:sub:internal}
M3 alone supplies the following, and it is the sharpest mathematical asset in
the group.

\begin{example}[Internal all-scale coherence need not force a polynomial]\label{found:ex:internal}
Let $K=\C((t^\Q))$, with ordered real part $F=\R((t^\Q))$. Define, using strong
Hahn summation,
\begin{equation}\label{found:eq:nonpoly}
 f(z)=\sum_{n\ge0}t^{n^2}z^n\qquad(z\in K).
\end{equation}
This is a \emph{nonpolynomial} function possessing a Hahn-coherent chart on the
whole ordinary complex plane at every positive radius $r\in F$.
\end{example}

\begin{proof}[Support proof]
Write a nonzero input as $z=c\,t^\sigma(1+\eta)$ with $\sigma\in\Q$,
$c\in\C^\times$, and $\eta$ of positive support. The base exponents
$n^2+n\sigma$ tend upward and are strictly increasing for all sufficiently large
$n$; the finite initial portion causes no obstruction. They therefore form a
well-ordered set. Moreover each base exponent comes from at most \emph{two}
indices, since $n^2+n\sigma=m^2+m\sigma$ with $n\ne m$ forces $n+m=-\sigma$.
Finite binomial expansion of $(1+\eta)^n$ has support contained in the positive
monoid generated by $\supp(\eta)$. The support lemma of
\cref{found:sub:positivesupport} therefore proves strong summability of
\eqref{found:eq:nonpoly}; the value at zero is $1$.

Apply the same argument with $z=r$ to see that the coefficient family
$t^{n^2}r^n$ is strongly summable. Regrouping
\[
 \sum_{n\ge0}t^{n^2}r^nw^n
\]
by Hahn exponent gives an \emph{ordinary polynomial in $w$} at each exponent,
because only finitely many indices contribute there. These entire coefficient
polynomials have one well-ordered support and define the claimed chart, and
evaluation at finite $w$ agrees with $f(rw)$ by the same finite-word support
calculation.

Finally, at $r=1$ the coefficient at exponent $n^2$ is $w^n$. A finite-degree
polynomial over $K$ has a uniform bound on the degrees of all of its ordinary
coefficient polynomials, so coefficient uniqueness on ordinary points rules out
equality with $f$. Hence $f$ is nonpolynomial.
\end{proof}

\begin{remark}[Why the rigidity proof cannot run here]\label{found:rem:whynotransfer}
The argument of \cref{found:sub:rigidityargument} needs a $\Delta$ in the value
group satisfying \eqref{found:eq:deltachoice} for \emph{all} $n$. Here
$v(a_n)=n^2$, which is unbounded in $\Q$. No such $\Delta$ exists in $\Q$. The
value group $\Q$ simply contains no exponent dominating all the values $n^2$.
\end{remark}

\subsection{Exact scope of the plus-sign example: a sharpness result, not a
refutation}
\label{found:sub:scope}
This point must not be compressed. \Cref{found:ex:internal} is a
\textbf{sharpness result}. It does \emph{not} refute the all-scale polynomial
rigidity theorem of the companion report \emph{Surcomplex Analysis: germs,
coherent Hahn sections, and infinitesimal zero geometry}, held in this
repository at \texttt{docs/surcomplex/analysis/}
\cite{found:CompanionAnalysis}. The two statements have different hypotheses and
both are true:

\begin{itemize}
\item That companion theorem is stated for \textbf{class functions}
$f:\SC\to\SC$ on the whole of $\No[i]$. Its hypothesis quantifies over a
\emph{proper class} of multiplicative scales, and that proper class is precisely
what supplies the dominating exponent $\Delta$ that
\eqref{found:eq:deltachoice} requires. Over the full surreal scale system the
dominating exponent always exists.
\item \Cref{found:ex:internal} is stated over a \textbf{fixed value group},
namely $\Q$. There the dominating exponent does not exist
(\cref{found:rem:whynotransfer}), and the conclusion genuinely fails: a
nonpolynomial function satisfies the internal all-scale condition.
\end{itemize}

So the example shows that the all-scale rigidity theorem is \emph{specific to
the full-class scale system} and does not transfer to a fixed Hahn field. That
is its entire content. M3 states this itself, in its own words: the example
``does not refute A: the value group $\Q$ has no exponent dominating all the
values $n^2$, and A explicitly quantifies over the larger surreal scale
system.'' M3 further records, and we carry verbatim, that the example ``is an
elementary comparison supplied here,'' that it ``is not attributed to an
unpublished theorem in the attached manuscripts,'' and that \emph{no priority
claim and no machine-verification claim is made for it}. Nothing in the present
report attaches a priority claim to it either, and nothing here should be read
as a correction to the companion report; that report is cited by title and
repository path, and none of its internal labels is referenced.

\subsection{Why the two are complementary, not contradictory}
\label{found:sub:complementary}
\Cref{found:ex:rescaling} and \cref{found:ex:internal} concern \emph{different
series}. One has exponent $-n^2$ and the other $+n^2$, over the same field, and
the sign of the exponent is the whole of the difference between them. They are
not variants of one example, and neither is a strengthening or weakening of the
other.

They are two \emph{different non-transfer theorems} about the same fixed field,
pointing in complementary directions:

\begin{itemize}
\item \Cref{found:ex:rescaling} says a fixed workspace may be too poor to
\textbf{evaluate} a given series. The exponents run downward out of the group,
and the series has no admissible argument at all. The repair
(\cref{found:prop:rescalingpositive}) is to enlarge the value group.
\item \Cref{found:ex:internal} says a fixed workspace may be too poor to
\textbf{rule out} a nonpolynomial all-scale-coherent function. The exponents run
upward out of reach of any single dominating element of the group, so the
rigidity argument has no $\Delta$ to work with, and the classification theorem
that holds over the full class fails over $\Q$.
\end{itemize}

Both are consequences of one fact about $\Q$: it is a value group in which
neither an unbounded-below family nor an unbounded-above family can be
dominated. Both are negative results about \emph{transfer}, not about the source
manuscript. Neither refutes A. Together they are the sharpest available evidence
for the report's recommendation that summability and ambient scale system be
separate explicit hypotheses, never connotations of the word ``surreal.''
\section{Topology: full-class, intrinsic, coefficientwise}
\label{found:sec:topology}

\subsection{A class neighbourhood relation is enough}
\label{found:sub:neighbourhood}
A's fine topology uses balls
\[
 B(a,r)=\{z\in\SC:|z-a|<r\},\qquad r\in\No_{>0},
\]
with radii ranging over positive \emph{surreals}. An individual ball is
generally a proper class. In class set theory one defines openness of a given
class $V$ by a formula --- for every $a\in V$ there is $r>0$ with
$B(a,r)\subseteq V$ --- so there is no need to form a set, or a class with class
members, of all open subclasses. Continuity and local analytic statements are
formulated through the same neighbourhood relation. On a set-sized workspace, of
course, an ordinary topology is a set and the usual definition applies (M2).

In Lean a universe-relative surreal carrier is an actual type, so an ordinary
\texttt{TopologicalSpace} on that type is available in principle, and its
subsets are Lean predicates. This does \emph{not} mean that every such subset is
small in the lower universe used for the cut construction. The apparent
discrepancy disappears once the two meanings of ``set'' are distinguished (M3).

\subsection{The degeneracy theorem, in its strongest form}
\label{found:sub:discrete}
All three audits prove that small subsets are discrete. M1's version has two
clauses (subsets, convergent nets); M2 and M3 both add a third clause about
Cauchy nets, proved by collecting the set of nonzero \emph{pairwise} distances
rather than distances to a limit. We print the strongest statement, with the
Cauchy clause credited to M2 and M3, and their common proof once.

\begin{theorem}[Small subsets and small-index nets are discrete]\label{found:thm:discrete}
Assume \cref{found:lem:bounds} for the permitted small families, and give the
carrier the ball topology with radii in the positive part of its own scalar
field. Then:
\begin{enumerate}
\item every set-sized (respectively permitted-small) subset is closed and
discrete;
\item every convergent net with a set-sized (permitted-small) directed index
set is eventually \emph{equal} to its limit;
\item every Cauchy net with such an index set is eventually constant.
\end{enumerate}
\end{theorem}
\begin{proof}
Let $A$ be such a subset and $p$ a point. The nonzero distances
$D_p=\{|p-a|:a\in A,\ a\ne p\}$ form a set of positive surreals, so by
\cref{found:lem:bounds} choose $\delta>0$ smaller than each of them. Then
$B(p,\delta)\cap A$ is empty if $p\notin A$ and is $\{p\}$ otherwise. This
proves closedness and discreteness.

For a convergent net, its range is a set; apply the same construction to the
nonzero distances of its terms from the proposed limit. Eventual membership in
the resulting ball forces eventual equality.

For a Cauchy net, instead collect the nonzero distances $|z_j-z_k|$ between all
pairs of its terms; that is still a set. Choosing $\delta$ below all of them,
the Cauchy condition for tolerance $\delta$ makes a tail have all pairwise
distances zero, hence constant.
\end{proof}

This is the precise size-dependent argument underlying A's topological
degeneracy theorem, with its essential size quantifier retained (all three
audits). A Lean theorem quantifying over arbitrary \texttt{Set Surreal} without
a smallness hypothesis would say something stronger and false (M2). It does
\emph{not} make the whole full-class topology discrete: every ball contains a
point distinct from its centre, for example $a+r/2$, and the carrier's own point
set is not one of the sets to which the theorem applies. At a universe-relative
level the correct statement is: every \emph{lower-universe-small} subset is
discrete in the full topology of the larger carrier. It is false with the
smallness qualifier removed, since the carrier is a subset of itself and has no
isolated points. This is one of the most important statement changes required in
a faithful Lean translation (M3).

\subsection{Two sharp boundaries: the index and the ambient}
\label{found:sub:boundaries}
\Cref{found:thm:discrete} has two independent sharpness statements, from two
different audits. Neither is implied by the other, and both are needed.

\begin{example}[The larger-index net counterexample]\label{found:ex:largerindex}
Let $D=F_{>0}$ and direct it by \emph{reverse} numerical order: $r\preceq s$
means $s\le r$. Define $x_r=r$. For any $\epsilon>0$, once $r\preceq s$ with
$r=\epsilon/2$, one has $0<x_s=s\le\epsilon/2<\epsilon$. Hence this net
converges to zero and is never zero.
\end{example}

In class theory $D$ is a proper class and the net is class-indexed. In Lean,
with $F=\No_u:\Type_{u+1}$, $D$ is a perfectly legitimate type in the larger
universe but is not a permitted $u$-small index type. There is no conflict with
\cref{found:thm:discrete}. The engineering consequence is stated by M2 exactly
as it should be carried: \emph{replacing ``sequence'' by ``net'' does not repair
a missing size annotation unless the index universe is also specified.} This
example is the exact complement to the discreteness theorem on the index side.

\begin{proposition}[Intrinsic and full-class subspace topologies differ]\label{found:prop:twotopologies}
Let $\Gamma$ be nonzero, divisible and set-sized, and let
$K_\Gamma\subset\SC$ be its workspace. Its subspace topology induced by the
full-class fine neighbourhoods is discrete. Its intrinsic field-valued ball
topology, with radii in $(F_\Gamma)_{>0}$, is not discrete.
\end{proposition}
\begin{proof}
The first conclusion follows from \cref{found:thm:discrete}, since $K_\Gamma$ is
a set. For the second, write $F_\Gamma=\R((t^\Gamma))$, so
$K_\Gamma=F_\Gamma[i]$; every intrinsic ball $B(0,r)$ with
$r\in(F_\Gamma)_{>0}$ contains the nonzero element $r/2$. Hence zero is not
isolated intrinsically.
\end{proof}

This is not a contradiction about subspace topology. The intrinsic field does
not contain the arbitrarily small \emph{external} radii used to isolate its
points --- which is the concrete content of the earlier remark that the $\delta$
of \cref{found:lem:bounds} need not lie in any previously specified subfield.
The intrinsic valuation topology is likewise nondiscrete for a nontrivial
ordered exponent group, since there are monomials of arbitrarily high valuation
\emph{within that group}; indeed $t^n\to0$ intrinsically in $\R((t^\Q))$ but not
finely (M2). An order embedding of workspaces preserves algebra and strong sums
but need not preserve the intended ambient topology in the naive subspace sense.
Formal theorems about continuity, limits or derivatives must specify which
topology is meant before invoking a library result.

The consequence for Lean is direct (M1): if a native type is used as the whole
carrier and every Lean subset is declared to satisfy the small-subset theorem,
taking the subset to be the whole carrier forces discreteness. To express the
full-class theorem at level $u$ one must restrict it to subsets with $u$-small
underlying type while the ambient $\No_u$ or $\SC_u$ lives at level $u+1$.
Alternatively one can express the theorem inside \texttt{ZFSet}, where an
internal set and a Lean predicate class are already distinguished.

\subsection{No countable ball basis, hence no real metric}
\label{found:sub:nometric}
M2 alone supplies the following, and its engineering consequence is the most
directly actionable statement in this section.

\begin{lemma}[No countable ball basis is coinitial]\label{found:lem:nometric}
The full fine topology is not discrete --- every ball around zero contains half
its radius --- yet its small subsets are discrete. No countable ball basis at
zero is coinitial.
\end{lemma}
\begin{proof}
A countable local basis would have a positive radius $r_n$ inside each of its
members. The radii $\{r_n:n\in\N\}$ form a set, so by \cref{found:lem:bounds}
choose $\delta$ with $\delta<r_n/2$ for every $n$. Then $r_n/2$ lies in the
$n$-th basic ball but $r_n/2\notin B(0,\delta)$, so no basic ball is contained
in $B(0,\delta)$.
\end{proof}

Ordinary real metrics always have a countable local basis, so \emph{no ordinary
real metric represents the fine topology}. The precise engineering consequence:
Mathlib's metric typeclasses are the wrong target, while
\texttt{TopologicalSpace} and filters are correct, since they are not restricted
to real metrics. This also prevents importing the ordinary intuition that
pointwise isolation on all small samples makes the whole space discrete.

\subsection{The fine derivative, over the actual scalar field}
\label{found:sub:derivative}
M2 gives the derivative in the only form that is correct here:
\begin{equation}\label{found:eq:derivative}
 \forall\epsilon\in F_{>0}\ \exists\delta\in F_{>0}\ \forall h,\quad
 0<|h|<\delta\ \Longrightarrow\
 \left|\frac{f(z+h)-f(z)}{h}-d\right|<\epsilon,
\end{equation}
where $\epsilon$ and $\delta$ range over the \emph{actual scalar field}
$F_{>0}$, not over $\R_{>0}$, $f$ is defined on a neighbourhood of $z$, and the
chosen ball of increments is required to stay in the domain. The derivative is
unique by ordered-field estimates. This definition requires no real-valued norm,
no sequence limit and no derivation on the scalar field.

\begin{remark}[Four notions that must be distinct types]\label{found:rem:fourtypes}
A function constant on each infinitesimal coset can be globally nonconstant and
have derivative zero everywhere, so the usual implication ``zero derivative on a
connected domain implies constant'' cannot be reused with an ordinary connected
base domain and a class-sized surreal thickening. A supplied analysis
manuscript imposes common-domain Hahn coherence precisely to obtain stronger
global conclusions. In a formalization, \texttt{FineHasDerivAt}, a formal-series
germ, a Hahn-coherent datum and a scalar-field derivation must be \emph{distinct
types or predicates} (M2). Proving implications among them is useful;
identifying them by notation is not.
\end{remark}

\subsection{Ordinary complex contours live in a third structure}
\label{found:sub:contours}
The ordinary complex parameter domain $U\subseteq\C^d$ retains its usual
topology and analytic structure. A Hahn-coherent datum is a family of ordinary
holomorphic functions on this common domain indexed by a well-ordered support;
its contour integral in A or R integrates ordinary coefficient functions and
then forms a Hahn series. It is not an integral along a nonconstant
fine-continuous map from a real interval into the full surcomplex class.

M3 makes the reason a corollary of \cref{found:thm:discrete}, and this survives
as a distinctive item.

\begin{corollary}[No nonconstant fine-continuous paths]\label{found:cor:nopaths}
The range of any real-parameter map into $\SC$ is a set, hence discrete by
\cref{found:thm:discrete}. A continuous image of a connected real interval in a
discrete space has at most one point. Therefore there are no nonconstant
fine-continuous paths, and the contour formalism necessarily lives in a third
structure, distinct from both the full fine topology and the intrinsic
$K_\Gamma$ topology.
\end{corollary}

Similarly, the global fine-analytic logarithm and the nonunique fine-analytic
exponentials described in A do not contradict classical complex analysis; their
local neighbourhoods and global coherence conditions differ from the ordinary
ones. A foundational formalization should preserve this distinction rather than
silently replacing A's functions with ordinary holomorphic maps into a normed
field.

\subsection{Set-coded germs versus quotients of class-function graphs}
\label{found:sub:germs}
M1 alone treats the representation question for germs. A describes fine-analytic
germs using equivalence of functions on sufficiently small fine balls. As a
semantic description this is natural; taken literally as a quotient whose
representatives are arbitrary class-function graphs, it is not automatically an
ordinary NBG class quotient of sets.

A safe implementation uses \emph{descriptions}. A formal coefficient sequence
$(a_n)$, a centre, and any needed admissibility data form a set code. Evaluation
is a definable class relation. Equality of two descriptions is eventual equality
of their evaluations on a smaller ball, expressed by a formula quantifying over
points and radii. One then uses canonical formal series if uniqueness of
expansion has been proved, or Scott codes (\cref{found:prop:scott}) if a more
general quotient is needed. No equivalence class of arbitrary proper-class
graphs is used as an element.

A's theorem identifying fine-analytic germs with $\SC[[X]]$ can support this
representation, but it must not be used circularly to justify the objects needed
to prove that very theorem: begin with formal-series codes and their evaluation
relation, prove the local uniqueness statement, then identify them with the
intended semantic germs. The resulting $\SC[[X]]$ is a \emph{class} of set-coded
sequences unless a workspace or universe has been fixed.

For a fixed class function $f:\SC\to\SC$, the assertion
\[
 \forall r\in\No_{>0}\ \exists s\;
 [\,s\text{ codes a suitable chart for }f(r\,\cdot)\,]
\]
quantifies over set-valued radii and set-coded charts with $f$ as a class
parameter. This is compatible with elementary class comprehension when the chart
condition has only set quantifiers; it needs no set of all radii or all charts.
If a proof needs a simultaneously selected class function $r\mapsto s_r$, that is
a further choice or definability question. A's all-scale argument can reason
about a chart at a specified radius without first choosing all charts globally,
and a formal development should retain that economical existential formulation.

\section{Coefficient rings, sheaves, and Hartogs coherence}
\label{found:sec:sheaves}

\subsection{Four rings that must not be merged}
\label{found:sub:fourrings}
M1 alone distinguishes the following four objects, and the distinction survives
in full with both of its worked separators.
\begin{align}
 \HH_\Gamma(U)&=\OO(U)((t^\Gamma)),\label{found:eq:commondomain}\\
 \Ag_{n,\Gamma}&=\C\{z_1,\ldots,z_n\}((t^\Gamma)),\label{found:eq:germring}\\
 \Fg_{n,\Gamma}&=\C[[z_1,\ldots,z_n]]((t^\Gamma)),\label{found:eq:formalcoeff}\\
 K_\Gamma[[z]]&=\text{formal power series over the Hahn field.}
 \label{found:eq:khahn}
\end{align}
In \eqref{found:eq:commondomain} all coefficient functions are defined on
\emph{one specified} ordinary open set $U$. In \eqref{found:eq:germring} every
coefficient is an ordinary convergent germ, but different coefficients need
\emph{not} have a common radius. In \eqref{found:eq:formalcoeff} no ordinary
convergence is required at all. All three are set-sized for fixed set $\Gamma$.

\emph{First separator.} Taking germs of common-domain families does not produce
\eqref{found:eq:germring}. AG's example
\[
 \sum_{r\ge1}\frac{t^{r\eta}}{1-rz},\qquad\eta>0,
\]
has coefficient poles approaching zero, so there is no common ordinary disk on
which all coefficients are holomorphic; it nevertheless belongs to
$\Ag_{1,\Gamma}$. Formalizing AG by imposing a single uniform analytic norm on
every coefficient would strengthen its hypotheses and lose this example.

\emph{Second separator.} Rearrangement of a double family does not identify
\eqref{found:eq:khahn} with the preceding rings. The identity
\[
 -\sum_{r\ge0}t^{-(r+1)\eta}z^r=(z-t^\eta)^{-1}
\]
is valid in $K_\Gamma[[z]]$, but its combined Hahn support is not well ordered,
so it is \emph{not} an element of $\Ag_{1,\Gamma}$, where $z-t^\eta$ has an
actual infinitesimal zero. A proof assistant should expose this distinction in
the types, not erase it by a convenient coercion.

\subsection{Sheaves: two different questions, two answers}
\label{found:sub:sheaves}
Two audits treat the sheaf question, and they are answering different questions.
Both answers survive, and the framing below is the point.

\medskip\noindent
\textbf{M1: the size of the totality of sections.} A describes a class-valued
sheaf over the ordinary complex analytic base. This can be read in GBC without a
set of all classes. Let $U$ range over ordinary open subsets of $\C$ and let
$s$ be a \emph{set code} for an admissible Hahn section. Define a single class
relation $\operatorname{Section}(U,s)$, and another class relation coding
restrictions; the sheaf axioms are then expressed using set-indexed covers and
compatible set-coded families. Each section code has set-sized support, while
the totality of section codes over a fixed $U$ may be a proper class. One must
not instead take an arbitrary class function as a \emph{member} of a section
set. M1 attaches the non-claim explicitly: this encoding is \emph{a
clarification of a workable reading of A, not a proof of its analytic gluing
theorem}; the latter still requires the coefficientwise overlap argument and A's
exact coherence hypotheses.

\medskip\noindent
\textbf{M3: the shape of an individual section.} A's gluing theorem on a
\emph{connected} ordinary domain obtains a common support by a strong identity
argument: on overlapping connected open sets a holomorphic coefficient that is
nonzero cannot vanish identically on a nonempty open overlap, so exact supports
agree along the connected overlap graph. On a \emph{disconnected} open set A
deliberately defines sections componentwise, and M3 shows this convention is
\emph{necessary}, not a convenience.

\begin{example}[The disjoint-disks necessity example]\label{found:ex:disjointdisks}
Let $U$ be a union of pairwise disjoint ordinary disks $D_n$, and prescribe the
constant section $t^{1/n}$ on $D_n$. Each local datum has singleton support, and
all overlap conditions are vacuous. A definition insisting on \emph{one} globally
well-ordered support for the whole disconnected $U$ would not admit this glue,
since $\{1/n:n\ge1\}$ is decreasing. The componentwise definition does admit it.
\end{example}

For fixed $\Gamma$ one formalizes this as a set- or type-valued sheaf with
\begin{equation}\label{found:eq:componentwise}
 \HH_\Gamma(U)=\prod_{V\in\pi_0(U)}\HH_\Gamma(V),
\end{equation}
the factors on connected components using the global-support convention; the
component index is an ordinary set. Alternatively one may use an appropriate
sheafification and prove agreement on connected domains. For the full surreal
exponent class, use set-coded individual sections and a schema of gluing
assertions, or move to a specified higher universe. One should not invoke an
ordinary category whose objects are literally all proper classes without
supplying its foundations.

\medskip\noindent
\textbf{The relation between the two.} M1's class-relation encoding addresses the
\emph{size} of the collection of sections; M3's componentwise convention
addresses the \emph{shape} of an individual section. M1 says nothing about
disconnected bases; M3 says nothing about the size of the totality of sections.
They are compatible, and M3's necessity example is a constraint binding on any
encoding \emph{including} M1's: an implementer who adopts the
$\operatorname{Section}(U,s)$ relation must still define sections componentwise
over $\pi_0(U)$.

\subsection{Hartogs coherence: where class replacement matters}
\label{found:sub:hartogs}
M3 alone audits R's automatic Hahn--Hartogs theorem, and the analysis survives
in three parts.

\emph{The decisive size step.} R begins with a class function on a full product
of thickenings. Its decisive step is
\begin{equation}\label{found:eq:hartogssupport}
 A=\bigcup_{c\in U}\supp(F(c)),
\end{equation}
where $U$ is an ordinary complex domain and therefore a set. Class replacement
makes the family of values $F(c)$ a set; another replacement and union make $A$
a set. This is exactly the use of class parameters supported by GBC. In ZFC the
same reasoning is a \emph{scheme} for definable $F$ --- see
\cref{found:sub:schema}.

\emph{The analytic step.} Once that is justified, the coefficient family
$(f_\gamma)_{\gamma\in A}$ is set-indexed, and the classical
separate-holomorphy theorem supplies joint holomorphicity. The automatic-support
argument then excludes a decreasing sequence of globally active exponents by
choosing a point where countably many nonzero holomorphic coefficient functions
are simultaneously nonzero. That step uses a Baire-category argument and an
appropriate classical choice principle to extract the decreasing sequence.

\emph{The step that is not another replacement argument.} The final extension
from ordinary tuples to arbitrary \emph{surcomplex} tuples is \emph{not} another
replacement argument. It proceeds one coordinate at a time using slice coherence
and coefficientwise identity, and its hypothesis must hold when all the other
fixed coordinates are surcomplex, not merely ordinary complex. Weakening this
hypothesis changes the theorem: arbitrary values at tuples with several
nonordinary coordinates would otherwise remain unconstrained.

A Lean specification should therefore expose a function on the relevant
universe-relative type, prove that its values on the small ordinary domain have
collectively small support, and then state the ordinary Hartogs and Baire lemmas
independently. \emph{Availability of those exact ordinary analytic lemmas in the
chosen dependency revision must be independently checked}; the existence of a
Hahn-series type does not provide them.

\subsection{Support certificates are the key reusable interfaces}
\label{found:sub:supportcertificates}
M1's summary of this section is the right implementation advice. The support
lemmas of \cref{found:sec:support} are ideal formalization targets: they are
independent of global surreal topology, and once proved they support many later
constructions. The proof pattern they enable is: replace the finitely many
retained coefficient germs by a finite ordinary analytic family; use the
classical identity for that family; compare finitely many coefficients; conclude
the Hahn coefficient identity. This is the formal counterpart of AG's
finite-parameter testing argument, and it does not require a single ordinary
contour for the original infinite family.

\section{Trigonometry, angles, and global phases}
\label{found:sec:trigonometry}

\emph{This entire section rests on a single audit and a single source. Only M2
was given the trigonometry manuscript T \cite{found:SrcTrig}, and no other audit
in this merge read it. Every statement below is tagged to that one reading.}

\subsection{The finite-angle coset group passes the size check}
\label{found:sub:cosets}
The finite-angle group $\OO_\R/(2\pi\Z)$ of T deserves a separate size check
(M2), and it passes for a reason that the raw-name equivalence classes of
\cref{found:sub:classquotient} do not: \emph{each coset $\theta+2\pi\Z$ is a
set, because it is indexed by $\Z$}. The totality of these cosets can therefore
be a class of sets. This is the clean contrast with the class-quotient trap:
there, zero had presentations of unbounded rank and its equivalence class was a
proper class; here every fibre is countable.

One may avoid the quotient altogether by using the unit-circle group, or by
choosing the unique representative in $[0,2\pi)$. Two warnings travel with the
latter choice. First, \emph{the interval $[0,2\pi)$ is itself a proper class in
the full surreal universe}: boundedness by ordinary reals does not imply
sethood, since even the interval of infinitesimals contains too many distinct
scales to be a set. Second, the chosen interval convention has an actual
\emph{branch jump}; it must not be mistaken for a globally analytic argument
function.

\subsection{Omnific integers, the finite-angle character, and its kernel}
\label{found:sub:omnific}
In the full surreal field the omnific integers have the normal-form description
\begin{equation}\label{found:eq:omnific}
 \Oz=\Pi\oplus\Z,
\end{equation}
where $\Pi$ consists of the purely infinite normal forms. For the finite-angle
$\cis$, put
\begin{equation}\label{found:eq:Uzero}
 U_0(p+u)=\cis(u),\qquad p\in\Pi,\quad u\in\OO_\R .
\end{equation}
Its kernel is $\Pi\oplus2\pi\Z=2\pi\Oz$, because multiplication by the nonzero
\emph{ordinary real} $2\pi$ preserves $\Pi$. Hence
\begin{equation}\label{found:eq:Ezero}
 E_0(x+iy)=\exp_{\No}(x)\,U_0(y)
\end{equation}
has kernel $2\pi i\Oz$.

\begin{proposition}[Why an infinite period is not a paradox]\label{found:prop:period}
Let $U:(\No,+)\to\mathbb T(\No)$ be a homomorphism extending the finite-angle
map, and assume the latter is already onto the full unit circle. Then every
coset of the finite additive subgroup contains an element of $\ker U$.
\end{proposition}
\begin{proof}
Given $x\in\No$, choose finite $a$ with $\cis(a)=U(x)^{-1}$. Then $U(x+a)=1$.
If $x$ is outside the finite subgroup, $x+a$ is still infinite.
\end{proof}

\subsection{The reconciliation with a canonical exponential}
\label{found:sub:canonicalexp}
T correctly separates finite-angle trigonometry from extensions determined only
by weak global group and local differential axioms, and proves an unavoidable
infinite period. Outside literature supplies the reconciling qualification (M2):
Ehrlich and Kaplan construct canonical surcomplex exponentials using an
\emph{initial integer part} \cite{found:EhrlichKaplanExp}. That is \emph{extra
structural information}, not a contradiction of nonuniqueness under weaker
hypotheses. Specifying the integer part kills the free purely infinite phase
component, and the kernel $2\pi i\Oz$ of \eqref{found:eq:Ezero} is the
normal-form description consistent with that canonical extension.

A formal library should expose both levels: finite-angle identities requiring no
infinite-phase choice, and the extra integer-part or character data used by a
selected global extension. \emph{Neither the foundation nor the bare field
axioms determine that data for free} (M2). The companion report
\emph{Trigonometry on the Surcomplex Plane}, held in this repository at
\texttt{docs/surcomplex/trigonometry/} \cite{found:CompanionTrig}, develops the
trigonometric theory itself; it is cited here by title and path only.

\subsection{The rational unit-circle parametrization is pure field algebra}
\label{found:sub:rationalcircle}
Much of T's geometric algebra --- determinant identities, the cosine-law
polynomial identity, Heron's factorization, circumcenter equations, Ptolemy's
inequality once a modulus is available --- is finite algebra over an ordered or
real-closed field. So is the rational parametrization of the unit circle:
\begin{equation}\label{found:eq:rationalcircle}
 t\longmapsto\frac{1-t^2}{1+t^2}+i\,\frac{2t}{1+t^2}.
\end{equation}
Representing an abstract direction by a circle point does not yet assert that a
chosen sine function parametrizes it by numerical finite angles; that is a later
theorem (M2). Much of the useful geometry is finite algebra even when the
lengths are infinite, and the computational complexity of evaluating a
particular surreal coefficient does not affect the logical validity of a
polynomial identity --- which makes generic algebra an unusually high-return
first formalization target.

\begin{remark}[Symbolic checks are not proofs]\label{found:rem:symbolic}
Validating finitely many Taylor coefficients cannot verify a support theorem, a
transfinite recursion, or an assertion about every scale. The symbolic checks
reported inside T are useful sanity checks of specific identities, \emph{not}
substitutes for these proofs (M2). The same limitation applies to the
illustrative Lean files in the M1 and M2 archives: they are intentionally much
narrower than the complete analysis project, and none of them was run
(\cref{found:sec:boundary}).
\end{remark}

\section{Axiomatizing the structure rather than constructing it}
\label{found:sec:axioms}

\subsection{Real closedness is an interface, not a characterization}
\label{found:sub:interface}
A type equipped with a real-closed-field structure supports much of finite
polynomial geometry after complexification. It does not thereby become the
canonical surreal field: $\R$, the real algebraic numbers and many
non-Archimedean fields satisfy the same first-order ordered-field sentences.
Adding an infinite element only rules out Archimedean examples; adding a
specified family of infinitesimals still does not recover Conway's simplicity
structure. The distinguished birthday function, simplicity tree, ordinal
embedding, normal-form monomials and small-cut filling need separate data and
laws. The same warning applies to an abstract algebraically closed field
advertised as ``surcomplex'': it may suffice for polynomial factorization, but
it does not supply a distinguished conjugation, an ordered fixed field, a Conway
monomial map, or an ordinary complex subfield (M2).

This suggests two legitimate axiomatic interfaces (M3). A \emph{generic algebra
interface} deliberately assumes only the field, order, valuation or closure
needed by a theorem. A \emph{surreal structure interface} also records the
simplicity and size data and has a proved model given by a construction. The
former is excellent for reusable mathematics; the latter is needed when one
wants to justify that the generic results apply to the intended surreal object.

\subsection{A safe abstract package}
\label{found:sub:package}
Merging M2's four ingredients with M3's interface clauses gives the following
specification. There is a carrier $S$ with:
\begin{enumerate}
\item an ordered-field structure, an ordinal-valued birthday map, a simplicity
relation with set-sized predecessor chains, and explicitly sized bounded
birthday strata;
\item a \emph{full sign-tree representation}: each element has a unique
ordinal-length sign sequence; every set-sized (or small) such sequence is
represented; numerical order is the termination-sensitive lexicographic order;
\item a cut operation for separated \emph{small} families, whose output lies
between its options, is characterized as the simplest eligible element, and
carries a rank bound; together with a compatibility theorem identifying the
designated simplicity relation with initial truncation in the sign
presentation;
\item arithmetic obeying the recursive Conway rules and respecting equality ---
not merely arbitrary field operations on the same carrier.
\end{enumerate}

This is a structural specification, \emph{not} a proposed minimal finite axiom
system (M2), and \emph{not} a claimed minimal or independent axiomatization
(M3). To establish categoricity one must state the precise notion of morphism
and prove an isomorphism to a canonical construction. Pure numerical order and
field operations are not enough, and adding an arbitrary rank label does not
force it to be the true birthday; conversely, specifying every cut by all
earlier numerical neighbours without a bound can make the options improperly
large. In a class theory the quantification over sign sequences is quantification
over set codes; in a universe-indexed theory it is restricted to the stated
universe. In a Lean project such an interface is best packaged as a structure
with fields and laws and then \emph{inhabited} by the constructed model; merely
declaring an uninhabited interface is harmless, but declaring a global inhabitant
as a new axiom shifts the entire consistency burden to that assumption (M3).

\subsection{Sign-tree categoricity}
\label{found:sub:signtree}
The following is M2's, and it is a \emph{foundational categoricity} statement.
It is renamed here from the source's ``rigidity'' to avoid collision with the
analytic all-scale polynomial rigidity of \cref{found:sub:rigidityargument},
which is a completely different theorem about a completely different subject;
see \cref{found:rem:rigidity}.

\begin{proposition}[Sign-tree categoricity]\label{found:prop:signtree}
Two realizations of the full sign-tree specification over the same ordinal and
set universe admit a unique sign-sequence-preserving isomorphism. If both use
the recursive Conway arithmetic, that isomorphism also preserves the arithmetic.
\end{proposition}
\begin{proof}
Map each element to the unique element of the second realization with the same
sign sequence. Existence and uniqueness of represented sequences make this a
bijection, and the order and simplicity relations are preserved by their
definitions. Compatibility with addition and multiplication follows by
well-founded induction on the input presentations, since the map identifies
corresponding option values and corresponding simplest cuts. No other
sign-preserving map can differ on any element.
\end{proof}

Ehrlich's work on simplicity hierarchies supplies more conceptual
axiomatizations and initial-embedding results
\cite{found:Ehrlich2001,found:EhrlichCorrection,found:EhrlichKaplanII}. Their
strength comes from the extra tree structure and the specified notion of initial
embedding, not from field axioms plus a vague maximality condition. The 2001
paper must be cited together with its published 2005 correction, which is listed
here as a separate bibliography item for exactly that reason (M1, M2).

\subsection{Universality caveats}
\label{found:sub:universality}
Universality claims need a specified class of inputs and embeddings (M3).
Embedding every \emph{set-sized} ordered field into the surreal class is
different from asserting a final object among all class-sized ordered fields.
Uniqueness of an initial simplicity-preserving embedding is different from
uniqueness of an arbitrary ordered-field embedding. The adjective ``absolute'' in
a structural account does not remove the need to say which parameters,
presentations and morphisms are preserved, and maximality in a structured
category must not be confused with the absence of every possible ordered-field
extension (M2). For the supplied manuscripts an all-purpose universality theorem
is not required: the concrete normal-form embedding of the selected Hahn
workspaces, plus the compatibility laws needed for each theorem, is both more
precise and more directly useful.

\subsection{A recent first-order program, at slide level, with its status boundary}
\label{found:sub:announcement}
All three audits cite the March 2026 announcement by Joel David Hamkins,
reporting joint work with Junhong Chen and Ruizhi Yang, of a first-order theory
of surreal arithmetic enriched by a birthday-precedence relation, with a reported
bi-interpretability with set theory \cite{found:Hamkins2026}. M1 reports the
lecture announcement posted February 4 2026 for a March 13 2026 CUNY Logic
Workshop talk; M2 reports the same and adds that it is not used as a proved
dependency. \textbf{Only M3 inspected the slides}, and its slide-level report is
therefore printed here as the substantive content, since it is the only part
actually read rather than inferred from an abstract
\cite{found:Hamkins2026Slides}.

From the slides (M3): the proposed theory uses \emph{schemes controlling
birthday-bounded definable cuts and definable families}. These restrictions are
the important foundational feature --- the proposed theory deliberately does
\emph{not} fill unrestricted definable cuts such as the entire carrier against
the empty side. That is precisely the restriction \cref{found:prop:allcuts}
forces on any consistent axiomatization.

The union of all three status boundaries applies, and every item survives:
\begin{itemize}
\item the slides explicitly describe the work as new, with details still being
checked; it is a promising research program, not a completed published
axiomatization;
\item it is not a verified consistency proof and not a substitute for a proved
model of the interface of \cref{found:sub:package};
\item it is \emph{not used as a proved dependency} of this report;
\item it is not a statement about the pure ordered-field theory or the pure
language of real-closed ordered fields; an expanded first-order structure can
encode information that the real-closed-field reduct cannot, and statements
about their theories must not be conflated;
\item the birthday-precedence relation of that program must be distinguished
from sign-prefix simplicity (\cref{found:rem:simplicity}) --- the one technical
warning M1 and M2 do not carry.
\end{itemize}
Its importance here is conceptual: enriching the field by construction history
can drastically change what its first-order language expresses.
\section{Dependent type theory and the universe boundary}
\label{found:sec:types}

\subsection{The universe shift is the construction, not an inconvenience}
\label{found:sub:universeshift}
In dependent type theory one constructs a family of surreal types relative to a
universe of small indexing types. Schematically,
\[
 \No_u:\Type_{u+1},\qquad I,J:\Type_u,
\]
with cut options $I\to\No_u$ and $J\to\No_u$: the larger universe contains the
carrier, while the smaller universe controls permissible branching. The illegal
cut $(\No_u,\varnothing)$ is excluded because its left index would generally
require the larger level. If all possible option indexing types were drawn from
the same size as the resulting carrier, one could try to use the whole carrier
as the left-option family and obtain \cref{found:prop:allcuts} again. The
constructor's arity and the carrier size must therefore be kept distinct. This
is the type-theoretic analogue of \cref{found:prop:universe}, and it is
explicitly \emph{not} a claim that types are literally sets in a particular
chosen model, nor an exact lower-bound theorem for the consistency strength of
any type theory or of Lean itself (M3).

An illustrative raw-game constructor has the shape
\begin{samepage}
\begin{lstlisting}
universe u

inductive RawGame : Type (u + 1) where
  | cut (L R : Type u)
      (left : L -> RawGame)
      (right : R -> RawGame) : RawGame
\end{lstlisting}
\end{samepage}
and all three audits print some form of it. The recursive occurrences are
strictly positive: they occur as function \emph{codomains}. Each branch is
structurally well founded even though the branching can be infinite;
``inductive'' does not mean ``finite birthday,'' and well-founded trees of
transfinite rank are represented by such inductive or W-type constructions (M2).
This is an illustration of the size discipline, not the current repository's
exact implementation (M1, M3).

\begin{remark}[The \texttt{Set RawGame} positivity and size warning]\label{found:rem:setrawgame}
M3 alone records the following. A tempting alternative constructor taking
\texttt{Set RawGame} on each side changes both positivity and size. A predicate
has the form \texttt{RawGame -> Prop}, so the recursive type occurs in a
function \emph{domain}: it is not the small positive family representation
above, and positivity is broken. Moreover a predicate on the entire larger
carrier need not have lower-universe-small support, so the size discipline is
broken as well. \emph{A constructor is not validated merely because its surface
notation resembles a mathematical set of options.}
\end{remark}

One should not try to put a not-yet-defined numerical comparison into the
constructor's hypotheses and expect an ordinary inductive declaration to solve
simultaneous order, quotient and numericity problems (M1).

\subsection{An ordinary-inductive route avoids primitive induction-recursion}
\label{found:sub:fivestep}
M2 gives the five-step route that removes the circular appearance of the cut
presentation:
\begin{enumerate}
\item define all raw games with small branching by an ordinary inductive
construction;
\item define the game comparison relation by well-founded recursion on pairs of
game trees;
\item define the numeric predicate recursively, requiring numeric options and
separation;
\item prove totality of comparison on numeric games, and prove the arithmetic
congruences;
\item quotient numeric games by mutual comparison.
\end{enumerate}
The comparison relation on arbitrary games need not be a total order; totality is
a theorem restricted to the numeric subtype. This restriction is indispensable:
games and surreal numbers are related objects, not synonymous types. The
alternative is a direct sign-sequence type indexed by a universe of ordinals,
which gives a canonical carrier and a straightforward simplicity relation but
still requires substantial recursion for the field operations. Neither approach
eliminates the mathematical work; it changes where that work is concentrated.
The Mizar development is evidence that no primitive induction-recursion
constructor is required (M2, M3) \cite{found:PakKaliszyk}.

\subsection{Quotient descent is a proof obligation}
\label{found:sub:quotient}
A function on raw games descends to numeric equivalence only after congruence is
proved; for addition this means
\[
 x\sim x',\ y\sim y'\ \Longrightarrow\ x+y\sim x'+y'.
\]
For a map selecting representatives, choice can supply an output, but the
mathematical result must not depend on which representative was selected. A
quotient recursor is the proper boundary: its respectfulness proof records
exactly this invariance.

Two different notions of equality are in play. Structural equality compares raw
presentations; quotient equality compares numerical values. A proof about the
options of a selected raw representative is not automatically a theorem about
arbitrary canonical options of the number, and a statement about the canonical
birthday of a number must not accidentally use the birthday of an arbitrarily
chosen noncanonical representative. The distinction matters especially for
birthday minimization and simplicity.

In a universe-relative quotient the collection of raw games is an actual type,
so the quotient constructor is a legitimate type-theoretic operation; this
avoids the literal proper-class-quotient problem of \cref{found:prop:scott}, but
it does not make all option families small and does not remove the need to prove
congruence (M3). Quotienting a large type usually leaves it large (M1).
Classical Lean has native quotients and a large quotient-based algebraic
library, so quotienting numeric games is a natural fit; the cost is that
reductions involving a representative chosen by choice are generally not
computational algorithms for surreal arithmetic (M2).

\subsection{Universe polymorphism is a schema, not a universal carrier}
\label{found:sub:polymorphism}
A declaration parameterized by universe $u$ can be instantiated at many levels.
It does not define a term $\sum_{u}\No_u$ over ``all universes,'' because
universe levels are not ordinary term-level elements of a single universal
universe. There is no top universe containing every universe at its own level,
and any purported final surreal type obtained by collapsing all levels must face
\cref{found:prop:allcuts} again. Lean has $\Type_u:\Type_{u+1}$ and is not
cumulative in the sense of automatic inclusion: explicit lifting constructions
such as \texttt{ULift} and universe-polymorphic declarations handle changes of
level \cite{found:LeanUniverses}. When comparing $\No_u$ with $\No_v$ one should
construct the embedding by lifting or re-encoding the underlying data, not assume
a definitional equality of carriers (M3).

\begin{remark}[\texttt{Small} is evidence to be proved]\label{found:rem:small}
Lifting a small type to a larger universe is harmless. \emph{Shrinking is
different.} A \texttt{Small} hypothesis is evidence that a type has an equivalent
representative in a specified smaller universe; Mathlib's mechanism records
exactly that, and a shrink construction uses such a witness. The proof may be
nonconstructive, but it is still a theorem to be supplied. It cannot be
manufactured to convert $\No_u$ into a legal index type for all of its own
options: the whole $\No_u$ cannot be declared $u$-small consistently with its
small-cut filling property and its linear order. A \texttt{Small} hypothesis on a
support should be an explicit constructor argument or theorem assumption; it
should not be inferred from the fact that the values happen to be surreal (M1,
M3).
\end{remark}

For a fixed-level surreal type the admissible cut interface should therefore
resemble
\[
 L,R:\Type_u,\quad \ell:L\to\No_u,\quad r:R\to\No_u,
 \quad\forall a,b\ \ \ell(a)<r(b),
\]
with $\No_u:\Type_{u+1}$. An interface using \emph{subsets} of $\No_u$ must
explicitly require their supports to be $u$-small. The same rule applies to Hahn
supports. \Cref{found:sub:cutdata} makes this concrete.

\subsection{Impredicative \texorpdfstring{$\Prop$}{Prop} is not \texorpdfstring{$\Type:\Type$}{Type : Type}}
\label{found:sub:prop}
Lean's $\Prop$ is impredicative: a proposition may quantify over arbitrary types
without moving to a larger $\Prop$, its proofs are definitionally
proof-irrelevant, and elimination from propositions into data is restricted.
These features are not an unrestricted universe of data containing itself
\cite{found:LeanPropositions,found:LeanUniverses}. The statement ``every surreal
in this universe has property $P$'' is well formed even when its quantifier
ranges over a large type; this does not produce a smaller type enumerating those
surreals. Logical quantification and construction of a data collection are
separate operations.

A predicate $P:\No_u\to\Prop$ defines a subtype at the existing data level; it
does not imply that the subtype is $u$-small (M2). This rules out the tempting
error of using the small universe of truth values to encode all surreals as
small data: the \emph{output} of a predicate is a proposition, but the domain,
graph, subtype and family of all predicates retain their appropriate universe
levels. Propositional resizing principles, in foundations that adopt them,
address proposition universes; they do not compress the ordinal-indexed surreal
carrier or its arbitrary supports, and resizing must be a stated principle with
a suitable scope rather than a consequence of impredicativity.

\subsection{Classical and computational content}
\label{found:sub:noncomputable}
A \texttt{noncomputable} definition in Lean is \emph{not} an unproved axiom
(M1). It can be fully kernel-checked while depending on classical choice.
Arbitrary real coefficients, arbitrary well-ordered supports and quotient
representatives naturally lead to such definitions. Efficient computation
requires additional presentations --- finite games, effective exponent
arithmetic, an enumeration procedure for coefficient dependencies --- and the
abstract existence of a finite contributing set is not automatically an
algorithm for finding it. Arbitrary well-ordered supports need not have a
computable enumeration or decidable coefficient equality; exact finite-prefix
computation requires a representation and a theorem that the prefix depends only
on the supplied finite data (M3). A formalization can first verify the
noncomputable mathematics and later refine effective subclasses.

Choice at a fixed universe level can select representatives from the types under
discussion. It does not form a single function quantified internally over all
universe levels, and identifying such an axiom with set-theoretic global choice
requires a specified interpretation of sets and classes: the common word
``choice'' alone is not a proof that the foundational principles are equivalent
(M3). ``Lean choice'' and ``NBG global choice'' are not literally the same axiom
in the same language (M2).

\subsection{Internal \texttt{ZFSet} as an alternative route}
\label{found:sub:zfset}
One can also formalize a set-theoretic universe inside Lean and define surreals
as sets in that model. Mathlib provides pre-sets \texttt{PSet} in
$\Type_{u+1}$, quotients them by extensional equivalence to obtain
\texttt{ZFSet}, and represents internal classes by \texttt{Set ZFSet}
\cite{found:MathlibPSet,found:MathlibZFBasic,found:MathlibClass}. The
distinction
\[
 \texttt{Set (ZFSet.\{u\})}\qquad\text{versus}\qquad \texttt{ZFSet.\{u\}}
\]
separates an external Lean predicate class from an internal set, and Mathlib
documents it explicitly. This route permits translating a conventional
set-theoretic argument almost literally, and it is well suited to interpretation
or conservativity results --- internal properness of $\No$, interpretations of
class arithmetic, or the dependence of a construction on a collection principle.

The cost is a bridge to ordinary algebraic types. An internal field described by
set-coded operations is not definitionally a Lean \texttt{Field} instance on its
external element type: one must extract the carrier, transport the operations,
and establish the relation between internal membership and external predicates.
Each ordinary field, polynomial, function and set operation may need an
internal-to-native translation. The existence of a ZFC model in Lean also does
not mean that Lean is ZFC under another spelling; the construction uses Lean's
ambient type theory and universe hierarchy.

\begin{remark}[The \texttt{Definable} warning]\label{found:rem:definable}
Mathlib's technical predicate \path{ZFSet.Definable}, and its theorem
\path{Classical.allZFSetDefinable}, must \emph{not} be read as saying that
every external function is definable by a first-order formula in the language of
set theory. The API concerns representative-level liftability of functions
through pre-set representatives; the documentation expressly notes the
distinction. One must not use that identifier as evidence that a native Lean
function has been translated into a ZFC formula. First-order definability for a
metamathematical interpretation must be formalized with actual syntax and
satisfaction, which is a further layer. M1 and M3 record the warning for
\path{ZFSet.Definable}; \textbf{M2 alone also names
\path{Classical.allZFSetDefinable}}.
\end{remark}

To compare GBC and KM inside Lean one must specify a two-sorted language,
formulas, models or interpretations, and the exact comprehension and replacement
schemas (M1). If classes are interpreted by all external Lean predicates on an
internal universe, one must prove which schemas are satisfied and at what
metatheoretic strength; the fact that the predicate type exists is not a proof
that the intended minimal class theory and its exact consistency strength have
been reproduced. Unrestricted external quantification over predicates can
provide more semantic resources than elementary class comprehension, so a
development claiming conservativity or precise axiom calibration must reason
about formal proof systems or explicitly restricted model expansions.

\section{Constructive set theory and univalent foundations}
\label{found:sec:constructive}

\subsection{IZF, CZF, and type-theoretic interpretations}
\label{found:sub:izf}
IZF retains a power-set style set theory with intuitionistic logic; CZF is
designed for predicative constructive mathematics and replaces some classical
set-existence principles with constructive collection principles. Type-theoretic
interpretations of constructive set theories provide a well-developed connection
with dependent types and W-type methods \cite{found:GambinoAczel}, and Aczel's
type-theoretic interpretation with its treatment of inductive definitions is the
standard source \cite{found:Aczel}; regular extension axioms supply additional
closure resources in some treatments \cite{found:Rathjen}. In weak theories
without power set, Replacement and Collection can differ in ways that are
immaterial in ordinary ZFC but material to formalization
\cite{found:ZFCminus}.

These are genuine foundational options, but a classical sign-sequence argument
cannot be moved into them by deleting the words ``by excluded middle'' (M2). One
cannot obtain a constructive formalization simply by removing the word
``classical'' from a classical surreal proof (M3).

\subsection{Why the constructive order needs positive information}
\label{found:sub:apartness}
Classically, strict comparison can be recovered from a suitable negated
non-strict comparison. Constructively, absence of a proof of $y\le x$ need not
give positive evidence that $x<y$, and a definition based on that negation can
also violate the positivity requirements of the intended inductive
construction. The equivalence
\[
 x<y\quad\Longleftrightarrow\quad\neg(y\le x)
\]
must not be assumed without the relevant order principle. An \emph{apartness}
relation --- $x\mathrel{\#}y$ meaning $x<y$ or $y<x$ --- carries positive
information absent from mere inequality. Even familiar definitions of a field
differ constructively: inversion from positive apartness from zero is stronger
evidence than mere negation of equality, and ordinary Cauchy and Dedekind reals
need not coincide without additional principles. Classical ordinal
comparability, deciding whether an option lies above a point, and turning
negated universal statements into witnesses can all fail constructively.

\subsection{The higher inductive-inductive construction}
\label{found:sub:hiit}
The HoTT book develops surreal numbers by a higher inductive-inductive
definition \cite[\S11.6]{found:HoTTBook}: one simultaneously defines the
carrier, its non-strict order and its strict order. Small separated left and
right families create a number, and an antisymmetry path constructor identifies
$x$ and $y$ when both inequalities hold. The positive comparison pattern reads
\begin{align*}
 [\forall x^L\;(x^L<y)]\ \land\ [\forall y^R\;(x<y^R)]&\Longrightarrow x\le y,\\
 [\exists x^R\;(x^R\le y)]\ \lor\ [\exists y^L\;(x\le y^L)]&\Longrightarrow x<y,
\end{align*}
replacing the classical negative test \eqref{found:eq:comparison}. Positive
constructors for the order relations support recursion without requiring an
arbitrary lift of quotient-valued option families to raw representatives.

This is not the classical quotient construction with different syntax. In a
constructive quotient approach a family of equivalence classes need not have a
chosen family of representatives, and the higher inductive-inductive
construction addresses that in its elimination principle. The numerical and
order components depend on each other, so a standard W-type alone does not
present the whole construction. A formalization must specify the elimination
principle, the coherence of the resulting operations, and the status of the
order relations as propositions. It \emph{is not licensed merely by writing
mutually recursive symbols} (M1), and there is no blanket theorem that every
informal recursive signature is sound: general semantic work on higher inductive
types is important background but does not certify a particular signature
\cite{found:LumsdaineShulman,found:CubicalHIT}. A proof assistant supporting
some higher inductive types does not thereby support every higher
inductive-inductive signature with all its required computation rules.

\subsection{Plump ordinals and constructive comparison}
\label{found:sub:plump}
Constructively, classical well-orders, trichotomous ordinal notions, and the
ordinal objects naturally living inside surreals need not coincide without
further hypotheses. Shulman's discussion of \emph{plump ordinals} illustrates
the distinction and explains why the constructive definition treats strict and
non-strict comparison separately \cite{found:ShulmanPlump}. A constructive
comparison theorem must therefore name its ordinal notion; the slogan ``the
surreals contain all ordinals'' needs more interpretation than in classical ZFC.
Joseph develops an ordered commutative-ring structure with apartness in univalent
foundations \cite{found:Joseph}. These sources establish that the constructive
approach is mathematically substantial; they do \emph{not} by themselves
constitute a Lean 4 implementation of the complete surcomplex analytic theory
(M2), and the constructive route is not claimed to already supply the
normal-form, algebraic-closure and analytic package of the supplied manuscripts
(M3).

\subsection{Why full univalence cannot simply be added to Lean's equality}
\label{found:sub:univalence}
Lean's ordinary equality is proposition-valued and proof-irrelevant, so for
fixed types $A,B$ the identity type $A=B$ has at most one proof. Full univalence
for a universe with nontrivial automorphisms has a different homotopical
equality structure, in which distinct self-equivalences correspond to distinct
equality paths. M2 and M3 both prove the incompatibility as a numbered
proposition by the same two-element-type argument; M1 states it in prose only.
We print the proof once.

\begin{proposition}[The elementary incompatibility test]\label{found:prop:univalence}
Suppose a type theory has proof-irrelevant equality between types, contains a
two-element type $B$, and asserts full univalence for a universe containing $B$.
These assumptions are incompatible.
\end{proposition}
\begin{proof}
Proof irrelevance makes the type of equality proofs $B=B$ have at most one
element. Univalence makes its map to the type of self-equivalences
$B\simeq B$ an equivalence, so $B\simeq B$ would also be a subsingleton and all
self-equivalences of $B$ would be equal. But the identity map and the map
exchanging the two elements are not equal, as evaluation at either element of
$B$ shows.
\end{proof}

The conclusion is drawn carefully, and in exactly this form: \emph{it is not
that univalent mathematics is inconsistent}. It is that incompatible equality
mechanisms cannot be superimposed on Lean's native proof-irrelevant type
equality without changing the foundation. Propositional extensionality is a
different, much weaker and compatible assertion \cite{found:LeanAxioms}.

M1 names the three legitimate alternatives, and they survive: translate the
set-level result to quotient-based mathematics; encode the syntax and semantics
of an HoTT object theory \emph{inside} Lean; or conduct the higher-inductive
development in a system designed for it. None of these changes Lean's kernel
equality. For the classical algebra and analysis of the supplied manuscripts,
ordinary quotient-based Lean mathematics is a better fit than rebuilding the
analysis on a new homotopical notion of equality.

\subsection{Categorical and algebraic-set-theoretic perspectives}
\label{found:sub:categorical}
An elementary topos with a natural numbers object supplies a useful environment
for much ordinary mathematics, but does not by itself guarantee the unrestricted
cumulative hierarchy or the class of all classical ordinals. Structural set
theories can be strengthened with replacement-like principles, universes, or
classes of small maps; in such a framework a surreal construction must specify
which well-founded trees are small and why the relevant recursive constructions
preserve smallness \cite[\S\S13--15]{found:ShulmanSize}. Algebraic approaches to
surreal structures have been developed in a class-set-theoretic setting
\cite{found:RangelMariano}.

Such frameworks do \emph{not} remove the diagonal obstruction (M1, M2).
Internally in any logic where $<$ is irreflexive, a rule producing a strict upper
bound for every subobject of the same carrier fails on its maximal subobject. A
categorical account must restrict its small maps or its admissible cuts, and must
state its semantics, quotient existence and inductive constructions explicitly
rather than hiding them behind the word ``universal.'' Likewise a categorical
description of $\No$ as a class-directed union of substructures is useful, but
one must not use an ordinary small-diagram colimit constructor on a proper-class
index category (\cref{found:rem:colimit}).

\subsection{Choosing the constructive route for the right purpose}
\label{found:sub:constructivepurpose}
A constructive development is attractive when the computational meaning of
order, choice-free construction, or the elimination principle is central; a
classical quotient-based Lean development is attractive when the immediate goal
is to reuse algebra, ordinary complex analysis and established classical
libraries. For a constructive version of the supplied analytic theory the axiom
ledger would have to revisit the extraction of decreasing sequences, Baire
arguments, holomorphic zero-set reasoning, and the choice of roots or
representatives. The aim need not be to recover the classical statements
unchanged: one can prove a constructive theorem with explicit apartness,
locatedness or supplied enumeration hypotheses, and then recover a classical
specialization when excluded middle and suitable choice are added. A constructive
project should specify its ordinal notion, its inductive or W-type existence
principles, its quotient policy, and its use of choice or propositional resizing.
The two routes can inform each other through explicit interpretation theorems;
neither is a drop-in implementation of the other.

\section{The current Lean ecosystem: one inspection, three readers}
\label{found:sec:ecosystem}

\subsection{The aggregation warning, in plain language}
\label{found:sub:aggregation}
Everything in this section comes from \textbf{one act of source reading,
performed three times}. The three audits inspected the same repository at the
same immutable revision on the same day, through the public web interface,
reading source only. They agree. That agreement is evidence about what the
source \emph{says}, and about nothing else.

It is not evidence that the repository builds. It is not evidence that its
theorems mean what their names suggest. It is not evidence about its axiom
closure. Three concurring readings of one web page corroborate each other only
about the contents of that page. No claim in this section was kernel-checked,
and the findings below are printed once, not three times, precisely so that
their repetition cannot be mistaken for replication. See
\cref{found:sec:boundary}.

\subsection{The pinned revision}
\label{found:sub:pinned}
\begin{center}\small
\begin{tabular}{@{}P{0.28\textwidth}P{0.62\textwidth}@{}}
\toprule
Item & Recorded value\\\midrule
Repository & \texttt{vihdzp/combinatorial-games} (a separate Lean 4 project,
\emph{not} upstream Mathlib)\\
Revision inspected & \texttt{02b4a908ea2ecfefffecb438f691951a814a5264}\\
Commit timestamp & September 19, 2026, 10:25:40 UTC\\
Inspection date & September 21, 2026\\
Lean toolchain & \texttt{leanprover/lean4:v4.35.0-rc2}\\
Mathlib revision & \texttt{dec5b2b780537b6eaf7f5e5f000c12f7387fb24d}\\
Method & public GitHub interface, source reading only; no clone, no build, no
comprehensive placeholder search, no axiom audit\\
\bottomrule
\end{tabular}
\end{center}

All three audits record the repository revision and the toolchain; all three
record the Mathlib revision, each reading it from the \path{lake-manifest.json}
at that pinned commit rather than from a search-engine snippet or an unrelated
cached branch page
\cite{found:GamesRepo,found:GamesToolchain,found:GamesManifest}. Files inspected
at that revision, with pinned blob URLs in the bibliography:
\path{CombinatorialGames.lean} (umbrella module inventory, M3);
\path{CombinatorialGames/Surreal/Basic.lean};
\path{CombinatorialGames/Surreal/Division.lean};
\path{CombinatorialGames/Surreal/HahnSeries/Basic.lean};
\path{CombinatorialGames/Surreal/Cut.lean} (M3 only); plus
\path{lean-toolchain} and \path{lake-manifest.json}.

\begin{remark}[Historical caution]\label{found:rem:lean3}
Old Lean 3 surreal documentation describes an earlier stage of the development
and is not a reliable inventory of the current separate Lean 4 project.
Conversely, finding a theorem in that project does not make it available by
importing upstream Mathlib alone (M3).
\end{remark}

\subsection{The surreal carrier and the field instance}
\label{found:sub:carrier}
At the pinned revision, \path{CombinatorialGames.Surreal.Basic} defines the
carrier in universe $u+1$ by antisymmetrizing the subtype of numeric games:
\begin{samepage}
\begin{lstlisting}
def Surreal : Type (u + 1) :=
  Antisymmetrization (Subtype Numeric) (...)
\end{lstlisting}
\end{samepage}
where the ellipsis stands schematically for the comparison relation. The same
module supplies the quotient map, numerical equality of representatives, a
linear order, additive structure, a choice-based representative operation
(\texttt{Surreal.out}, using quotient choice), and a set-based cut interface
carrying smallness conditions on the option families. This is the
universe-sensitive pattern developed in \cref{found:sec:types}, not a single
small type accepting arbitrary subsets of itself
\cite{found:GamesBasic}.

The inspected \path{CombinatorialGames.Surreal.Division} contains an
\emph{actual} declaration
\begin{samepage}
\begin{lstlisting}
noncomputable instance : Field Surreal where
  ...
\end{lstlisting}
\end{samepage}
whose inverse descends from numeric-game inversion by a quotient map and a
congruence theorem, with the field cancellation proof transported through
numerical equivalence and inverse correctness established by well-founded
induction over numeric games \cite{found:GamesDivision}. It is therefore
incorrect to describe surreal multiplication or inversion as wholly
unimplemented in the current inspected ecosystem, or to describe current Lean
surreal work solely by older Lean 3 documentation.

\begin{remark}[The ``complete ordered field'' prose]\label{found:rem:completeprose}
Prose in the basic module uses the phrase ``complete ordered field.'' It must
\emph{not} be read as a proved Dedekind-completeness statement;
\cref{found:prop:incomplete} gives the elementary obstruction. The field
structure and the actual theorem signatures are what matter. This is a useful
example of why a source audit cannot stop at a module's introductory paragraph.
\end{remark}

\subsection{The small-support surreal Hahn layer}
\label{found:sub:hahnlayer}
The inspected \path{CombinatorialGames.Surreal.HahnSeries.Basic} defines
\texttt{SurrealHahnSeries} in universe $u+1$ as a small-support subfield of
lexicographically ordered real Hahn series on the \emph{order dual} of the
surreal exponent type. Its \texttt{mk} constructor requires \emph{both} a
\texttt{Small.\{u\}} support witness \emph{and} reverse well-foundedness of the
support. It provides field, linear-order and strict-ordered-ring infrastructure,
coefficient extraction, support, truncation, and ordinal indexing of the small
support \cite{found:GamesHahn}.

Two details matter. First, it lives at $\Type_{u+1}$, matching the surreal
carrier. Second, support is required to be $u$-small although the exponent
carrier is larger. This is precisely the relative-smallness discipline of
\cref{found:sub:toolarge}, and it is an important \emph{implemented} safeguard:
the library does not identify the unrestricted Hahn-series type with the surreal
carrier.

The module presents itself as a translation layer. In the inspected file,
including its end, \emph{no completed ordered-field isomorphism between the
surreal carrier and this Hahn type was verified}. A constructor for the target
type, or a comment describing an intended future isomorphism, is not that
isomorphism: normal-form existence, surjectivity of evaluation and compatibility
with multiplication remain proof obligations unless a separately located theorem
discharges them. A finished bridge needs maps both ways, inverse laws,
compatibility with arithmetic, and support preservation.

Likewise, \emph{no real-closedness instance for the concrete game-based carrier,
no surcomplex algebraic-closure specialization, and no Gonshor exponential was
verified in the inspected modules}. These must be checked or supplied before
specializing the algebraically closed-field parts of the supplied manuscripts.

\begin{remark}[Scope of the negative findings]\label{found:rem:scopenegative}
Every negative finding in this subsection is \emph{scope-limited by
construction}: it is a statement about what these particular files, read on this
particular date, were observed to contain. It is \textbf{not} a claim that these
results are absent from any other file, branch, repository, project or later
revision. All three audits state this qualification, and it applies to every
sentence of the form ``was not verified'' in this section.
\end{remark}

\subsection{What \texttt{Surreal/Cut.lean} actually contains}
\label{found:sub:cutfile}
\textbf{M3 alone inspected this file}, and it is the module most likely to be
misread. The inspected \path{CombinatorialGames.Surreal.Cut} concerns
complementary lower and upper \emph{sections of the entire universe-relative
surreal order}. It builds a complete linear order of such sections via
concept-lattice machinery, and a number corresponds to two adjacent sections
according to which side contains it \cite{found:GamesCut}.

This is \emph{not} the same object as a small pair of Conway options. A section
of the full carrier may be large and need not be filled by a number in that
carrier. Therefore its completeness yields \emph{neither} an unrestricted Conway
cut constructor \emph{nor} Dedekind completeness of the surreal field, and it
does not contradict \cref{found:prop:allcuts} or
\cref{found:prop:incomplete}. M3's own instruction stands: \emph{a formal API
should name the option-pair interface and the section interface distinctly, even
though the literature calls both ``cuts.''} \Cref{found:sub:cutdata} adopts that
instruction.

\subsection{What generic Mathlib supplies}
\label{found:sub:mathlibhahn}
Mathlib's Hahn-series development already provides coefficient functions with
support restrictions, convolution operations, and formal summable families. In
the inspected documentation, \path{HahnSeries.SummableFamily} carries
\emph{both} a partially well-ordered union-of-supports condition \emph{and}
finite coefficient co-supports; the order hypotheses specialize that support
condition to the usual well-ordered condition used here, and its \texttt{hsum}
operation is the \emph{formal} sum, not an unqualified topological series
\cite{found:MathlibHahnBasic,found:MathlibHahnSummable}. Mathlib's documentation
explicitly distinguishes this from topological summability, so inventing a new
use of ordinary \texttt{tsum} would obscure an already-supported distinction
(M2). A surreal specialization must prove that the general partially-well-ordered
hypotheses yield the required well-ordered supports in the chosen linear order,
and must establish smallness of the resulting support, not only its
order-theoretic shape (M2, M3).

\textbf{M1 alone} inspected the \texttt{HEval} module. It supplies one-variable
\path{PowerSeries.heval}, evaluation of a formal power series at a
positive-order Hahn series, and its retrieved documentation explicitly lists the
finite-variable \path{MvPowerSeries.heval} as a \textbf{TODO}
\cite{found:MathlibHEval}. This is the single concrete extension target named
anywhere in the group. The observation concerns the inspected module and its
documentation; it is \emph{not} a proof that no other project contains a
multivariable evaluation theorem.

\textbf{M2 alone} inspected the real-closed-field page. Mathlib has an
\texttt{IsRealClosed} \emph{interface} \cite{found:MathlibRealClosed}. An
interface name is not an instance for a particular type. Before depending on real
closedness of the actual game-based \texttt{Surreal}, a project must locate or
prove the relevant instance and check its dependencies; similar care is needed
when deriving algebraic closedness of its quadratic extension.

Mathlib also supplies \texttt{PSet}, \texttt{ZFSet} and class infrastructure
(\cref{found:sub:zfset}).

\begin{remark}[Rolling documentation]\label{found:rem:rolling}
All Mathlib generated documentation pages and all official Lean ``latest''
reference pages cited in this report are \emph{moving} sources consulted on
September 21, 2026. They are \textbf{not} asserted to be generated from the
pinned dependency commit, and they are not substitutes for the repository's
dependency lock. Exact module imports and instance hypotheses must be re-tested
against the pinned Mathlib revision, and exact typeclass names should be taken
from the pinned dependency source rather than invented to match a mathematical
phrase. All three audits record this.
\end{remark}

\subsection{Mizar as an independent completed milestone}
\label{found:sub:mizar}
P\k{a}k and Kaliszyk's ITP 2024 paper describes a substantial Mizar formalization
combining complementary presentations of surreal numbers and reaching Conway
normal form, discussing how ordinary transfinite methods replace an otherwise
troublesome inductive-recursive presentation and how sign-based and Conway-style
reasoning support different parts of the development
\cite{found:PakKaliszyk}. It demonstrates that the representation bridge is a
realistic, substantial formalization task, not an optional piece of notation.

Mizar's Tarski--Grothendieck foundation and its set-theoretic proof style differ
from Lean's native dependent type theory. The development is therefore a
valuable mathematical and formal-engineering reference, \emph{not} a source of
Lean proof terms that can be imported unchanged: theorems do not transfer between
libraries merely because their informal names agree. No Mizar run was performed
here; no formal article in that development was inspected individually; its axiom
closure was not reconstructed.

M2 and M3 both flag one further terminological point: the abstract's description
of the surreal field as ``algebraically closed'' must be distinguished from the
actual real-versus-complex field statements. The ordered surreal field cannot
have a root of $X^2+1$; algebraic closedness belongs to $\No[i]$
(\cref{found:sub:realclosed}).

\subsection{Terminological warnings, collected}
\label{found:sub:terminology}
Four naming traps recur, and all four are recorded here in one place.
\begin{enumerate}
\item ``\textbf{complete ordered field}'' in source prose is not a
Dedekind-completeness theorem (\cref{found:rem:completeprose}).
\item ``\textbf{cut}'' names at least two different implemented objects: a small
Conway option pair, and a complementary section of the whole order
(\cref{found:sub:cutfile}).
\item \texttt{\textbf{IsRealClosed}} is an interface, not an instance
(\cref{found:sub:mathlibhahn}).
\item \texttt{\textbf{ZFSet.Definable}} and
\texttt{\textbf{Classical.allZFSetDefinable}} are representative-level
liftability conditions, not first-order definability by a formula
(\cref{found:rem:definable}).
\end{enumerate}
To these add the Mizar abstract's ``algebraically closed'' wording
(\cref{found:sub:mizar}).

\subsection{The evidence table, with its third column}
\label{found:sub:evidencetable}
The following is M1's evidence table, reproduced with its third column
\textbf{verbatim}. That column --- ``What is not claimed'' --- is the group's
only artifact built specifically to prevent a reader from converting inspected
declarations into verified mathematics, and it is used here as the organizing
device of this whole section.

\Needspace{18\baselineskip}
\begin{longtable}{@{}P{0.27\textwidth}P{0.34\textwidth}P{0.33\textwidth}@{}}
\toprule
Component & Evidence inspected & What is not claimed\\
\midrule\endhead
Numeric-game quotient & Current \texttt{Surreal.Basic} definition & A universe-free type of all surreals\\[0.35em]
Surreal field & Explicit \texttt{Field Surreal} declaration & Independent local rebuild and complete dependency audit\\[0.35em]
Surreal Hahn layer & Small-support subfield, coefficients, truncations, ordinal indexing & Finished normal-form equivalence inferred from the module title\\[0.35em]
Formal Hahn sums & \texttt{SummableFamily} and \texttt{hsum} APIs & Topological convergence of every formal sum\\[0.35em]
Power-series evaluation & One-variable \path{PowerSeries.heval}; multivariable TODO in the inspected module & A proof that multivariable evaluation exists nowhere else\\[0.35em]
Internal set theory & \texttt{PSet}, \texttt{ZFSet}, and class infrastructure & An absolute consistency proof of ZFC or Lean\\[0.35em]
The supplied analysis & Mathematical proofs in the supplied manuscripts & Already machine-verified analytic geometry and residues\\
\bottomrule
\end{longtable}

\noindent Two rows must be read with \cref{found:sub:ledger} in hand: ``the
supplied analysis'' in M1's original table meant the two manuscripts M1 read,
and in this merged report the row covers five manuscripts, of which only one was
read by all three audits.

\subsection{Per-archive Lean inventory}
\label{found:sub:archiveinventory}
The per-archive inventory --- which files ship, in what state, with what import
footprint --- is \cref{found:sub:archives}, in the verification-boundary
section, where it belongs. Its single most important line bears repeating here:
M1 ships three uncompiled Lean files, one of which imports \texttt{Init} only;
M2 ships two uncompiled Lean files, one of which imports all of Mathlib; and M3
ships \emph{no Lean files at all}, its five code blocks existing only as
typeset listings inside its PDF. None of the five files was compiled by the audit that shipped it; three of
the five have since been compiled for this report (\cref{found:sub:compiled}).
\section{A layered Lean architecture}
\label{found:sec:architecture}

\emph{All module names, layers, architectures, completion criteria and
implementation sequences in this section are design suggestions. None of them is
a claim that files with those paths exist.}

\subsection{Separate the global number type from the local algebra}
\label{found:sub:twotracks}
M1's organizing decision is to run two tracks. The \emph{foundational track}
uses the existing numeric-game quotient, sign-expansion interfaces and
small-support normal forms to justify a universe-indexed surreal
interpretation. The \emph{analytic track} works over arbitrary set-sized Hahn
fields and ordinary complex germ rings. They meet at explicit embedding and
compatibility theorems.

This avoids making the entire analytic program depend on finishing every global
surreal theorem first. An abstract theorem for $\C((t^\Gamma))$ is already
genuine mathematics, and its eventual application to $\No[i]$ has a clearly
identified dependency: the normal-form embedding and its preservation of the
constructions used. Conversely the local track must not claim that an arbitrary
ordered abelian group $\Gamma$ is already a subgroup of $\No$; it can prove
theorems for abstract $\Gamma$, and the embedding is additional data at the
application boundary (\cref{found:rem:nocover}).

M3's complementary decomposition is by \emph{kind of definition}, in four
layers: \emph{representation} (raw games, numerical equivalence, canonical sign
data, birthdays, universe lifts); \emph{algebra} (ordered-field laws, real
closedness, complexification, finite polynomial and matrix structures);
\emph{Hahn structure} (monomials, small normal forms, support-preserving maps,
valuation, residue, strong sums); and \emph{analysis} (ordinary holomorphic
coefficient rings, coherent evaluation, finite-domain division, contours,
full-scale hypotheses). Dependencies must point in that direction: a proof of a
residue determinant should import finite algebra, not full surreal topology; a
proof that a support recursion stays in $S+E^*$ should import ordered-group
support theory, not Hartogs. The full-scale rigidity theorem may then combine
the small-cut bound from the first layer with the common-support fact from the
fourth.

\subsection{Layers A--E, with completion tests}
\label{found:sub:layers}
M2's lettered layering is the most operational of the three and is adopted here
as the working spine.

\paragraph{Layer A: reusable finite algebra.} Start with an arbitrary ordered
field $F$, not a concrete surreal datatype. Define its quadratic extension with
the correct cross-term multiplication; prove conjugation, norm-square
identities, inverses and the Gram identity \eqref{found:eq:gram}. Add real
closedness only where positive square roots or polynomial factorization are
genuinely required. This layer already carries most of the geometric algebra of
\cref{found:sec:trigonometry}, including the rational circle parametrization
\eqref{found:eq:rationalcircle}.

\paragraph{Layer B: set-sized Hahn analysis.} Use $F_\Gamma$ and $K_\Gamma$ as
ordinary types, with the ordered value group, its divisibility when needed, and
the coefficient field as explicit parameters. Reuse the existing
formal-summability API. Establish standard part for finite series from the
coefficient at exponent zero, and distinguish finite elements from
positive-valuation infinitesimals. Implement evaluation as an operation whose
input includes an admissibility proof, or as a total operation with an
accompanying theorem guaranteeing its meaning on admissible inputs; the former
is clearer:
\[
 \operatorname{evaluate}\bigl(A,\ h,\ \text{proof that }(a_nh^n)_n
 \text{ is strongly summable}\bigr).
\]
An ordinary polynomial needs no infinite-sum proof; a formal power series does.
Keeping these constructors distinct prevents accidental evaluation outside the
proven domain. The next targets are the positive-support geometric series, the
binomial series, the infinitesimal exponential and logarithm, and Taylor lifting
of ordinary analytic germs.

\begin{remark}[The Layer B recentering warning]\label{found:rem:recentering}
M2 alone records this, and it is easy to get wrong. At a finite point
$r+\epsilon$, first evaluate the ordinary function at $r$ and \emph{then}
substitute $\epsilon$ into its Taylor series. Do \emph{not} replace this by the
ordinary Maclaurin series at the nonzero real $r$: infinitely many of those
terms can occupy the same Hahn exponent, and the second clause of
\eqref{found:eq:summability} fails.
\end{remark}

\paragraph{Layer C: support-controlled coherent functions.} For an ordinary open
$U\subseteq\C$, represent a coherent datum as a Hahn series with coefficients in
the ring $\OO(U)$ of ordinary holomorphic sections. Equality of sections is
equality on $U$; values outside $U$ must not become part of the object.
Coefficientwise differentiation does not enlarge the common support, and
integration along a fixed ordinary admissible contour also acts coefficientwise
and does not enlarge it --- which is why ordinary complex-analytic theorems can
be applied to the coefficients and then reassembled. Arbitrary composition does
\emph{not} follow: for that one must control the supports of the inner
perturbations and verify that the ordinary part maps into the common coefficient
domain. A useful composition interface has an ordinary base map $g_0:V\to U$ and
an infinitesimal perturbation whose coefficient support lies in a common positive
well-ordered set $T$; Taylor expansion then produces supports in $S+T^*$, and
the formalization must prove local finiteness of coefficient contributions as
well as well-ordering.

\paragraph{Layer D: the universe-indexed surreal bridge.} Establish explicit maps
among numeric games, canonical sign sequences and small-support normal forms.
Deliverables: preservation of order, field operations, ordinary real
coefficients, Conway monomials and admissible sums, with the support universe
visible in the statement of the strongest comparison theorem. The reason a
small-support subtype is required is \cref{found:sub:toolarge}.

\paragraph{Layer E: topology, functions, and enrichment.} Define the fine
topology and the derivative \eqref{found:eq:derivative} using all radii of the
chosen carrier. State the discreteness and eventual-constancy theorems
\emph{only} for small subsets or small index types, and ship the larger-index
counterexample \cref{found:ex:largerindex} alongside them. Workspace embeddings
get separate algebraic, order, summability and topological compatibility
theorems; the last is not automatic. After the normal-form bridge, add canonical
finite sine and cosine, the finite-angle quotient and the circle logarithm; add
a real exponential as a separately proved structure; for a global exponential,
require a specified character or the stronger initial-integer-part structure of
\cref{found:sub:canonicalexp}. Do not package all of this into one enormous
typeclass: a theorem about Heron's identity should not acquire an accidental
dependency on a global exponential or a class-choice axiom.

\Needspace{14\baselineskip}
\begin{center}\small
\begin{tabular}{@{}P{0.28\textwidth}P{0.63\textwidth}@{}}
\toprule
Proposed module group (M2) & A meaningful completion test\\\midrule
\texttt{Algebra/Quadratic} & Pair multiplication, conjugation, inverse, norm-square identities.\\
\texttt{Hahn/Support} & Strong sum API, positive-support substitutions, smallness preservation.\\
\texttt{Hahn/Germs} & Formal evaluation and recentering with explicit domain witnesses.\\
\texttt{Hahn/Coherent} & Coefficientwise identity, differentiation, and contour theorems.\\
\texttt{Surreal/Bridge} & Normal-form comparison with algebra and order compatibility.\\
\texttt{Surreal/FineTopology} & Small-index statements plus the larger-index counterexample.\\
\texttt{Surcomplex/Angles} & Full unit-circle parametrization and finite-angle kernel.\\
\texttt{Surcomplex/Enrichment} & Real exponential, strip functions, specified global phases.\\
\bottomrule
\end{tabular}
\end{center}

\subsection{Proposed module layers with contracts}
\label{found:sub:contracts}
M1's nine-row table gives the complementary view: each proposed layer with the
contract it must establish and its typical dependencies. Each contract should
state its universe parameters, its hypotheses, and whether it concerns a formal
sum or a topological limit. A theorem using ordinary residue duality should take
the exact proved structure as input, rather than a loosely named ``analytic
field'' typeclass that silently includes every desired conclusion.

\Needspace{16\baselineskip}
\begin{longtable}{@{}P{0.24\textwidth}P{0.40\textwidth}P{0.28\textwidth}@{}}
\toprule
Proposed layer (M1) & Contract to establish & Typical dependencies\\
\midrule\endhead
\texttt{Complexify} & Pair or quadratic-quotient field; conjugation; squared modulus & Ordered-field algebra\\[0.35em]
\path{Hahn.Support} & Admissible-family certificates; reindexing; finite coefficient dependency & Mathlib Hahn and well-order APIs\\[0.35em]
\path{Hahn.MultivariateEval} & Evaluation at finite tuples of positive-order series & Finite-word support lemma; multivariate formal series\\[0.35em]
\path{Analytic.CoeffGerm} & A precise convergent complex germ ring and derivative maps & Ordinary complex analysis\\[0.35em]
\path{Analytic.HahnGerm} & $\Ag_{n,\Gamma}$, evaluation, translation, finite jets & Previous two layers\\[0.35em]
\path{Analytic.Division} & Support-controlled division and preparation & Ordinary Weierstrass division; operator inversion\\[0.35em]
\path{Analytic.FiniteAlgebra} & Finite bases, residue pairing, trace, local lengths & Division; ordinary residue inputs; linear algebra\\[0.35em]
\path{Workspace.Extension} & Embeddings and finite-quotient base change & Field and algebra maps; support compatibility\\[0.35em]
\path{Surreal.Interpretation} & Small-support normal forms and class/universe interpretation & Existing surreal development; normal-form bridge\\
\bottomrule
\end{longtable}

\subsection{Encode the permitted size of a cut in its data}
\label{found:sub:cutdata}
Two of the three audits present a \texttt{SmallCutData} structure, and they
differ only in whether the value universe is left general or already fixed. We
print M1's two-universe form as the \emph{general} form, and record M3's as the
intended instantiation; M1 itself prescribes $v=u+1$, so this is a
specialization, not a rival.

\begin{samepage}
\begin{lstlisting}
universe u v

-- General form (M1): indices in u, values in a separate universe v.
structure SmallCutData (X : Type v) (lt : X -> X -> Prop) where
  Left : Type u
  Right : Type u
  left : Left -> X
  right : Right -> X
  separated : forall l r, lt (left l) (right r)
\end{lstlisting}
\end{samepage}

\noindent The intended instantiation takes $v=u+1$ and a typeclass order, which
is exactly M3's form:

\begin{samepage}
\begin{lstlisting}
universe u

-- Intended instantiation (M3): the value universe is fixed at u+1.
structure SmallCutData' (A : Type (u + 1)) [LT A] where
  Left : Type u
  Right : Type u
  left : Left -> A
  right : Right -> A
  separated : forall l r, left l < right r
\end{lstlisting}
\end{samepage}

Either structure records cut \emph{input} data; neither asserts that every such
cut is realized. A constructor consuming it can fill exactly the
lower-universe-small families. What must \emph{not} be exposed is a filler for
every unrestricted \texttt{Set A}, and it would be invalid to introduce a
realization axiom uniformly for all relations and carriers, or to instantiate the
interface at an illicit same-size self-family. One may also work with a
predicate-based set of options plus a \texttt{Small} witness for its subtype; the
two interfaces should be related by a proved reindexing theorem.

M2 ships no such structure. Instead it proves, over \texttt{Set F} with
\texttt{[Preorder F]}, a subset-based negative guard --- the version it actually
delivered as Lean source, uncompiled:

\begin{samepage}
\begin{lstlisting}
-- Subset-based variant (M2), shipped as Lean source; compiled for this report, see \S2.5.
theorem noUnrestrictedUpperBounds (F : Type u) [Preorder F] :
    Not (forall s : Set F, exists x : F, forall y in s, y < x) := by
  intro h
  obtain (x, hx) := h Set.univ
  exact (lt_irrefl x) (hx x (Set.mem_univ x))
\end{lstlisting}
\end{samepage}

\noindent M1's corresponding guard uses \emph{predicates} rather than subsets of
a set, and imports \texttt{Init} only:

\begin{samepage}
\begin{lstlisting}
-- Predicate-based variant (M1), shipped as Lean source; compiled for this report, see \S2.5.
theorem noUniversalStrictBound
    {X : Type u} (lt : X -> X -> Prop)
    (irrefl : forall x, Not (lt x x))
    (bound : (X -> Prop) -> X)
    (above : forall A x, A x -> lt x (bound A)) : False := by
  let all : X -> Prop := fun _ => True
  exact irrefl (bound all) (above all (bound all) True.intro)
\end{lstlisting}
\end{samepage}

Both are the machine-level form of \cref{found:prop:allcuts}. Neither had
been run when the archives were prepared; both have since been compiled and
axiom-audited for this report, and each depends on no axioms whatever
(\cref{found:sub:compiled}). A smallness test belongs in the statement
of every analogue of \cref{found:lem:bounds} and \cref{found:thm:discrete}: for
a function from a small indexing type the range is small, but for a subtype of
$\No_u$ smallness is an additional assertion, not a fact inferred from the word
\texttt{Set}.

\begin{remark}[Name the two cut interfaces distinctly]\label{found:rem:twocuts}
By \cref{found:sub:cutfile}, ``cut'' already names two different implemented
objects. A formal API should give the option-pair interface and the
whole-order-section interface distinct names, even though the literature calls
both ``cuts.'' Add to this \cref{found:rem:setrawgame}: a \texttt{Set RawGame}
constructor breaks positivity, because the predicate puts the recursive type in
a function domain.
\end{remark}

\subsection{Complex pairs as independently reusable algebra}
\label{found:sub:complexpair}
The coordinate data can be introduced with no surreal-specific assumptions:
\begin{samepage}
\begin{lstlisting}
universe u
structure ComplexPair (F : Type u) where
  re : F
  im : F
namespace ComplexPair
variable {F : Type u} [CommRing F]
def multiply (z w : ComplexPair F) : ComplexPair F :=
  { re := z.re * w.re - z.im * w.im
    im := z.re * w.im + z.im * w.re }
end ComplexPair
\end{lstlisting}
\end{samepage}
The next steps are coordinate extensionality, ring laws, the real embedding,
conjugation, and the quadratic norm $x^2+y^2$; prove positivity over an ordered
field, then implement inversion and its laws. A quadratic quotient is an
alternative with different API advantages. \emph{Neither design should inherit
the coordinatewise product-ring instance} (\cref{found:rem:productring}). Real
closedness and algebraic closedness should enter only where needed: construct
the field before introducing a square-root-valued modulus, and prove polynomial
algebra generically over the algebraically closed extension only after the
real-closedness bridge is available.

\subsection{Formal sums carried by certified data}
\label{found:sub:certifieddata}
Reuse \path{HahnSeries.SummableFamily} when its hypotheses match the desired
operation. For a full surreal-support interface, add the
lower-universe-smallness certificate where it is not already implied by the
input types. If $I$ is small and each $A_i$ has small support, prove smallness
of the dependent union as a \emph{separate lemma}; do not leave it to an
informal cardinality argument at each use (M3).

The central transport results should have these meanings: coefficients of a sum
are finite coefficient sums; embeddings commute with the certified sum; a
positive correction has inverse supported in its generated monoid; and
reindexing preserves the sum under the stated finiteness conditions. These are
stronger and more precise than declaring a generic infinite-sum notation and
hoping elaboration will choose the intended convergence theory.

M1's finite-coefficient-dependency interface is the reusable lemma to aim for:
for prescribed source support $S$, positive perturbation support $E$ and
exponent $\gamma$, there is a \emph{finite} set of words
$(s,e_1,\ldots,e_r)$ with total $\gamma$, and that finite set controls all
allowed operator expansions contributing to the coefficient. The statement should
make the finiteness witness reusable, supporting the proof pattern of
\cref{found:sub:supportcertificates}. A separate theorem should handle
$K_\Gamma$-linearity of an operator initially extended coefficientwise from a
complex-linear map; it follows from finite convolution and admissible
rearrangements, not from the map's being called coefficientwise.

\subsection{Representing ordinary analytic germs without a global radius field}
\label{found:sub:germrepresentation}
M1 compares two representations. One uses ordinary function germs at zero, for
instance a germ construction on the neighbourhood filter, restricted to germs
with an analytic representative; equality is eventual equality on an ordinary
neighbourhood, and differentiation and multiplication must be proved
representative-independent. The other represents a convergent germ by a
multivariate formal power series together with an existential positive
convergence radius, making coefficient extraction, finite jets and formal
differentiation transparent, but still requiring a theorem identifying these
codes with ordinary holomorphic germs in the desired number of variables.

The design rule is the important part: for $\Ag_{n,\Gamma}$, use a Hahn series
with this germ ring as coefficient type, and \emph{do not} make the radius a
global field of the Hahn-series structure. Such a field would enforce the
excluded common-radius hypothesis and destroy the first separator of
\cref{found:sub:fourrings}. A radius can instead be selected locally when a
finite list of coefficient germs is being compared with an ordinary analytic
theorem.

\subsection{Finite-dimensional algebra before analytic topology}
\label{found:sub:finitefirst}
Once a finite basis of a quotient has been obtained, ordinary matrix APIs express
multiplication, characteristic polynomials, traces and nondegenerate bilinear
forms, and the zero-count descent of \cref{found:thm:finitepoints} becomes a
generic algebraic theorem needing no global surreal metric space and no theory of
fine-continuous paths. The effective order of attack: verify the
support-controlled division operator; prove spanning of the quotient; prove
independence using the residue Gram matrix \eqref{found:eq:residue}; then
establish trace and base-change results. AG explicitly separates spanning from
independence, and preserving that separation in Lean avoids accidentally assuming
the constant dimension one is trying to prove (M1).

\subsection{The normal-form bridge has several independent obligations}
\label{found:sub:bridge}
A usable normal-form bridge needs more than a map between two carriers (M3). It
must prove injectivity and surjectivity, agreement with canonical monomials,
preservation of addition and multiplication, order compatibility on real
coefficients, and preservation of the specified strong sums. To support the
analytic manuscripts one also needs coefficient extraction, standard part and
Taylor evaluation to commute with the bridge. If the bridge is first established
at a fixed exponent group, prove its coherence under group enlargement; if at a
universe-relative small-support Hahn type, prove the corresponding smallness and
universe-lift compatibility. These are mathematical obligations, not cosmetic
coercion lemmas: a noncanonical field isomorphism may fail to preserve the
distinguished monomials or the strong summation.

\subsection{Three implementation strategies, four implementation strategies,
and the recommendation to combine}
\label{found:sub:strategies}
The two audits that surveyed implementation strategy carved the space
differently, and both survive.

M1 compares three:
\begin{itemize}
\item \textbf{Game-first} follows the existing Lean representation most
directly. Strengths: recursive arithmetic and compatibility with current game
theory. Costs: representative management, and the later normal-form bridge.
\item \textbf{Sign-first} makes canonical equality, birthdays and simplicity
transparent. Costs: defining arithmetic efficiently, and proving agreement with
the well-developed cut interface. Especially appropriate for foundational
questions about construction history.
\item \textbf{Hahn-first} is strongest for the analytic goals. Strengths:
coefficientwise algebra, explicit support, valuation, finite-dimensional
methods. Limitation: an abstract Hahn field is not by itself the full surreal
field.
\end{itemize}
M1's conclusion is that the recommended project \emph{combines} game and sign
interfaces for interpretation with Hahn-first proofs for analysis, rather than
choosing one representation exclusively.

M3 tabulates four:

\Needspace{12\baselineskip}
\begin{longtable}{@{}P{0.21\textwidth}P{0.34\textwidth}P{0.36\textwidth}@{}}
\toprule
Strategy (M3) & Best use & Principal proof burden\\
\midrule\endhead
Extend the current surreal library & Canonical numbers, birthdays, small cuts, full-scale theorems & Validate the current APIs; locate or prove real-closedness and normal-form bridges.\\[0.55em]
Develop over generic Hahn fields & Polynomial valuations, support recursion, coherent coefficients & Prove the needed closure hypotheses and later embed the workspace into the intended surreal carrier.\\[0.55em]
Use a bounded surreal field & A concrete set-sized target with strong algebraic closure & Prove the chosen birthday cutoff and all operations stay inside it; expose its scale limitations.\\[0.55em]
Internal set-theoretic model in Lean & Formal foundational interpretation and comparison of set axioms & Bridge modeled sets and operations to external algebraic types; maintain internal/external semantics.\\
\bottomrule
\end{longtable}

\noindent M3's recommendation is a \emph{hybrid of the first two}: the polynomial
and support lemmas can progress over generic fields while the full surreal bridge
is developed. A statement whose proof really needs all small cuts must wait for
the corresponding carrier theorem rather than being smuggled into a generic Hahn
interface as an axiom. The two recommendations agree in substance: do not commit
the entire project to one representation.

\subsection{A theorem-sized implementation sequence}
\label{found:sub:sequence}
M3's sequence, with its useful failure mode. First lock the toolchain and rebuild
the existing carrier and field imports. Next prove or locate the needed
small-cut and universe-lift laws, and independently implement generic
complexification. Then establish the real-closedness or algebraic-closure input,
followed by the normal-form and support transport bridge. In parallel, develop
generic polynomial identities and the finite residue pairing. After that, prove
support-local recursion, construct ordinary holomorphic coefficient rings, and
formalize coherent evaluation and fixed-domain division. Only then combine the
layers for automatic Hartogs coherence and all-scale rigidity.

The useful failure mode: missing normal-form or closure theorems do not block
generic algebra, but they do prevent incorrectly advertised surreal
specializations. Each milestone should produce named theorems with audited
hypotheses rather than a growing collection of unproved global instances. M1's
ordering agrees: begin with multivariable positive-order evaluation and the
support-controlled operator calculus, which directly address the analytic needs
at the set-field level; then formalize germ division, finite-dimensional
deformation, residues and finite base change; and let the full-class
interpretation preserve only the properties for which compatibility has been
proved.

\begin{remark}[The closing rule]\label{found:rem:closingrule}
A dependency gate that has not yet been proved must remain a clearly stated
hypothesis in downstream work. It must never be replaced by a new unexplained
axiom called a theorem. Declaring one axiom named after each desired theorem
makes later deductions \emph{conditional}; it does not turn a manuscript into a
verified result. Conversely, an explicitly quantified hypothesis --- such as a
proved ordinary Weierstrass-division structure --- is not an inconsistency: it
makes the theorem conditional until that structure is instantiated. The final
analytic theorems remain conditional until their ordinary complex-analytic
dependencies, in particular Hartogs and Baire, have also been checked in the
chosen dependency revision.
\end{remark}

\section{Trust, relative consistency, and reproducible verification}
\label{found:sec:trust}

\subsection{What a construction can guarantee}
\label{found:sub:relativeconsistency}
If a base theory $T$ proves that a definable structure satisfies a collection of
axioms $A$, the interpretation is a \emph{relative} consistency justification: a
contradiction obtained from the interpreted mathematics would yield a
contradiction in $T$. Definitional extensions add no new mathematical existence
assumptions beyond what their definitions and existence proofs establish. A
construction in ZFC shows that the added surreal notation can be interpreted in
that foundation without an additional surreal-specific axiom; a universe-relative
type construction gives a corresponding interpretation relative to the chosen
type theory and any semantic universe assumptions used to justify it; and an
axiomatic surreal interface becomes safer when one proves that a known
construction inhabits it.

None of these is an absolute proof that ZFC, GBC, KM, a universe principle or
Lean's full foundations are consistent. A strong effectively axiomatized theory
cannot in general prove its own consistency under the usual hypotheses of the
second incompleteness theorem, and a Lean proof of a model or interpretation
carries its ambient foundation and universe assumptions. It would also be
incorrect to infer a specific minimal consistency strength merely from the use of
universe levels in source code: a set-theoretic model with inaccessibles is a
sufficient \emph{semantic} account of Lean universes, but not a theorem that the
kernel assumes the corresponding large-cardinal sentences
(\cref{found:sub:leanuniverses}). These qualifications are not reasons to
distrust surreal arithmetic; they are the same qualifications that apply to
ordinary formalized mathematics.

A checked construction of an internal model can yield a relative consistency
implication \emph{after} the appropriate syntactic soundness theorem has been
formalized. It does not imply that Lean's own foundation is absolutely
consistent, and there is no conflict with incompleteness in proving a weaker
theory's consistency in a stronger metatheory. The relevant statement records
the metatheory and the model interpretation instead of presenting a circular
guarantee (M1).

\subsection{Kernel checking and mathematical faithfulness are separate}
\label{found:sub:faithfulness}
M2 states the point that no amount of kernel checking addresses. Lean's kernel
checks proof terms against definitions, inference rules and admitted axioms, and
tactics that synthesize such terms are not substitutes for the kernel. But a
final theorem \emph{can be kernel-checked and still mathematically
misidentified} if its formal statement uses the wrong notion of summability,
omits a smallness hypothesis, strengthens the base assumptions, or proves the
result only for a restricted surrogate field.

A reliable verification report therefore records both a successful build and a
statement-level correspondence with the intended theorem. It should say whether
the carrier is $\No_u$, a birthday fragment, a Hahn workspace or an abstract
real-closed field, and whether a function is fine-local, Hahn-coherent, or
equipped with a chosen global phase (\cref{found:rem:fourtypes}).

\subsection{Axioms should be audited transitively, not hidden}
\label{found:sub:axiomaudit}
Lean's official documentation distinguishes its logical axioms --- choice,
propositional extensionality, quotient soundness --- from computational trust
mechanisms and from unfinished proofs, and explains how to inspect a theorem's
axiom dependencies \cite{found:LeanAxioms,found:LeanValidation}. A declaration
accepted by the parser is not a proof, and a theorem depending on
\texttt{sorryAx} is not a completed result: that axiom can prove arbitrary
statements and does not belong in a finished proof. The relevant pattern is
\begin{samepage}
\begin{lstlisting}
#print axioms theoremName
\end{lstlisting}
\end{samepage}
with \texttt{theoremName} the actual exported theorem being audited.

The audit must extend through imports and must inspect the \emph{transitive}
dependency set, not merely search the final file for the word \texttt{axiom}. A
trusted library theorem, an imported assumed instance, or a computationally
generated axiom can carry a dependency invisible in the local proof script.
\emph{Absence of the token \texttt{sorry} in one short file is not enough if a
dependency contains a placeholder or custom postulate.} Tactics that produce
ordinary proof terms leave those terms to the kernel; native evaluation can
involve additional trust in the evaluation or compilation mechanism, and a
project should state whether its certification target permits such dependencies
rather than treating all successful tactic invocations as identical trust
events.

\begin{remark}[What an axiom listing is not]\label{found:rem:axiomlisting}
An axiom listing is \emph{not} a complete description of the kernel's rules or
of the semantic strength of its universe hierarchy, and an empty list does not
establish the consistency of the kernel. It is, however, a valuable way to
detect unproved placeholders or unexpected project-specific axioms in a
purported theorem (M2).
\end{remark}

\subsection{What to record for a reproducible formal result}
\label{found:sub:reproducible}
M1's checklist. A publishable verification record should contain: the exact
repository and dependency revisions; the toolchain; a clean build command and
its exit status; theorem names and types; the transitive axiom report; and the
precise interpretation theorem linking the implementation to the mathematical
statement. It should distinguish kernel-checked proof terms from numerical
experiments, tactic benchmarks and external computation.

Useful regression tests for the supplied manuscripts include a finite zero
cluster, a higher-rank positive-support geometric series, the non-common-radius
coefficient example of \cref{found:sub:fourrings}, the rescaling obstruction
\cref{found:ex:rescaling}, and a negative test rejecting an unrestricted cut.
These test different layers. \emph{Passing a finite example does not prove a
general Hahn-support lemma; rejecting a self-cut does not prove real
closedness.} The full merged regression suite is
\cref{found:app:shortcuts}, and it consists of \emph{specification checks}, not
randomized evidence for any general theorem.

\subsection{This report's own verification status}
\label{found:sub:ownstatus}
Stated once, in \cref{found:sec:boundary}, and not repeated here beyond the
two sentences that govern everything above. Of the three source audits: three
read-only inspections, zero compilations, zero rebuilds, zero axiom audits,
nothing kernel-checked, the \LaTeX{} builds the only executed step. Of the
merged report: three of the five shipped Lean files compiled and
\texttt{\#print axioms} run on every declaration in them, under a toolchain
that is not the pinned one, as recorded in \cref{found:sub:compiled} --- and
nothing else machine-checked, no repository rebuilt, no manuscript verified. The per-audit build evidence and the
per-archive inventories are in
\cref{found:sub:executed,found:sub:archives}, together with M2's instruction not
to run \texttt{lake update} and then describe the result as a check of the
pinned revision, and the record of the missing \texttt{SHA256SUMS} file.

\section{Comparison and recommendations}
\label{found:sec:comparison}

\subsection{One reconciled foundation table}
\label{found:sub:foundationtable}
The three audits produced three overlapping foundation-comparison tables --- nine
rows (M1), ten rows (M2), seven rows (M3). Printing all three would be a third
kind of duplication. They are reconciled here into one wider table, with a final
column recording which audits rated each foundation.

\Needspace{22\baselineskip}
\begin{longtable}{@{}P{0.17\textwidth}P{0.30\textwidth}P{0.34\textwidth}P{0.10\textwidth}@{}}
\toprule
Foundation & Principal advantage here & Principal qualification to keep explicit & Rated by\\
\midrule\endhead
ZF / ZFC with definable classes & Economical classical construction; canonical set-coded numbers; full $\No$ as a definable class; local set recursion; no inaccessible required & Class functions and global assertions are formula-based; quantification over arbitrary classes becomes a definability schema; every recursive graph and set-bounding step must be justified; choice use depends on the proof & M1, M2, M3\\[0.5em]
GB / GBC / specified NBG & Direct language for $\No$, $\SC$ and fixed global maps; explicit class-function parameters; replacement on ordinary domains; closest to the manuscripts & Classes cannot be elements; predicative comprehension only; no unrestricted class-member quotients or collection of all classes; arbitrary ETR and hyperclasses are not automatic; the choice convention must be specified & M1, M2, M3\\[0.5em]
KM / MK and stronger class theories & Impredicative class comprehension; stronger class recursion and satisfaction constructions & Greater foundational commitment than the present local algebra requires; still no unrestricted cuts, no naive hyperclasses, no class of all proper classes; class collection and other variants must be named & M1, M2, M3\\[0.5em]
Grothendieck universes / TG & Treat one level's classes as larger sets; externally set-sized but internally large surreal systems; systematic size closure; rich set operations & Additional universe assumptions beyond ZFC; smallness must survive every change of level; no single small carrier fills all of its external cuts; relative totality is not absolute totality & M1, M2, M3\\[0.5em]
Feferman reflection (ZFC/S) & Some large-object reasoning without blanket inaccessibles & Definability and family-closure conditions are not those of a full universe; conservativity rests on each proof using finitely many instances & M1 only\\[0.5em]
Birthday fragments & Concrete set-sized construction stages; clear ordinal bookkeeping; explicit real-closed fields at suitable epsilon bounds & An arbitrary stage need not be a field (dyadic rationals); analytic closure and closure under arbitrary families require specific theorems; full surreal scales are not automatic & M2, M3\\[0.5em]
Set-sized Hahn fields & Direct, modular access to support-controlled algebra and analysis; specified value group; transparent supports & Not the entire surreal universe; full-class rescaling and topology do not automatically survive restriction (\cref{found:sec:twoseries}) & M2\\[0.5em]
Classical dependent type theory & Native arithmetic libraries, quotients, well-founded recursion, explicit universe polymorphism; small-branching raw games and numeric quotient & Smallness certificates, quotient congruence and universe maps are mathematical obligations; quotient representatives are not computational normal forms & M1, M2, M3\\[0.5em]
Constructive set / type theory & Exposes witnesses, apartness and computational content & Classical trichotomy, choice and ordinal identifications need explicit hypotheses; classical order arguments must be replaced or added explicitly & M1, M2, M3\\[0.5em]
Univalent HoTT / higher induction & Values and extensional comparisons generated together; numerical equivalence can be built into identity & Requires the relevant higher inductive-inductive principles with their computation rules; not native Lean equality (\cref{found:prop:univalence}) & M1, M2, M3\\[0.5em]
Native Lean with universes & Existing surreal quotient; strong reusable local algebra & Small supports, lifts and interpretation theorems remain essential & M1\\[0.5em]
Internal \texttt{ZFSet} in Lean & Formalizes the foundational comparison itself; suitable for internal models and conservativity results & More encoding and transport work; bridges to native algebra and analysis; still relative to Lean's metatheory; not needed for every application & M1, M2\\
\bottomrule
\end{longtable}

\subsection{The split recommendation}
\label{found:sub:split}
The spine of the recommendation is M2's \emph{split}, adopted here because it is
the only framing among the three that names the two questions separately.

\paragraph{For the mathematics.} A precisely specified class theory, or an
explicitly definable-class ZFC presentation, organized by M1's three-part
architecture \eqref{found:eq:architecture}: set-sized algebra and analysis, plus
explicit size-aware comparison maps, plus a class or universe interpretation.
The first component carries most proofs; the second prevents accidental changes
of meaning; the third supplies the unbounded interpretation. GBC with explicitly
stated Global Choice is equally acceptable. More powerful class or universe
systems are useful alternatives, not mandatory cures for a paradox, and a change
of foundation should be motivated by an operation actually needed --- such as
manipulating arbitrary families of class functions --- rather than by the sheer
size of a surreal exponent.

\paragraph{For the implementation.} A universe-indexed classical Lean
development with a set-sized Hahn analytic core. Use classical Lean unless
constructive content is itself a research goal. Build on the current separate
surreal library rather than an outdated snapshot. Keep finite polynomial algebra
and Hahn support calculus \emph{generic}, so they remain useful independently of
the full surreal model. Treat the small-support normal-form comparison as a named
bridge. Establish the surcomplex pair or quadratic-field layer generically so it
applies both to abstract real-closed fields and to the chosen surreal
realization. Adopt M3's hybrid of the first two strategies together with M1's
conclusion that game, sign and Hahn strategies should be combined rather than
chosen between (\cref{found:sub:strategies}).

\paragraph{Specific requirements.} Require a proved normal-form bridge, not just
a Hahn-series datatype. Require a proved real-closedness or algebraic-closure
input where root existence is used, not just a field instance. Represent coherent
data on ordinary complex domains with explicit support certificates and the
componentwise sheaf convention \eqref{found:eq:componentwise}. State full-class
topological and all-scale results as lower-universe-small-family theorems, or as
precise class-theoretic schemes. Keep the entire-carrier cut outside the
constructor's domain by design.

\paragraph{These recommendations are compatible, not competing.} One is a
convenient mathematical language; the other is a proof-engineering
representation with a documented interpretation.

\subsection{What all three independently agreed on}
\label{found:sub:convergence}
Three points are common to all three audits:
\begin{enumerate}
\item class theory, or definable-class ZFC, is the right language for the
written mathematics;
\item universe-indexed classical Lean with a set-sized Hahn analytic core is the
right shape for the code;
\item \emph{no large-cardinal axiom is required merely because $\No$ is a proper
class.}
\end{enumerate}

This convergence should be read for exactly what it is: \textbf{three
independent judgements that happen to coincide, not a single argued position
reached by one author and endorsed twice}. The three audits carved the
recommendation three different ways --- M1 as a default organizing language plus
a separate remark about ordinary classical Lean for the local layers, M2 as an
explicit and load-bearing split, M3 as a layered construction plus a choice among
four strategies. This report adopts M2's framing for its clarity, and does not
claim that any single consensus text existed.

\subsection{The hierarchy of completion claims}
\label{found:sub:milestones}
M1's five milestones must not be relabelled as one another:
\begin{enumerate}
\item a legal representation;
\item an ordered field with an interpretation theorem;
\item a normal-form and workspace bridge;
\item the local analytic theorems;
\item the full-class or universe-coherent interpretation of those theorems.
\end{enumerate}
Each is meaningful progress. None should be promoted to the next merely because
a notation or a typeclass name suggests it.

\subsection{The five items that must stay visible}
\label{found:sub:fiveitems}
M3's closing discipline. The most robust formalization keeps five items visible
in every theorem statement:
\begin{enumerate}
\item the carrier's universe or class status;
\item the permissible index sizes;
\item the exact algebraic closure available;
\item the summability certificate;
\item the topology or coefficient structure used.
\end{enumerate}
With those stated, the foundational choices become transparent alternatives
rather than competing informal interpretations of a single ambiguous notation.

\subsection{The unifying rule, and the three meanings of ``all''}
\label{found:sub:unifying}
M2's unifying rule is simple and has substantial consequences:
\begin{center}
\emph{an individual object can have small data without the totality of all such
objects being small.}
\end{center}
Applied consistently it explains why surreal cuts, surcomplex fields, normal
forms and strong sums are legitimate; why full-class analysis differs from
analysis in one field; and why a formalization must preserve the size
annotations rather than erase them.

The surreal numbers do not force unrestricted comprehension (M3). They force
clarity about what is small. Once every cut, quotient, support, family and
function has an explicit size boundary, the apparent paradoxes become either
short contradictions to an overstrong specification or ordinary theorems about
different ambient structures, and surcomplexification adds a finite algebraic
layer rather than a new foundational difficulty. The architecture proposed here
does not weaken the interesting mathematics: it explains why finite polynomial
theory can be developed in an ordinary set-sized field, why positive-support
preparation needs only set recursion, why the Hartogs support argument needs
replacement on the ordinary domain, and why all-scale rigidity uses a resource
unavailable in an arbitrary fixed Hahn field.

The central lesson, in M1's formulation, is not to avoid proper classes. It is
to prevent three different meanings of ``all'' from being confused:
\begin{enumerate}
\item all members of a specified \emph{set};
\item all values at a specified \emph{universe level}, subject to
smaller-index rules;
\item the entire \emph{definable class}.
\end{enumerate}
With those meanings separated, surreal and surcomplex mathematics fits into
rigorous foundations without new set-theoretic paradoxes, while retaining the
scales and support phenomena that motivate the theory.
\appendix

\section{The source-manuscript ledger}
\label{found:app:ledger}

Three audits were merged into this report. They were given different source
manuscripts, and this appendix records exactly which. It is referenced from
\cref{found:sub:ledger}, and it is the authority for the per-audit tagging
convention used throughout the body.

\Needspace{16\baselineskip}
{\small
\begin{description}[leftmargin=2.2em,style=sameline,itemsep=0.7em]

\item[\textbf{A}] \emph{Surcomplex Analysis: Infinitesimal Calculus,
Hahn-Coherent Holomorphy, and Contour Theory}.\\
\textsc{Filenames used:} \path{surcomplex_analysis(3).tex} by M1 and M3;
\path{surcomplex_analysis(2).tex} by M2.\\
\textsc{Recorded SHA-256:} \texttt{e6e2ef4a8b8d920fffa72564}
\texttt{90f98aa312023078fec966c33b3f204cb17748f4}, recorded identically by M2
and by M3; none recorded by M1. 105{,}194 bytes (M3).\\
\textsc{Read by:} \textbf{M1, M2 and M3.}

\item[\textbf{AG}] \emph{Infinitesimal Analytic Geometry over the Surcomplex
Numbers: Radius-Free Hahn Germs, a Nullstellensatz, and Conservation of
Singular Zeros}.\\
\textsc{Filename used:} \path{surcomplex_analytic_geometry.tex}.\\
\textsc{Recorded SHA-256:} none.\\
\textsc{Read by:} \textbf{M1 only.}

\item[\textbf{T}] \emph{Trigonometry on the Surcomplex Plane: Canonical Angles,
Infinite Triangles, Infinitesimal Contact, and Hahn-Analytic Oscillation}.\\
\textsc{Filename used:} \path{surcomplex_trigonometry.tex}.\\
\textsc{Recorded SHA-256:} \texttt{e61577fffcb23c1dabe60a22}
\texttt{57ab2e490a30ab274315c1e6edaeab7639ad4600} (M2).\\
\textsc{Read by:} \textbf{M2 only.}

\item[\textbf{P}] \emph{Surcomplex Polynomial Algebra: Order Geometry, Newton
Profiles, and Scale-Resolved Stability}.\\
\textsc{Filename used:} \path{surcomplex_polynomial_algebra(2).tex}.\\
\textsc{Recorded SHA-256:} \texttt{9d2ad437715ab5a31804a282}
\texttt{6911e5f9138fa394aa330c3735ea3dfbfe516005}; 112{,}766 bytes (M3).\\
\textsc{Read by:} \textbf{M3 only.}

\item[\textbf{R}] \emph{Hartogs Coherence, Finite Maps, and Residue Duality
over the Surcomplex Numbers}.\\
\textsc{Filename used:} \path{surcomplex_research.tex}.\\
\textsc{Recorded SHA-256:} \texttt{9130327331d5635a6e95bb2b}
\texttt{b3f972f92f18fe0dbea9889a8d32f121c857c712}; 98{,}818 bytes (M3).\\
\textsc{Read by:} \textbf{M3 only.}

\end{description}
}

\subsection*{The \texorpdfstring{(2)}{(2)}/\texorpdfstring{(3)}{(3)} filename divergence}
M2 labels the Analysis manuscript \path{surcomplex_analysis(2).tex}, while M1
and M3 label it \path{surcomplex_analysis(3).tex}. The SHA-256 values recorded
independently by M2 and by M3 are \emph{identical}, so it is one file under two
labels. This report does not silently pick a filename, and asserts nothing about
which of the two labels is the correct one. M3 additionally records the file as
105{,}194 bytes.

\subsection*{Consequences of the ledger}
\begin{itemize}
\item \textbf{A is the only manuscript read by all three audits.} Claims derived
from A may be attributed jointly.
\item \textbf{Every claim derived from AG, T, P or R rests on a single audit.} In
the body, each such claim carries the tag of that single audit: (M1) for AG,
(M2) for T, (M3) for P and R.
\item No sentence in this report says ``the supplied manuscripts'' when it means
a text that only one member saw.
\item \textbf{A fourth, unsupplied manuscript.} P refers to a separate
\emph{global-divisors} manuscript that was not supplied for any of the three
tasks. Its contents are not treated as inspected evidence and no independent
assessment of it is offered (M3).
\item \textbf{Other archives were not examined.} Archives named in the
bibliographies of the supplied manuscripts, in particular in T's bibliography,
were not part of the materials examined (M2); the same holds for M1 and M3.
\item All five are user-supplied research expositions, not independently
certified libraries or refereed publications. Their original files are not
duplicated in any of the three archives, and none of them was machine-verified.
\end{itemize}

\subsection*{The three merged audits}
For completeness, the three audit archives merged into this report are recorded
in this repository's provenance material as
\texttt{01-size-ledger-and-universe-audit.tex} (M1, 110{,}177 bytes, 1073
lines), \texttt{03-set-and-type-theory-comparison.tex} (M2, 100{,}299 bytes, 923
lines) and \texttt{06-classes-universes-and-lean-audit.tex} (M3, 125{,}084
bytes, 1109 lines), together with their READMEs and audit records. M3 also
recorded SHA-256 values for its own \texttt{.tex}
(\texttt{0f51c217e7b090089c245c46b0f8985cc16587d94b3d0796d5ddb361fcb31e17}) and
its own \texttt{.pdf}
(\texttt{aa7ea0641f50b40de9fe4a5146aefa6291cef0c7e5f53273ac0c83f62b6fc3d6}); it
is the only audit that hashed its own outputs. None of those files is
redistributed with this report; the disposition of every result they contain is
recorded in the provenance material that remains.

\section{Forbidden shortcuts and regression tests}
\label{found:app:shortcuts}

Everything in this appendix consists of \textbf{specification checks}. None of
it is randomized evidence for a general theorem. These lists are especially
effective immediately before a mathematically strong typeclass instance is about
to be introduced. Passing a finite example does not prove a general Hahn-support
lemma; rejecting a self-cut does not prove real closedness.

\subsection{Twelve forbidden shortcuts}
\label{found:app:twelve}
M2's list, reproduced in full. Each item records a specific size, domain,
algebraic or verification obligation established in the body; none is an ad hoc
prohibition.
\begin{enumerate}
\item Do not apply the cut operation to all surreals at its own operative size
level. (\cref{found:prop:allcuts})
\item Do not treat a proper-class equivalence class as an element of an NBG
class. (\cref{found:sub:classquotient})
\item Do not identify surreal cut interpolation with Dedekind completeness.
(\cref{found:prop:incomplete})
\item Do not use an arbitrary birthday bound as a field instance without closure
proofs. (\cref{found:sub:cutoffs})
\item Do not identify the self-Hahn normal-form theorem with an unrestricted
function-space definition. (\cref{found:sub:hahnworkspace})
\item Do not confuse well-ordered support with small support, or either with
finite coefficient incidence. (\cref{found:sub:certificates})
\item Do not replace strong Hahn summation by a topological \texttt{tsum}.
(\cref{found:sub:tsum})
\item Do not infer workspace stability from the fact that the inputs lie in a
workspace. (\cref{found:sub:twohalves})
\item Do not remove the small-index qualification from the fine-net theorem.
(\cref{found:ex:largerindex})
\item Do not substitute a real-valued norm for the actual surreal-valued
modulus. (\cref{found:sub:modulus})
\item Do not derive global analytic uniqueness from a local differential
equation alone. (\cref{found:rem:fourtypes})
\item Do not call illustrative source, symbolic checks, or a module comment a
machine-checked proof of the intended theorem.
(\cref{found:sec:boundary}, \cref{found:rem:symbolic})
\end{enumerate}

\subsection{A ten-row diagnostic table}
\label{found:app:regressiontable}
M3's table pairs a proposed statement with the exact diagnostic input that
refutes it.

\Needspace{18\baselineskip}
\begin{longtable}{@{}P{0.36\textwidth}P{0.57\textwidth}@{}}
\toprule
Proposed statement & Diagnostic input or contradiction\\
\midrule\endhead
Every pair of subsets of the carrier has a Conway separator & The carrier itself on the left and the empty set on the right (\cref{found:prop:allcuts}).\\[0.45em]
The ordered surreal field is algebraically closed & The polynomial $X^2+1$ (\cref{found:sub:realclosed}).\\[0.45em]
The surreal field is Dedekind complete & The set of natural numbers; subtract one from a proposed least upper bound (\cref{found:prop:incomplete}).\\[0.45em]
Every well-ordered support defines a same-level surreal & The too-large ordinal-indexed support \eqref{found:eq:toolarge}; require relative smallness.\\[0.45em]
Individual legal supports imply a legal family sum & The singleton supports $\{1/n\}$ (\cref{found:ex:failed}).\\[0.45em]
Well-ordered union suffices for a Hahn family sum & Infinitely many copies of the constant series one (\cref{found:ex:failed}).\\[0.45em]
Every internal all-scale coherent function is polynomial & \Cref{found:ex:internal} in $\C((t^\Q))$ --- note the \emph{plus} sign, and see \cref{found:sub:complementary}.\\[0.45em]
Every subset of a universe-relative surreal type is discrete & The entire carrier; only the lower-universe-small subsets satisfy \cref{found:thm:discrete}.\\[0.45em]
The product ring on two real-field copies is surcomplex & $(1,0)(0,1)=0$ in the product ring, unlike complex multiplication (\cref{found:rem:productring}).\\[0.45em]
Every algebra homomorphism preserves infinite Hahn evaluation & Require a strong-sum preservation theorem, or replace the step by the finite divided-difference identity \eqref{found:eq:divideddifference}.\\
\bottomrule
\end{longtable}

\subsection{A layered regression suite}
\label{found:app:layeredsuite}
M1's suite tests different layers, and passing one says nothing about another:
\begin{enumerate}
\item a finite zero cluster (finite algebra layer);
\item a higher-rank positive-support geometric series (support layer);
\item the non-common-radius coefficient example
$\sum_{r\ge1}t^{r\eta}/(1-rz)$, which lies in $\Ag_{1,\Gamma}$ but has no common
ordinary radius (coefficient-ring layer, \cref{found:sub:fourrings});
\item the rescaling obstruction \cref{found:ex:rescaling}, with its
\emph{minus} sign (workspace layer);
\item a negative test rejecting an unrestricted cut (interface layer).
\end{enumerate}
To these the merge adds two further separators that were not in any single
suite: the valuation-versus-modulus separator \eqref{found:eq:separator}, and
the disjoint-disks sheaf necessity example \cref{found:ex:disjointdisks}.

\section{Dependency ledgers: three complementary views}
\label{found:app:dependencies}

The three audits each produced a dependency view of the same graph, from a
different direction. All three are printed, because none subsumes another.

\subsection{Construction, input, and what it does not supply}
\label{found:app:ledgerthree}
M3's eleven-row ledger.

\Needspace{20\baselineskip}
\begin{longtable}{@{}P{0.28\textwidth}P{0.34\textwidth}P{0.30\textwidth}@{}}
\toprule
Construction or theorem & Explicit input & What it does not supply\\
\midrule\endhead
Canonical sign carrier & Ordinal-indexed two-sign functions; bounded function sets & Field operations or normal form.\\[0.45em]
Numeric-game quotient & Well-founded game relation, numericity, equivalence, congruence & Smallness of arbitrary subsets of the quotient.\\[0.45em]
Small-cut filling & Small separated option families; birthday control & Filling arbitrary proper-class or larger-universe cuts.\\[0.45em]
Complexification & Ordered-field laws; pair multiplication & Algebraic closedness without real closedness.\\[0.45em]
Hahn field & Ordered exponent group and coefficient field; support calculus & Normal-form equivalence with the surreal carrier.\\[0.45em]
Algebraically closed workspace & Divisible group; algebraically closed coefficients; Hahn closure theorem & Every full-scale analytic theorem.\\[0.45em]
Strong Hahn evaluation & Well-ordered union; finite coefficient dependence; small support & Topological convergence of partial sums.\\[0.45em]
Positive-support recursion & The set $E^*$, finite convolution, an invertible linear step & General class-valued ETR.\\[0.45em]
Automatic Hartogs support & Set-indexed holomorphic coefficients, Baire and choice; class replacement for the initial restriction & Coherence when only ordinary slices have been specified.\\[0.45em]
Full-scale polynomial rigidity & One chart support per scale; a bound for the entire small valuation family & The same conclusion in a fixed value group such as $\Q$ (\cref{found:ex:internal}).\\[0.45em]
Kernel-verified specialization & All required instances constructed, dependency audit, pinned build & Absolute consistency of the foundational kernel and axioms.\\
\bottomrule
\end{longtable}

\subsection{Theorem-by-theorem audit of the source manuscripts}
\label{found:app:theorembytheorem}
M2's eleven-row audit. It concerns foundational and representation
requirements; it does \emph{not} certify all analytic proof details in the
original texts. Its rows were drawn from the two manuscripts M2 read (A and T);
the rows are retained with that scope.

\Needspace{22\baselineskip}
\begin{longtable}{@{}P{0.25\textwidth}P{0.32\textwidth}P{0.35\textwidth}@{}}
\toprule
Source result family & Sufficient working layer & Obligation that must remain explicit\\
\midrule\endhead
Finite vector and triangle identities & Generic ordered or real-closed field; quadratic extension & Positive square roots and nonzero denominators; genuine angle interpretation added separately.\\[0.4em]
Strong sums and infinitesimal Taylor lifts & A set-sized Hahn workspace & Well-ordered union, finite coefficient incidence, and correct ordinary base value.\\[0.4em]
Polar decomposition by finite angles & Hahn infinitesimal exponential and logarithm, plus the ordinary circle & Surjectivity on every unit point, exact kernel, and quotient size.\\[0.4em]
Arbitrary-coefficient small-radius evaluation & Full class, or universe-relative small-family boundedness & A dominating exponent may lie outside the initial workspace; \cref{found:ex:rescaling}.\\[0.4em]
Fine-topology degeneracy & The small positive lower-bound principle & Smallness of the subset or net index, not only the existence of a Lean type.\\[0.4em]
Coherent Cauchy and identity theorems & Ordinary holomorphic coefficients on one common domain & Support compatibility, ordinary domain hypotheses, and coefficientwise contour definition.\\[0.4em]
Preparation and root-cluster conservation & Finite polynomial and algebraic data plus support recursion & Algebraic closure; finite coefficient equations; root supports may enlarge beyond factor supports.\\[0.4em]
Finite Fourier and Fej\'er--Riesz results & Finite sums and real-closed-field algebra, with circle interpretation & Ordinary finite degree; positivity on the entire relevant circle, not only its real shadow.\\[0.4em]
Strip exponential and hyperbolic geometry & A Gonshor exponential plus finite phases & The real exponential structure is additional to the bare field interface.\\[0.4em]
Global phase classification & Additive normal-form splitting and a specified class or type homomorphism & Quantify over characters at the right sort; do not form an illegal class of classes (\cref{found:rem:secondsort}).\\[0.4em]
All-scale rigidity & A full class, or universe-relative quantification over scales & No inference from a single fixed workspace; each coefficient family must be small at the operative level.\\
\bottomrule
\end{longtable}

\subsection{A minimal theorem-dependency diagram}
\label{found:app:diagram}
M1's recommended logical order for the analytic-geometry programme:
\[
\begin{gathered}
\text{well-ordered supports and finite co-support}\\
\Downarrow\\
\text{formal sums, evaluation, translation, operator inversion}\\
\Downarrow\\
\text{ordinary division input}\ \Longrightarrow\ \text{Hahn-germ division}\\
\Downarrow\\
\text{finite spanning}\quad+\quad\text{ordinary residue duality}\\
\Downarrow\\
\text{finite basis and perfect pairing}\\
\Downarrow\\
\text{trace and finite base change}\ \Longrightarrow\ \text{zero invariance}.
\end{gathered}
\]
The normal-form embedding into $\SC$ supplies the interpretation boundary; it
does not replace any of these local implications. The topological and all-scale
theorems form a separate branch with additional smallness arguments.

\subsection{Two source-specific checklists}
\label{found:app:checklists}
M1's checklists, retained with their single-audit scope.

\paragraph{For A (read by all three; the checklist is M1's).} Represent
individual surcomplex values by set codes; specify normal-form support and the
sign of the Hahn exponent convention; give class replacement or a
universe-smallness theorem for each family of outputs; encode formal coefficient
sequences as sets; allow the scale chosen by the small-rescaling theorem to
enlarge the initial workspace; and state the fine-topology degeneracy theorem
only for small sets and small nets. Keep fine-local series, common-domain
holomorphic coefficients and all-scale chart conditions distinct. Encode a
global class function by its graph, not as an element of an unconstructed
function space. Give a set-code interpretation to analytic germs before
quotienting. Define the contour functional on its coefficientwise data and prove
its laws without silently replacing it by a path integral in the full fine
topology. \emph{These conditions clarify the intended objects; the substantive
theorems still need their own proof terms.}

\paragraph{For AG (M1 only).} Fix the exponent group before constructing the
germ ring, its ideals and finite quotients. Preserve the ``one ordinary
convergence radius per coefficient germ'' condition rather than imposing a
uniform radius. Make both clauses of strong summability explicit. Prove that
coefficientwise maps preserve support and interact correctly with Hahn scalar
multiplication. Formalize the positive-support inverse independently of
convergence of its partial sums. For finite deformation, prove spanning and
independence separately; define the perturbed residue before using the roots to
interpret it; state ordinary residue duality and finite-parameter trace
identities with exactly the hypotheses needed; prove
\eqref{found:eq:finitebasechange} in the finite situation rather than assuming
global tensor-product equality; and only then invoke
\cref{found:thm:finitepoints} to exclude zeros at new surreal scales.
Distinguish the formal coefficient ring from completion at an infinitesimal
point.

\section{Notation and implementation vocabulary}
\label{found:app:notation}

M1's glossary, extended with items from M2 and M3.

\Needspace{24\baselineskip}
\begin{longtable}{@{}P{0.25\textwidth}P{0.67\textwidth}@{}}
\toprule
Symbol or term & Meaning in this report\\
\midrule\endhead
$\No$, $\SC$ & The full definable class surreal field and its complexification $\SC=\No[i]$\\[0.3em]
$\No_\U$, $\No_u$ & A universe-relative realization; the latter is schematic type-theoretic notation\\[0.3em]
$\No_{<\lambda}$ & The birthday cutoff $\{x:\bd(x)<\lambda\}$, a set in ZFC\\[0.3em]
$\bd(x)$ & Birthday: the ordinal length of the sign sequence\\[0.3em]
$\prec_s$, $\prec_b$ & Simplicity (proper initial segment) and birthday precedence; the second does not imply the first\\[0.3em]
$\mathbin{\#}$ & Hessenberg natural sum, used as a pair-recursion measure\\[0.3em]
$\Gamma$, $K_\Gamma$, $F_\Gamma$ & A set-sized ordered exponent group; its complex Hahn field $\C((t^\Gamma))$; its real Hahn field $\R((t^\Gamma))$\\[0.3em]
$t^\gamma$ & $\omega^{-\gamma}$ after the surreal normal-form interpretation; supports are well ordered in the \emph{increasing} $t$-exponent order\\[0.3em]
$R_n$, $\Ag_{n,\Gamma}$ & Ordinary convergent complex germs $\C\{z_1,\ldots,z_n\}$, and Hahn series of those germs (radius-free)\\[0.3em]
$\Fg_{n,\Gamma}$ & Hahn series with arbitrary formal complex coefficient series\\[0.3em]
$\HH_\Gamma(U)$ & $\OO(U)((t^\Gamma))$: Hahn series whose coefficient functions are holomorphic on \emph{one} specified ordinary open $U$\\[0.3em]
$E^*$ & The additive monoid of finite sums of members of a positive well-ordered support $E$\\[0.3em]
Strong sum & A coefficientwise sum satisfying \emph{both} conditions in \eqref{found:eq:summability}\\[0.3em]
Coherent datum & A Hahn series with coefficients in $\OO(U)$ for one ordinary open $U$; a third structure, distinct from both topologies below\\[0.3em]
Fine topology & The full-surreal-radius neighbourhood language, unless a universe is specified\\[0.3em]
Intrinsic topology & A workspace's own neighbourhood structure, using only its own tolerances $(F_\Gamma)_{>0}$\\[0.3em]
$\Oz$, $\Pi$ & The omnific integers $\Oz=\Pi\oplus\Z$; $\Pi$ the purely infinite normal forms\\[0.3em]
\texttt{SmallCutData} & A structure recording cut \emph{input} data with the index universe separated from the value universe; it records an input, and asserts no realization\\[0.3em]
option-cut vs.\ section-cut & A small pair of Conway options, versus a complementary section of the whole order; the literature calls both ``cuts''\\[0.3em]
\texttt{Small} & Evidence that a type is equivalent to one in a specified smaller universe; it must be \emph{proved}, never manufactured\\[0.3em]
\texttt{Set X} & A Lean predicate on \texttt{X}; not an internal ZFC set by default\\[0.3em]
\texttt{ZFSet} & Mathlib's internal extensional set universe, relative to a universe level\\[0.3em]
\path{ZFSet.Definable} & A representative-level liftability condition, \emph{not} first-order definability by a formula\\[0.3em]
\texttt{noncomputable} & May use noncomputational classical choices; \emph{not} synonymous with an unproved axiom\\[0.3em]
\texttt{sorryAx} & A placeholder dependency that prevents any claim of a completed proof\\[0.3em]
ETR & Elementary transfinite recursion: recursions producing a whole \emph{class} at each stage; genuinely stronger than GBC in its usual forms\\[0.3em]
GB / GBC / NBG & See \cref{found:sub:gbconvention}; the convention is fixed before the letters are used\\[0.3em]
``Paradox-free'' & Only the explicit exclusion of the specific invalid constructions analysed here; expressly not an absolute consistency claim\\
\bottomrule
\end{longtable}

\section{Precise limitations of the present verification}
\label{found:app:limitations}

This appendix is the union of the three audits' limitation lists, deduplicated.
Every item is load-bearing. Dropping any one of them silently upgrades the
report's apparent strength.

\subsection{On verification and build status}
\label{found:app:limbuild}
\begin{enumerate}
\item No Lean or Lake executable was available in any of the three preparation
environments.
\item No Lean file in any archive was compiled or kernel-checked \emph{by
the three audits}. Three of the five were compiled for the merged report; see
\cref{found:sub:compiled}, and item~\ref{found:item:compiledscope} below for
what that does and does not cover.
\item The external Lean repository was not rebuilt.
\item No transitive \texttt{\#print axioms} audit was executed. The
\texttt{\#print axioms} commands present in M1's and M2's files were unexecuted
source in all three archives; their output is now recorded, for three of the
five files, in \cref{found:sub:compiled}. A per-declaration axiom list is not a
transitive audit of the libraries those declarations import.
\item No Mizar run was performed; the Mizar development was not inspected
article by article and its axiom closure was not reconstructed.
\item None of the five supplied manuscripts was machine-verified.
\item The archives are articles and implementation designs, not certified Lean
packages.
\item A successful future execution of the small illustrative examples would not
constitute verification of the full surreal theory or of any manuscript.
\item Source inspection and compilation are different verification levels.
\item Only the \LaTeX{} builds were executed by the audits themselves. That
is the only step common to all three. The merged report adds three Lean
compilations and their axiom probes, and nothing else.
\label{found:item:compiledscope}
\item M2's README refers to a \texttt{SHA256SUMS} file which its archive did
ship and which this repository does not carry, checksum manifests being outside
the collection's scope.
\item M2's clone/build/run commands are a reproduction recipe, not commands that
were executed. Do not run \texttt{lake update} and then describe the result as a
check of the pinned revision.
\end{enumerate}

\subsection{On the pinned source inspection}
\label{found:app:liminspection}
\begin{enumerate}\setcounter{enumi}{12}
\item The inspection covered named files at one revision via the public GitHub
interface, reading source only: no clone, no build, no comprehensive placeholder
search, no axiom audit.
\item Module names alone do not prove their strongest imaginable theorems. This
study did not verify an installed real-closedness theorem or a complete
surcomplex-analysis library by reading module names.
\item \texttt{\#synth Field Surreal} is not proof of real closedness,
normal-form equivalence, or any analytic theorem.
\item A constructor for a Hahn target type, or a comment describing an intended
future isomorphism, is not that isomorphism. Normal-form existence, surjectivity
of evaluation and compatibility with multiplication remain proof obligations.
\item No complete ordered-field normal-form equivalence, no real-closedness
instance for the concrete game-based carrier, no surcomplex algebraic-closure
specialization, and no Gonshor exponential was verified in the inspected
modules. \emph{This is a scope-limited observation, not a claim of their absence
from any other file, branch, repository, project or later revision.}
\item \texttt{IsRealClosed} is an interface, not an instance for any particular
type.
\item \path{ZFSet.Definable} and \path{Classical.allZFSetDefinable} are
representative-level liftability conditions, not first-order definability by a
formula.
\item The phrase ``complete ordered field'' in source prose is not a
Dedekind-completeness theorem.
\item The \path{MvPowerSeries.heval} TODO concerns one inspected module and
its documentation; it is not a proof that no other project contains a
multivariable evaluation theorem.
\item Mathlib and Lean generated or official ``latest'' documentation pages are
\emph{rolling} sources consulted on the inspection date. They are not asserted to
be generated from the pinned dependency commit, and they are not substitutes for
the repository's dependency lock.
\item Old Lean 3 surreal documentation is not a reliable inventory of the
current separate Lean 4 project, and finding a theorem in that project does not
make it available by importing upstream Mathlib alone.
\item Exact typeclass names should be taken from the pinned dependency source
rather than invented to match a mathematical phrase.
\item Theorems do not transfer between Mizar and Lean merely because their
informal names agree; the two differ in foundational and representation choices.
\end{enumerate}

\subsection{On consistency and ``paradox-free''}
\label{found:app:limconsistency}
\begin{enumerate}\setcounter{enumi}{25}
\item All consistency assurances are relative to the chosen foundation and its
metatheory.
\item No absolute consistency proof of ZFC, GBC, KM, any universe principle, or
Lean's foundations is claimed.
\item ``Paradox-free'' means the explicit exclusion of the specific invalid
constructions analysed --- the self-cut, the proper-class-member quotient, the
unrestricted support sum, incompatible equality principles --- and nothing more.
\item The second incompleteness theorem forbids advertising a construction in
ZFC or Lean as an absolute proof that its foundation cannot derive a
contradiction.
\item An internal model gives relative consistency only, and only after the
appropriate syntactic soundness theorem has been formalized.
\item A specific minimal consistency strength cannot be inferred from the use of
universe levels in source code. A set-theoretic model with inaccessibles is a
sufficient semantic account of Lean universes, but not a theorem that the kernel
assumes the corresponding large-cardinal sentences.
\item GBC's conservativity over ZFC for first-order set sentences must not be
strengthened to the claim that every model of ZFC expands to GBC with the same
sets, nor that every model of ZFC has a definable global well-order.
\item Feferman reflection conservativity relies on any one proof using finitely
many instances. It is not the assertion that ZFC proves the existence of a single
transitive set model of all of ZFC in one internal satisfaction statement.
\item No completed translation into NFU or any universal-set theory is claimed.
\item A Grothendieck universe assumption has a cost: an inaccessible is not a
theorem of ZFC, and arbitrarily many universes is stronger still. The
inaccessible bound is a convenient sufficient condition, not a claim of
minimality.
\item A completed Mizar proof establishes a result in Tarski--Grothendieck;
extracting a proof in a weaker theory requires dependency analysis.
\end{enumerate}

\subsection{On the supplied manuscripts}
\label{found:app:limmanuscripts}
\begin{enumerate}\setcounter{enumi}{36}
\item All five manuscripts are user-supplied research expositions, not
independently certified libraries or refereed publications. Their original files
are not duplicated in any archive.
\item Their analytic results are \emph{targets and supplied arguments}, not
imported formal certificates.
\item This report is a foundational audit, not a line-by-line independent
verification of any manuscript's analytic proofs.
\item Other archives referenced in their bibliographies were not examined.
\item The fourth, global-divisors manuscript mentioned inside P was never
supplied, and no assessment of it is offered.
\item The Analysis manuscript A is the only source read by all three audits.
Every claim about the analytic-geometry, trigonometry, polynomial-algebra or
Hartogs manuscripts rests on a single audit.
\item Size-correctness of a ring does not certify its Noetherianity, and
formalization of an abstract field does not certify a proposed analytic residue
formula.
\end{enumerate}

\subsection{On the mathematics proved here}
\label{found:app:limmath}
\begin{enumerate}\setcounter{enumi}{43}
\item The size obstructions, localization arguments, finite-algebra descent,
recursion-locality proposition, support-local recursion, topology
counterexamples and the two $t^{\pm n^2}$ examples are elementary foundational or
algebraic deductions included to make the architecture precise. \emph{No novelty
claim, no priority claim, and no resolution of any published open conjecture is
attached to any of them.}
\item \Cref{found:ex:internal} specifically carries no priority claim and no
machine-verification claim. It is a sharpness result about the scope of the
all-scale rigidity theorem of the companion report at
\texttt{docs/surcomplex/analysis/}, not a refutation of it
(\cref{found:sub:scope}).
\item No sharp reverse-mathematical calibration of normal forms, real
closedness, Hahn algebraic closedness or analytic duality is claimed, and no
least axiom system for any source manuscript is identified. Determining such a
minimum would be a separate investigation.
\item Particular proofs using arbitrary representatives, well-orderings of a
coefficient field, or extraction of decreasing sequences may require choice
principles not analysed here.
\item \Cref{found:thm:workspace} does not say that one set field contains all of
$\SC$, nor that a set of workspaces covers the class.
\item Neither \cref{found:thm:finitepoints} nor its source asserts global
tensor-product equality for the whole infinite-dimensional analytic ring.
\item \Cref{found:prop:recursion} is not a blanket certification of every
informal class recursion.
\item \Cref{found:prop:closure} does not produce closure under all set-indexed
Hahn sums; those have unbounded arity.
\item \Cref{found:thm:discrete} is not a claim that the whole full-class fine
topology is discrete; the carrier itself has no isolated points.
\item \Cref{found:prop:univalence} concludes only that incompatible equality
mechanisms cannot be superimposed --- \emph{not} that univalent mathematics is
inconsistent.
\item The proposed working surreal interface is an interface, not a claimed
minimal or independent axiomatization; establishing categoricity would require a
stated notion of morphism and a proved isomorphism to a canonical construction.
\item The four-part abstract package is a structural specification, not a
proposed minimal finite axiom system.
\item Exact ordinal closure bounds for exponential subfields are theorems from
the literature, not consequences of notation, and must be ported with their
hypotheses and the van den Dries--Ehrlich erratum.
\item The constructive and HoTT routes are not claimed to already supply the
full normal-form, algebraic-closure and analytic package; bare CZF is not
claimed to prove every desired large inductive construction, and constructive
surreal variants are not claimed equivalent to the classical sign-sequence class
without additional logical hypotheses.
\item The cited HoTT, Shulman and Joseph sources do not by themselves constitute
a Lean 4 implementation of the complete surcomplex analytic theory.
\item A categorical small-maps framework does not remove the diagonal
obstruction.
\item The 2026 Hamkins/Chen/Yang program is work in progress with details still
being checked; it is not a completed published axiomatization, not a verified
consistency proof, not a proved dependency of this report, and not a statement
about the pure real-closed-ordered-field language.
\end{enumerate}

\subsection{On design proposals}
\label{found:app:limdesign}
\begin{enumerate}\setcounter{enumi}{60}
\item All proposed module names, layers, architectures, completion criteria and
implementation sequences are design suggestions, not claims that files with those
paths exist.
\item A dependency gate not yet proved must remain a stated hypothesis in
downstream work, never be replaced by a new axiom called a theorem.
\item Declaring one axiom named after each desired theorem makes later
deductions conditional; it does not turn a manuscript into a verified result.
\item Passing a finite example does not prove a general Hahn-support lemma;
rejecting a self-cut does not prove real closedness.
\item Symbolic checks of specific identities reported in a source manuscript are
sanity checks, not substitutes for support theorems, transfinite recursions or
all-scale assertions.
\item The regression tests and forbidden-shortcut lists are specification
checks, not randomized evidence for any general theorem.
\item Absence of the token \texttt{sorry} in one short file is insufficient if a
dependency carries a placeholder or a custom postulate.
\item A theorem can be kernel-checked and still mathematically misidentified if
its statement uses the wrong summability notion, omits a smallness hypothesis, or
proves the result only for a restricted surrogate field.
\item An axiom listing is not a description of the kernel's rules, and an empty
list does not establish kernel consistency.
\item M1's class-relation sheaf encoding is a clarification of a workable
reading of A, not a proof of its analytic gluing theorem.
\item The illustrative Lean files contain no intended claim to implement surreal
arithmetic, Hahn evaluation, Weierstrass preparation or residue theory, and the
complexification sketch claims no \texttt{Field} instance and no real-closedness
theorem.
\item The raw-game constructor shown is an illustration of the size discipline,
not the current repository's exact implementation.
\end{enumerate}

\subsection{On what a full formalization would still require}
\label{found:app:limfuture}
\begin{enumerate}\setcounter{enumi}{72}
\item Several-variable Weierstrass division in the chosen germ model.
\item Support preservation for coefficientwise extension.
\item Finite-dimensional residue duality.
\item Comparison with ordinary finite analytic deformations.
\item Noetherian and local-completion arguments.
\item The precise finite base-change theorem
\eqref{found:eq:finitebasechange}.
\item Coherence and domain conditions for gluing.
\item Contour functionals, local evaluation, and global rigidity.
\item The ordinary Hartogs and Baire lemmas, whose availability in the chosen
dependency revision must be independently checked.
\item The final analytic theorems in the proposed architecture remain
conditional until those ordinary complex-analytic dependencies have been checked.
\end{enumerate}

\subsection{Recovered on audit}
An audit of this merge against its three sources found the following nine
limitations stated in a source and missing from the merged text. They are
restored here rather than silently dropped, since each narrows a claim the
report otherwise appears to make.

\begin{enumerate}[resume]
\item Strong Hahn summability is not ordinary convergence. A formal series with
arbitrary ordinary coefficients can be evaluated at an infinitesimal even when
its ordinary analytic radius is zero, and doing so does \emph{not} make it an
ordinary holomorphic germ (M2).
\item The simplest-number theorem, linking Conway option pairs to the unique
sign sequence, is a theorem requiring proof. It is not a definition licensing
circular recursion (M1).
\item A proof that selects the first nonzero coefficient of an arbitrary
real-coefficient sequence may conceal excluded middle or a decidability
assumption. The supplied support and normal-form proofs use exactly that step
(M1).
\item No proof-assistant setting turns a full proper class into an ordinary
small type merely by renaming it. This is not special to Lean's \texttt{Small}
(M1).
\item How high an operation's birthdays can grow and how large the collection
of descriptions below a bound becomes are two different questions. A bound
solving one need not solve the other (M1).
\item That $\No\cong\R((\omega^{\No}))$ does not mean a fixed-point equation
for an unknown field automatically has a unique solution (M2).
\item An abstract proof using a basis of a class-sized vector space would need
separate justification and must not be substituted for a concrete argument
without notice (M2).
\item The gluing theorem is not based on an arbitrary infinite union of
unrelated well-ordered supports; the connected-overlap hypothesis is doing work
(M3).
\item No bibliographic entry here attributes authorship to an uploaded
manuscript merely from its filename (M3, followed in practice throughout
\cref{found:app:ledger}).
\end{enumerate}

\subsection{On this merged document itself}
\label{found:app:limmerged}
\begin{enumerate}\setcounter{enumi}{82}
\item This report merges three independently written audits.
\item Their agreement on the pinned Lean revision reflects \emph{three readings
of one source}, not independent execution.
\item Of the claims in this report, exactly one has been kernel-checked:
\cref{found:prop:allcuts}, the unrestricted-cut obstruction, in the two Lean
encodings of \cref{found:sub:compiled}, each depending on no axioms. Nothing
else here was machine-checked, and that one check says nothing about the
surrounding mathematics.
\item Where the three audits differ in scope, sources or emphasis, the
differences are printed rather than averaged.
\item No merged sentence asserts more than the strongest single source proves.
\item The merged report's greater length, its five source manuscripts, its
proved propositions and its bibliography are not execution. The union of three
read-only audits is a larger blueprint, not a certificate.
\end{enumerate}

\begingroup\small
\begin{thebibliography}{99}
\addcontentsline{toc}{section}{References}
\setlength{\itemsep}{0.4em}\raggedright

\bibitem{found:SrcA}
\emph{Surcomplex Analysis: Infinitesimal Calculus, Hahn-Coherent Holomorphy, and
Contour Theory}. Supplied manuscript, September 21, 2026. Filed as
\texttt{surcomplex\_analysis(3).tex} by M1 and M3 and as
\texttt{surcomplex\_analysis(2).tex} by M2; one file under two labels, see
\cref{found:app:ledger}. Especially the class-theoretic conventions,
fine-topological degeneracy, small-scale evaluation and rescaling,
support-preserving gluing with its componentwise convention, affine charts,
contour theory, and all-scale polynomial and rational rigidity. Read by all
three audits. Not independently proof-assistant verified.

\bibitem{found:SrcAG}
\emph{Infinitesimal Analytic Geometry over the Surcomplex Numbers: Radius-Free
Hahn Germs, a Nullstellensatz, and Conservation of Singular Zeros}. Supplied
manuscript, September 21, 2026. File
\texttt{surcomplex\_analytic\_geometry.tex}. Hahn-workspace, radius-free-germ,
finite-deformation and workspace-invariance sections. \textbf{Read by M1 only};
no hash recorded. Not independently proof-assistant verified.

\bibitem{found:SrcTrig}
\emph{Trigonometry on the Surcomplex Plane: Canonical Angles, Infinite
Triangles, Infinitesimal Contact, and Hahn-Analytic Oscillation}. Supplied
manuscript, September 21, 2026. File \texttt{surcomplex\_trigonometry.tex},
SHA-256 \texttt{e61577ff\ldots7639ad4600}. Size and support conventions, the
angle group, and global phases. \textbf{Read by M2 only}. Not independently
proof-assistant verified.

\bibitem{found:SrcP}
\emph{Surcomplex Polynomial Algebra: Order Geometry, Newton Profiles, and
Scale-Resolved Stability}. Supplied manuscript, September 21, 2026. File
\texttt{surcomplex\_polynomial\_algebra(2).tex}, 112{,}766 bytes, SHA-256
\texttt{9d2ad437\ldots3dfbfe516005}. Set-sized workspaces; support and
summability; fixed-degree real-closed-field transfer; support-local Hensel
recursion; set-sized rings and polynomial systems; dependency and effectiveness
boundary. \textbf{Read by M3 only}. Not an independently verified publication.

\bibitem{found:SrcR}
\emph{Hartogs Coherence, Finite Maps, and Residue Duality over the Surcomplex
Numbers}. Supplied manuscript, September 21, 2026. File
\texttt{surcomplex\_research.tex}, 98{,}818 bytes, SHA-256
\texttt{91303273\ldots2f121c857c712}. Workspaces and evaluation; positive
operators; automatic support and class replacement in Hartogs; fixed-domain
division and parameters; finite free algebras; finite divided differences and
actual fibers. \textbf{Read by M3 only}. Not independently proof-assistant
verified.

\bibitem{found:CompanionAnalysis}
\emph{Surcomplex Analysis: Germs, Coherent Hahn Sections, and Infinitesimal Zero
Geometry --- a theory of holomorphic functions on $\No[i]$}. Companion report in
this repository at \texttt{docs/surcomplex/analysis/}. Its all-scale polynomial
rigidity theorem is stated for \emph{class} functions on the whole of $\No[i]$
and quantifies over a proper class of scales; see \cref{found:sub:scope}. Cited
by title and repository path; none of its internal labels is referenced here.

\bibitem{found:CompanionTrig}
\emph{Trigonometry on the Surcomplex Plane: Canonical Phases, Arbitrary-Scale
Geometry, and Infinitesimal Degeneration}. Companion report in this repository at
\texttt{docs/surcomplex/trigonometry/}. Cited by title and repository path only.

\bibitem{found:Conway}
John H. Conway. \emph{On Numbers and Games}. Second edition, A K Peters, 2001;
first edition, Academic Press, 1976. Foundational constructions and arithmetic of
surreal numbers.

\bibitem{found:Gonshor}
Harry Gonshor. \emph{An Introduction to the Theory of Surreal Numbers}. London
Mathematical Society Lecture Note Series 110, Cambridge University Press, 1986.
\href{https://doi.org/10.1017/CBO9780511629143}{doi:10.1017/CBO9780511629143}.
Sign expansions, arithmetic, and normal forms.

\bibitem{found:Alling}
Norman L. Alling. \emph{Foundations of Analysis over Surreal Number Fields}.
North-Holland Mathematics Studies 141, North-Holland, 1987.

\bibitem{found:Poonen}
Bjorn Poonen. ``Maximally complete fields.'' \emph{L'Enseignement
Math\'ematique} (2) \textbf{39} (1993), 87--106.
\href{https://math.mit.edu/~poonen/papers/amsval.pdf}{Author-hosted text}.
Corollary 4 supplies the divisible-group, algebraically closed-coefficient
Hahn-field input used in \cref{found:eq:hahncomplex} and
\cref{found:thm:workspace}.

\bibitem{found:Neumann}
B. H. Neumann. ``On ordered division rings.'' \emph{Transactions of the American
Mathematical Society} \textbf{66} (1949), 202--252. The ordered-support and
finite-word summability machinery behind $E^*$.

\bibitem{found:Ehrlich2001}
Philip Ehrlich. ``Number systems with simplicity hierarchies: a generalization of
Conway's theory of surreal numbers.'' \emph{Journal of Symbolic Logic}
\textbf{66} (2001), 1231--1258.
\href{https://doi.org/10.2307/2695104}{doi:10.2307/2695104}. \emph{Must be read
together with the correction below.}

\bibitem{found:EhrlichCorrection}
Philip Ehrlich. Correction to ``Number systems with simplicity hierarchies: a
generalization of Conway's theory of surreal numbers.'' \emph{Journal of Symbolic
Logic} \textbf{70}(3) (2005), 1022.
\href{https://doi.org/10.2178/jsl/1122038926}{doi:10.2178/jsl/1122038926}.
Listed as a separate item because both M1 and M2 insist that detailed claims of
the 2001 paper be cited with it.

\bibitem{found:EhrlichKaplanII}
Philip Ehrlich and Elliot Kaplan. ``Number systems with simplicity hierarchies
II.'' \href{https://arxiv.org/abs/1512.04001}{arXiv:1512.04001}; published in
\emph{Journal of Symbolic Logic} (2018). Used for the structural
simplicity-hierarchy viewpoint, not for an asserted Lean implementation.

\bibitem{found:EhrlichKaplanExp}
Philip Ehrlich and Elliot Kaplan. ``Surreal ordered exponential fields.''
\emph{Journal of Symbolic Logic}, published online 2021.
\href{https://arxiv.org/abs/2002.07739}{arXiv:2002.07739}. Especially \S11 and
Theorem 11.1 on the canonical surcomplex exponential with omnific-integer
kernel, via an \emph{initial integer part}; see \cref{found:sub:canonicalexp}.

\bibitem{found:vdDE}
Lou van den Dries and Philip Ehrlich. ``Fields of surreal numbers and
exponentiation.'' \emph{Fundamenta Mathematicae} \textbf{167} (2001), 173--188.
\href{https://doi.org/10.4064/fm167-2-3}{doi:10.4064/fm167-2-3}. \emph{Read
together with the erratum below.}

\bibitem{found:vdDEerratum}
Lou van den Dries and Philip Ehrlich. Erratum to ``Fields of surreal numbers and
exponentiation.'' \emph{Fundamenta Mathematicae} \textbf{168} (2001), 295--297.
\href{https://doi.org/10.4064/fm168-3-5}{doi:10.4064/fm168-3-5}. Detailed
ordinal support estimates from the original article require this correction; a
separate item for that reason.

\bibitem{found:Tarski}
Lou van den Dries. ``Alfred Tarski's elimination theory for real closed fields.''
\emph{Journal of Symbolic Logic} \textbf{53} (1988), 7--19.
\href{https://doi.org/10.2307/2274424}{doi:10.2307/2274424}. First-order transfer
is used only in its specified language; see \cref{found:sub:qe}.

\bibitem{found:ShulmanSize}
Michael A. Shulman. ``Set theory for category theory.''
\href{https://arxiv.org/abs/0810.1279}{arXiv:0810.1279}, 2008. Used for
foundations and size-framework comparisons --- universes, reflection, small maps
--- not as a surreal construction theorem.

\bibitem{found:Williams}
Kameryn J. Williams. \emph{The Structure of Models of Second-order Set
Theories}. Ph.D.\ dissertation, City University of New York, 2018.
\href{https://arxiv.org/abs/1804.09526}{arXiv:1804.09526}. Source for
distinctions among GBC, KM, class realizations, and recursion principles.

\bibitem{found:WilliamsMin}
Kameryn J. Williams. ``Minimum models of second-order set theories.''
\emph{Journal of Symbolic Logic} \textbf{84} (2019), 589--620.
\href{https://doi.org/10.1017/jsl.2019.27}{doi:10.1017/jsl.2019.27};
\href{https://arxiv.org/abs/1709.03955}{arXiv:1709.03955}. See in particular the
definitions of GB, GBC, KM and ETR.

\bibitem{found:HamkinsETR}
Joel David Hamkins. ``Second-order transfinite recursion is equivalent to
Kelley--Morse set theory over GBC.'' Author's research exposition, July 23, 2017.
\href{https://jdh.hamkins.org/second-order-transfinite-recursion-is-equivalent-to-kelley-morse-set-theory/}{Online exposition}.

\bibitem{found:HamkinsRecursion}
Joel David Hamkins. ``Determinacy for proper-class clopen games is equivalent to
transfinite recursion along proper-class well-founded relations.'' Author's
mathematical exposition, 2014.
\href{https://jdh.hamkins.org/determinacy-for-proper-class-clopen-games-is-equivalent-to-transfinite-recursion-along-proper-class-well-founded-relations/}{Online exposition}.
Used for the distinction between set-like recursion and unrestricted class
recursion.

\bibitem{found:GitmanHamkins}
Victoria Gitman and Joel David Hamkins. ``Open determinacy for class games.''
Author preprint, 2015. \href{https://arxiv.org/abs/1509.01099}{arXiv:1509.01099}.
Class-game determinacy and elementary transfinite recursion.

\bibitem{found:ClassForcing}
Victoria Gitman, Joel David Hamkins, Peter Holy, Philipp Schlicht and Kameryn J.
Williams. ``The exact strength of the class forcing theorem.'' \emph{Journal of
Symbolic Logic}, published online July 29, 2020;
\href{https://arxiv.org/abs/1707.03700}{arXiv:1707.03700}. Establishes the role
of $\ETR_{\Ord}$; the strength of class-recursive principles over GBC.

\bibitem{found:ZFCminus}
Victoria Gitman, Joel David Hamkins and Thomas A. Johnstone. ``What is the
theory ZFC without power set?''
\href{https://arxiv.org/abs/1110.2430}{arXiv:1110.2430}, 2011. Used only for the
warning about weakening set-existence and Collection principles.

\bibitem{found:Aczel}
Peter Aczel. ``The type theoretic interpretation of constructive set theory:
inductive definitions.'' In \emph{Logic, Methodology and Philosophy of Science
VII}, Studies in Logic and the Foundations of Mathematics 114, Elsevier, 1986,
17--49.
\href{https://doi.org/10.1016/S0049-237X(09)70683-4}{doi:10.1016/S0049-237X(09)70683-4}.

\bibitem{found:Rathjen}
Michael Rathjen. ``On the regular extension axiom and its variants.''
\emph{Mathematical Logic Quarterly} \textbf{49} (2003), 511--518.
\href{https://doi.org/10.1002/malq.200310054}{doi:10.1002/malq.200310054}.

\bibitem{found:GambinoAczel}
Nicola Gambino and Peter Aczel. ``The generalised type-theoretic interpretation
of constructive set theory.'' \emph{Journal of Symbolic Logic} \textbf{71}
(2006), 67--103.
\href{https://doi.org/10.2178/jsl/1140641163}{doi:10.2178/jsl/1140641163}.

\bibitem{found:HoTTBook}
The Univalent Foundations Program. \emph{Homotopy Type Theory: Univalent
Foundations of Mathematics}. Institute for Advanced Study, 2013, \S11.6 ``The
surreal numbers.'' \href{https://homotopytypetheory.org/book/}{Book page};
\href{https://github.com/HoTT/book/blob/master/reals.tex}{public chapter source}.
Higher inductive-inductive construction with a universe-relative small-option
condition.

\bibitem{found:ShulmanPlump}
Michael Shulman. ``The surreals contain the plump ordinals.'' HoTT research blog,
February 22, 2014.
\href{https://homotopytypetheory.org/2014/02/22/surreals-plump-ordinals/}{Author's exposition}.
Constructive order evidence and ordinal comparison.

\bibitem{found:LumsdaineShulman}
Peter LeFanu Lumsdaine and Michael Shulman. ``The semantics of higher inductive
types.'' Author preprint, 2017.
\href{https://arxiv.org/abs/1705.07088}{arXiv:1705.07088}. General background,
not a claim that every surreal higher inductive-inductive signature is
automatically supported.

\bibitem{found:CubicalHIT}
Thierry Coquand, Simon Huber and Anders M\"ortberg. ``On higher inductive types
in cubical type theory.'' Author preprint, 2018.
\href{https://arxiv.org/abs/1802.01170}{arXiv:1802.01170}. Cubical
higher-inductive background.

\bibitem{found:Joseph}
Jean S. Joseph. ``The surreal numbers as an ordered commutative ring with an
apartness: a development in univalent foundations.''
\href{https://arxiv.org/abs/1812.00051}{arXiv:1812.00051}, 2018.

\bibitem{found:RangelMariano}
Dimi Rocha Rangel and Hugo Luiz Mariano. ``An algebraic (set) theory of surreal
numbers, I.'' \href{https://arxiv.org/abs/1911.12726}{arXiv:1911.12726}, 2019.

\bibitem{found:PakKaliszyk}
Karol P\k{a}k and Cezary Kaliszyk. ``Conway normal form: bridging approaches for
comprehensive formalization of surreal numbers.'' In \emph{15th International
Conference on Interactive Theorem Proving (ITP 2024)}, LIPIcs \textbf{309},
Article 29, 29:1--29:18, 2024.
\href{https://doi.org/10.4230/LIPIcs.ITP.2024.29}{doi:10.4230/LIPIcs.ITP.2024.29};
\href{https://arxiv.org/abs/2410.00065}{arXiv:2410.00065}. Mizar formalization
on a Tarski--Grothendieck foundation. The abstract's ``algebraically closed''
wording is \emph{not} used as a field theorem here: the carrier $\No$ is real
closed, and algebraic closedness belongs to $\No[i]$.

\bibitem{found:Hamkins2026}
Joel David Hamkins, reporting joint work with Junhong Chen and Ruizhi Yang.
``Surreal arithmetic is bi-interpretable with set theory.'' CUNY Logic Workshop
lecture announcement, posted February 4, 2026; lecture March 13, 2026.
\href{https://jdh.hamkins.org/surreal-arithmetic-cuny-logic-workshop-march-2026/}{Author's announcement}.
\textbf{Work in progress}; cited with that status, and not used as a proved
dependency, a checked axiom list, or a consistency proof.

\bibitem{found:Hamkins2026Slides}
Joel David Hamkins. Slides, ``Elementary theory of surreal arithmetic,'' CUNY,
March 2026.
\href{https://jdh.hamkins.org/wp-content/uploads/Slides-Elementary-theory-of-surreal-arithmetic-Hamkins-CUNY-March-2026.pdf}{Slides}.
\textbf{Inspected by M3 only}; the slides explicitly mark the work as new with
details still being checked. Listed separately from the announcement because
only the slides were actually read, and only by one audit.

\bibitem{found:LeanUniverses}
Lean development team. \emph{The Lean Language Reference}, ``Universes'' (The
Type System).
\href{https://lean-lang.org/doc/reference/latest/The-Type-System/Universes/}{Official documentation}.
\emph{Rolling documentation, inspected September 21, 2026.}

\bibitem{found:LeanPropositions}
Lean development team. \emph{The Lean Language Reference}, ``Propositions'' (The
Type System).
\href{https://lean-lang.org/doc/reference/latest/The-Type-System/Propositions/}{Official documentation}.
\emph{Rolling documentation, inspected September 21, 2026.}

\bibitem{found:LeanAxioms}
Lean development team. \emph{The Lean Language Reference}, ``Axioms.''
\href{https://lean-lang.org/doc/reference/latest/Axioms/}{Official documentation}.
\emph{Rolling documentation, inspected September 21, 2026.} Logical axioms,
unfinished proofs, and computational trust.

\bibitem{found:LeanValidation}
Lean development team. \emph{The Lean Language Reference}, ``Validating Proofs.''
\href{https://lean-lang.org/doc/reference/latest/ValidatingProofs/}{Official documentation}.
\emph{Rolling documentation, inspected September 21, 2026.} Axiom inspection and
validation boundaries.

\bibitem{found:MathlibHahnBasic}
Mathlib contributors. \texttt{Mathlib.RingTheory.HahnSeries.Basic},
\texttt{Multiplication} and \texttt{Lex}.
\href{https://leanprover-community.github.io/mathlib4_docs/Mathlib/RingTheory/HahnSeries/Basic.html}{Basic documentation};
\href{https://leanprover-community.github.io/mathlib4_docs/Mathlib/RingTheory/HahnSeries/Lex.html}{Lex documentation}.
\emph{Rolling documentation, inspected September 21, 2026.}

\bibitem{found:MathlibHahnSummable}
Mathlib contributors. \path{Mathlib.RingTheory.HahnSeries.Summable}.
\href{https://leanprover-community.github.io/mathlib4_docs/Mathlib/RingTheory/HahnSeries/Summable.html}{Generated documentation}.
\emph{Rolling documentation, inspected September 21, 2026.} Formal summable
families, coefficient finiteness and \texttt{hsum}; not a generic topological
convergence API.

\bibitem{found:MathlibHEval}
Mathlib contributors. \texttt{Mathlib.RingTheory.HahnSeries.HEval}.
\href{https://leanprover-community.github.io/mathlib4_docs/Mathlib/RingTheory/HahnSeries/HEval.html}{Generated documentation}.
\emph{Rolling documentation, inspected September 21, 2026, by M1 only.} The
finite-variable \path{MvPowerSeries.heval} TODO was read from this page.

\bibitem{found:MathlibRealClosed}
Mathlib contributors. \path{Mathlib.FieldTheory.IsRealClosed.Basic}.
\href{https://leanprover-community.github.io/mathlib4_docs/Mathlib/FieldTheory/IsRealClosed/Basic.html}{Generated documentation}.
\emph{Rolling documentation, inspected September 21, 2026, by M2 only.} Presence
of the interface is not a verification of a surreal instance.

\bibitem{found:MathlibPSet}
Mathlib contributors. \texttt{Mathlib.SetTheory.ZFC.PSet}.
\href{https://leanprover-community.github.io/mathlib4_docs/Mathlib/SetTheory/ZFC/PSet.html}{Generated documentation}.
\emph{Rolling documentation, inspected September 21, 2026.}

\bibitem{found:MathlibZFBasic}
Mathlib contributors. \path{Mathlib.SetTheory.ZFC.Basic}.
\href{https://leanprover-community.github.io/mathlib4_docs/Mathlib/SetTheory/ZFC/Basic.html}{Generated documentation}.
\emph{Rolling documentation, inspected September 21, 2026.} The modeled ZFC set
type, Lean's predicate-based sets, and the technical \texttt{Definable}
interface are distinct.

\bibitem{found:MathlibClass}
Mathlib contributors. \texttt{Mathlib.SetTheory.ZFC.Class}.
\href{https://leanprover-community.github.io/mathlib4_docs/Mathlib/SetTheory/ZFC/Class.html}{Generated documentation}.
\emph{Rolling documentation, inspected September 21, 2026.}

\bibitem{found:GamesRepo}
Violeta Hern\'andez Palacios and contributors. \emph{Combinatorial Games}, Lean 4
repository \texttt{vihdzp/combinatorial-games} --- a separate project, not
upstream Mathlib.
\href{https://github.com/vihdzp/combinatorial-games/tree/02b4a908ea2ecfefffecb438f691951a814a5264}{Pinned source tree}
at revision \texttt{02b4a908ea2ecfefffecb438f691951a814a5264}, commit timestamp
September 19, 2026, 10:25:40 UTC; inspected September 21, 2026 by source reading
only.

\bibitem{found:GamesToolchain}
Same repository and revision. \path{lean-toolchain}.
\href{https://github.com/vihdzp/combinatorial-games/blob/02b4a908ea2ecfefffecb438f691951a814a5264/lean-toolchain}{Pinned file}.
Records \texttt{leanprover/lean4:v4.35.0-rc2}.

\bibitem{found:GamesManifest}
Same repository and revision. \path{lake-manifest.json}.
\href{https://github.com/vihdzp/combinatorial-games/blob/02b4a908ea2ecfefffecb438f691951a814a5264/lake-manifest.json}{Pinned file}.
Records the Mathlib revision
\texttt{dec5b2b780537b6eaf7f5e5f000c12f7387fb24d}.

\bibitem{found:GamesUmbrella}
Same repository and revision. \path{CombinatorialGames.lean}.
\href{https://github.com/vihdzp/combinatorial-games/blob/02b4a908ea2ecfefffecb438f691951a814a5264/CombinatorialGames.lean}{Pinned file}.
Umbrella module inventory; an import list is not an independent proof of each
module's mathematical scope.

\bibitem{found:GamesBasic}
Same repository and revision. \path{CombinatorialGames/Surreal/Basic.lean}.
\href{https://github.com/vihdzp/combinatorial-games/blob/02b4a908ea2ecfefffecb438f691951a814a5264/CombinatorialGames/Surreal/Basic.lean}{Pinned file}.
Carrier at $\Type_{u+1}$ by antisymmetrization of the numeric-game subtype;
quotient map, order, additive structure, choice-based representative, and a
set-based cut interface with smallness conditions.

\bibitem{found:GamesDivision}
Same repository and revision. \path{CombinatorialGames/Surreal/Division.lean}.
\href{https://github.com/vihdzp/combinatorial-games/blob/02b4a908ea2ecfefffecb438f691951a814a5264/CombinatorialGames/Surreal/Division.lean}{Pinned file}.
The numeric inverse congruence and an actual
\texttt{noncomputable instance : Field Surreal} were inspected in source.

\bibitem{found:GamesHahn}
Same repository and revision.
\path{CombinatorialGames/Surreal/HahnSeries/Basic.lean}.
\href{https://github.com/vihdzp/combinatorial-games/blob/02b4a908ea2ecfefffecb438f691951a814a5264/CombinatorialGames/Surreal/HahnSeries/Basic.lean}{Pinned file}.
Small-support Hahn carrier over order-dual exponents; \texttt{mk} requires a
\texttt{Small.\{u\}} witness and reverse well-foundedness; coefficients,
truncation and ordinal support indexing. \emph{No completed normal-form
isomorphism was verified in this inspected file.}

\bibitem{found:GamesCut}
Same repository and revision. \path{CombinatorialGames/Surreal/Cut.lean}.
\href{https://github.com/vihdzp/combinatorial-games/blob/02b4a908ea2ecfefffecb438f691951a814a5264/CombinatorialGames/Surreal/Cut.lean}{Pinned file}.
Complementary lower and upper sections of the entire fixed-universe order, built
via concept lattices --- distinct from small Conway option pairs.
\textbf{Inspected by M3 only.}

\end{thebibliography}
\endgroup
\end{document}
